% ============================================================
%  THE ECOLOGY OF DISTINCTIONS
%  Author: Flyxion
%  Engine: LuaLaTeX
%  Build:  lualatex main && biber main &&
%          lualatex main && makeindex main.idx &&
%          lualatex main && lualatex main
% ============================================================
\documentclass[12pt,twoside,openright]{book}

% ---- Fonts -------------------------------------------------
\usepackage{fontspec}
\setmainfont{TeX Gyre Pagella}
\setsansfont{TeX Gyre Heros}
\setmonofont[Scale=0.85]{TeX Gyre Cursor}
\usepackage{unicode-math}
\setmathfont{TeX Gyre Pagella Math}

% ---- Geometry ----------------------------------------------
\usepackage[paperwidth=17cm,paperheight=24cm,
  top=2.5cm,bottom=2.8cm,inner=2.8cm,outer=2.2cm,
  marginparwidth=1.4cm,marginparsep=0.3cm]{geometry}

% ---- Colour ------------------------------------------------
\usepackage[dvipsnames,svgnames,x11names]{xcolor}
\definecolor{DistBlue}  {HTML}{1B4F72}
\definecolor{DistTeal}  {HTML}{117A65}
\definecolor{DistAmber} {HTML}{B7770D}
\definecolor{DistRed}   {HTML}{922B21}
\definecolor{DistGray}  {HTML}{566573}
\definecolor{DistGold}  {HTML}{7D6608}
\definecolor{DistLight} {HTML}{EAF4FB}
\definecolor{DistLightG}{HTML}{E9F7EF}
\definecolor{DistGoldBg}{HTML}{FEF9E7}

% ---- TikZ --------------------------------------------------
\usepackage{tikz,pgfplots}
\pgfplotsset{compat=1.18}
\usetikzlibrary{arrows.meta,positioning,shapes.geometric,
  decorations.pathmorphing,backgrounds,fit,calc,matrix,patterns}
\tikzset{
  thmnode/.style={rectangle,rounded corners=4pt,
    draw=DistTeal,fill=DistLight,font=\sffamily\footnotesize,
    text width=4.2cm,align=center,inner sep=4pt},
  lawnode/.style={rectangle,rounded corners=4pt,
    draw=DistGold,fill=DistGoldBg,
    font=\sffamily\footnotesize\bfseries,
    text width=4.2cm,align=center,inner sep=4pt},
  thmedge/.style={-Stealth,DistTeal,thick},
  distnode/.style={circle,draw=DistBlue,fill=DistLight,
    minimum size=7mm,font=\small\sffamily},
  distedge/.style={-Stealth,DistTeal,thick},
}

% ---- Chapter/section style ---------------------------------
\usepackage{titlesec}
\titleformat{\chapter}[display]
  {\sffamily\bfseries\color{DistBlue}\huge}
  {\color{DistTeal}\sffamily\Large\chaptername\ \thechapter}
  {0.5ex}{\titlerule[1pt]\vspace{0.8ex}}[\vspace{0.4ex}\titlerule]
\titleformat{\section}
  {\sffamily\bfseries\large\color{DistBlue}}{\thesection}{1em}{}
\titleformat{\subsection}
  {\sffamily\bfseries\normalsize\color{DistGray}}{\thesubsection}{1em}{}

% ---- Headers/footers ---------------------------------------
\usepackage{fancyhdr}
\pagestyle{fancy}\fancyhf{}
\fancyhead[LE]{\small\sffamily\color{DistGray}\leftmark}
\fancyhead[RO]{\small\sffamily\color{DistGray}\rightmark}
\fancyfoot[C]{\small\sffamily\thepage}
\renewcommand{\headrulewidth}{0.4pt}

% ---- Maths -------------------------------------------------
\usepackage{amsmath,amssymb,amsthm,mathtools}
\usepackage[most]{tcolorbox}

% ---- Theorem environments ----------------------------------
\tcbset{thmbase/.style={enhanced,breakable,
  left=4pt,right=4pt,top=6pt,bottom=4pt,
  fonttitle=\sffamily\bfseries,separator sign={.}}}

% Level 0 — Gold
\newtcbtheorem[number within=chapter]{axiom}{Axiom}{
  thmbase,colback=DistGoldBg,colframe=DistGold,coltitle=white,
  borderline west={3pt}{0pt}{DistGold},
  attach boxed title to top left={yshift=-2mm,xshift=4mm},
  boxed title style={colback=DistGold,rounded corners}}{axm}
\newtcbtheorem[use counter from=axiom]{law}{Law}{
  thmbase,colback=DistGoldBg,colframe=DistGold!80,coltitle=white,
  borderline west={3pt}{0pt}{DistGold!80},
  attach boxed title to top left={yshift=-2mm,xshift=4mm},
  boxed title style={colback=DistGold!80,rounded corners}}{law}
\newtcbtheorem[use counter from=axiom]{principle}{Principle}{
  thmbase,colback=DistGoldBg,colframe=DistGold!60,coltitle=white,
  borderline west={3pt}{0pt}{DistGold!60},
  attach boxed title to top left={yshift=-2mm,xshift=4mm},
  boxed title style={colback=DistGold!60,rounded corners}}{prc}

% Level 1 — Teal
\newtcbtheorem[number within=chapter]{theorem}{Theorem}{
  thmbase,colback=DistLight,colframe=DistTeal,coltitle=white,
  attach boxed title to top left={yshift=-2mm,xshift=4mm},
  boxed title style={colback=DistTeal,rounded corners}}{thm}

% Level 2 — Blue
\newtcbtheorem[use counter from=theorem]{proposition}{Proposition}{
  thmbase,colback=DistLight,colframe=DistBlue,coltitle=white,
  attach boxed title to top left={yshift=-2mm,xshift=4mm},
  boxed title style={colback=DistBlue,rounded corners}}{prop}

% Level 3 — Gray
\newtcbtheorem[use counter from=theorem]{lemma}{Lemma}{
  thmbase,colback=DistLight,colframe=DistGray,coltitle=white,
  attach boxed title to top left={yshift=-2mm,xshift=4mm},
  boxed title style={colback=DistGray,rounded corners}}{lem}
\newtcbtheorem[use counter from=theorem]{corollary}{Corollary}{
  thmbase,colback=DistLight,colframe=DistTeal!50,coltitle=DistBlue,
  attach boxed title to top left={yshift=-2mm,xshift=4mm},
  boxed title style={colback=DistTeal!50,rounded corners}}{cor}

% Definition / Example / Remark
\newtcbtheorem[number within=chapter]{definition}{Definition}{
  thmbase,colback=DistLightG,colframe=DistTeal,coltitle=white,
  attach boxed title to top left={yshift=-2mm,xshift=4mm},
  boxed title style={colback=DistTeal,rounded corners}}{defn}
\newtcbtheorem[number within=chapter]{example}{Example}{
  thmbase,colback=orange!5,colframe=DistAmber!80,coltitle=white,
  attach boxed title to top left={yshift=-2mm,xshift=4mm},
  boxed title style={colback=DistAmber!80,rounded corners}}{ex}
\newtcbtheorem[number within=chapter]{remark}{Remark}{
  thmbase,colback=yellow!6,colframe=DistAmber,coltitle=white,
  attach boxed title to top left={yshift=-2mm,xshift=4mm},
  boxed title style={colback=DistAmber,rounded corners}}{rmk}

\renewcommand{\qedsymbol}{\textcolor{DistTeal}{$\blacksquare$}}

% ---- Listings ----------------------------------------------
\usepackage{listings}
\lstset{basicstyle=\ttfamily\small,
  keywordstyle=\color{DistBlue}\bfseries,
  commentstyle=\color{DistGray}\itshape,
  stringstyle=\color{DistAmber},
  frame=single,rulecolor=\color{DistGray!50},
  numbers=left,numberstyle=\tiny\color{DistGray},
  breaklines=true,captionpos=b,tabsize=2,
  showstringspaces=false}
\lstdefinelanguage{Spherepop}{
  morekeywords={pop,refuse,collapse,bind,emit,reach,repair},
  sensitive=false,morecomment=[l]{//},morestring=[b]"}

% ---- Tables / figures --------------------------------------
\usepackage{graphicx,booktabs,caption,subcaption}
\captionsetup{font=small,labelfont={sf,bf,color=DistBlue}}

% ---- Margin icons ------------------------------------------
\usepackage{fontawesome5,enumitem}
\newcommand{\csicon} {\marginpar{\tiny\color{DistBlue}\faCode}}
\newcommand{\bioicon}{\marginpar{\tiny\color{DistTeal}\faDna}}
\newcommand{\physicon}{\marginpar{\tiny\color{DistGray}\faAtom}}

% ---- Index -------------------------------------------------
\usepackage{makeidx}
\makeindex

% ---- Glossaries --------------------------------------------
\usepackage[acronym,toc,nonumberlist]{glossaries}
\makeglossaries
\newacronym{rsvp}{RSVP}{Relativistic Scalar Vector Plenum}
\newacronym{tartan}{TARTAN}{Trajectory-Aware Recursive Tiling with Annotated Noise}
\newacronym{cmb}{CMB}{Cosmic Microwave Background}
\newacronym{hydra}{HYDRA}{Hybrid Dynamic Reasoning Architecture}
\newglossaryentry{distinction}{name={distinction},
  description={A partition of domain $X$ into non-identical
  regions; the primitive act from which objects, information,
  and entropy arise (Ch.\,1)}}
\newglossaryentry{recoverability}{name={recoverability},
  description={The degree to which dispersed distinctions remain
  reconstructible (Ch.\,5)}}
\newglossaryentry{repair}{name={repair},
  description={Operator $\repair:\mathcal{D}\to\mathcal{D}$
  restoring a damaged distinction; possible iff $\reco>0$ (Ch.\,7)}}
\newglossaryentry{admissibility}{name={admissibility},
  description={Property of trajectories preserving admissible
  future distinction-producing volume (Ch.\,14)}}
\newglossaryentry{generative-admissibility}{
  name={generative admissibility},
  description={Condition $\frac{d}{dt}\Vol(\adm(t))\ge0$;
  the deepest evaluative criterion (Ch.\,15, Ch.\,31)}}

% ---- Hyperlinks (late) -------------------------------------
\usepackage[colorlinks=true,linkcolor=DistBlue,
  citecolor=DistTeal,urlcolor=DistAmber,
  pdftitle={The Ecology of Distinctions},
  pdfauthor={Flyxion}]{hyperref}
\usepackage{cleveref}

% ---- Bibliography ------------------------------------------
\usepackage[backend=biber,style=authoryear,
  sorting=nyt,maxbibnames=6]{biblatex}
\addbibresource{distinctions.bib}

% ============================================================
%  CUSTOM COMMANDS
% ============================================================
\newcommand{\vfield}{\vec{v}}
\newcommand{\partition}{\mathcal{P}}
\newcommand{\hist}{\mathcal{H}}
\newcommand{\reach}{\mathcal{R}}
\newcommand{\adm}{\mathcal{A}}
\newcommand{\repair}{\mathfrak{R}}
\newcommand{\reco}{\mathfrak{rec}}
\newcommand{\Vol}{\operatorname{Vol}}

\newenvironment{epigraph}[2]{%
  \vspace{1em}\begin{flushright}
  \begin{minipage}{0.72\textwidth}
  \itshape\small #1\\[0.4em]
  \upshape\footnotesize\sffamily\color{DistGray}--- #2
  \end{minipage}\end{flushright}\vspace{1.5em}}{}

\newenvironment{objectives}{%
  \begin{tcolorbox}[colback=DistBlue!5,colframe=DistBlue,
    title={\faCheckCircle\enspace Learning Objectives},
    fonttitle=\sffamily\bfseries\color{white},
    attach boxed title to top left={yshift=-2mm,xshift=4mm},
    boxed title style={colback=DistBlue},breakable]
  \begin{itemize}[leftmargin=*,nosep]}{%
  \end{itemize}\end{tcolorbox}\vspace{1em}}

\newenvironment{chapsummary}{%
  \begin{tcolorbox}[colback=DistBlue!5,colframe=DistBlue,
    title={\faBookOpen\enspace Summary},
    fonttitle=\sffamily\bfseries\color{white},
    attach boxed title to top left={yshift=-2mm,xshift=4mm},
    boxed title style={colback=DistBlue,rounded corners}]
  \begin{itemize}[nosep,leftmargin=*]}{%
  \end{itemize}\end{tcolorbox}}

\newcounter{exercise}[chapter]
\newenvironment{exercise}[1][]{%
  \refstepcounter{exercise}%
  \noindent\textbf{\sffamily\color{DistBlue}%
  Exercise~\thechapter.\theexercise%
  \ifx&#1&\else~(#1)\fi.}\enspace}{\par\medskip}

% ============================================================
\begin{document}
\frontmatter

% ---- Title page --------------------------------------------
\begin{titlepage}
\centering\vspace*{2.5cm}
{\sffamily\Huge\bfseries\color{DistBlue}
  The Ecology of Distinctions}\\[0.9em]
{\sffamily\Large\color{DistGray}
  From Information and Entropy\\
  to Repair, Regeneration, and Admissibility}\\[2.2cm]
{\large Flyxion}\\[0.4em]
{\sffamily\color{DistGray}Independent Researcher}\\[2.5cm]
\begin{tikzpicture}
  \node[distnode,minimum size=1.8cm](X)at(0,0){$X$};
  \node[distnode,minimum size=1.2cm,fill=DistTeal!20]
    (A)at(-2.5,-2){$P_1$};
  \node[distnode,minimum size=1.2cm,fill=DistAmber!20]
    (B)at(0,-2){$P_2$};
  \node[distnode,minimum size=1.2cm,fill=DistRed!10]
    (C)at(2.5,-2){$P_3$};
  \draw[distedge](X)--(A);\draw[distedge](X)--(B);
  \draw[distedge](X)--(C);
  \node[below=.3cm of A,font=\tiny\sffamily\color{DistGray}]
    {marked};
  \node[below=.3cm of B,font=\tiny\sffamily\color{DistGray}]
    {unmarked};
  \node[below=.3cm of C,font=\tiny\sffamily\color{DistGray}]
    {boundary};
\end{tikzpicture}\\[2cm]
{\sffamily\color{DistGray}First Edition~·~2025}
\end{titlepage}

\thispagestyle{empty}\vspace*{\fill}
{\small\noindent\textcopyright{} 2025 Flyxion.
Released under CC~BY-NC-SA~4.0. Typeset with Lua\LaTeX.}
\clearpage

\thispagestyle{empty}\vspace*{5cm}
\begin{center}\itshape
For everyone who noticed that the categories were wrong\\
and tried to fix them anyway.
\end{center}\clearpage

% ---- Conceptual Backbone -----------------------------------
\chapter*{The Conceptual Backbone}
\addcontentsline{toc}{chapter}{The Conceptual Backbone}
\thispagestyle{empty}
\begin{tcolorbox}[colback=DistGoldBg,colframe=DistGold,
  fonttitle=\sffamily\bfseries\color{white},
  title={\faLayerGroup\enspace Eight Structural Results},breakable]
\begin{enumerate}[leftmargin=*,itemsep=0.6em,label=\textbf{\arabic*.}]
  \item \textbf{Axiom of Distinction} (Ch.\,1)\\
  {\small Every observation presupposes a distinction.
  Objects, information, cost, and blind spots arise
  simultaneously from the act of partitioning.}
  \item \textbf{Information--Distinction Theorem} (Ch.\,2)\\
  {\small Information is the quantitative expression of
  distinctions induced by a partition of possibility space.}
  \item \textbf{Distinction--Entropy Duality} (Ch.\,3)\\
  {\small Entropy measures latent multiplicity hidden beneath
  a distinction structure. The Second Law is a theorem about
  distinction erosion.}
  \item \textbf{Law of Historical Compression} (Ch.\,4)\\
  {\small States are projections of histories. History
  ontology strictly contains state ontology.}
  \item \textbf{Law of Recoverability} (Ch.\,5)\\
  {\small Dispersal and destruction are not equivalent.
  Recoverability is the exact precondition for repair.}
  \item \textbf{Principle of Repair} (Ch.\,7)\\
  {\small Persistent distinctions require repair. Repair
  is possible iff recoverability is strictly positive.}
  \item \textbf{Principle of Regeneration} (Ch.\,11)\\
  {\small Regenerative systems preserve their repair capacity.
  Mere continuation is not regeneration.}
  \item \textbf{Generative Admissibility Principle} (Ch.\,31)\\
  {\small A trajectory is valuable insofar as it preserves
  the capacity for future distinction-production:
  $\frac{d}{dt}\Vol(\adm(t))\ge0$.}
\end{enumerate}
\end{tcolorbox}\clearpage

% ---- Theorem Dependency Graph ------------------------------
\chapter*{Theorem Dependency Graph}
\addcontentsline{toc}{chapter}{Theorem Dependency Graph}
\thispagestyle{empty}
\begin{center}
\begin{tikzpicture}[node distance=1.05cm and 0.2cm,
  every node/.style={font=\sffamily\scriptsize,align=center,
    text width=3.8cm,inner sep=3pt}]
\node[lawnode](A1){Axiom of Distinction\\{\tiny Ch.\,1}};
\node[lawnode,below=of A1](A2){Information--Distinction\\{\tiny Ch.\,2}};
\node[lawnode,below=of A2](A3){Distinction--Entropy Duality\\{\tiny Ch.\,3}};
\node[lawnode,below=of A3](A4){Historical Compression\\{\tiny Ch.\,4}};
\node[lawnode,below=of A4](A5){Law of Recoverability\\{\tiny Ch.\,5}};
\node[lawnode,below=of A5](A6){Principle of Repair\\{\tiny Ch.\,7}};
\node[lawnode,below=of A6](A7){Principle of Regeneration\\{\tiny Ch.\,11}};
\node[lawnode,below=of A7](A8){Generative Admissibility\\{\tiny Ch.\,31}};
\foreach \u/\v in {A1/A2,A2/A3,A3/A4,A4/A5,A5/A6,A6/A7,A7/A8}
  \draw[thmedge](\u)--(\v);
\node[thmnode,right=2.4cm of A1](B1){Projection Loss\\{\tiny Ch.\,2}};
\node[thmnode,right=2.4cm of A2](B2){Entropy Production\\{\tiny Ch.\,3}};
\node[thmnode,right=2.4cm of A3](B3){History Primacy\\{\tiny Ch.\,4}};
\node[thmnode,right=2.4cm of A4](B4){Semantic Horizon\\{\tiny Ch.\,5}};
\node[thmnode,right=2.4cm of A5](B5){Repair Closure\\{\tiny Ch.\,7}};
\node[thmnode,right=2.4cm of A6](B6){Regeneration Thm\\{\tiny Ch.\,11}};
\node[thmnode,right=2.4cm of A7](B7){Reachability Mono.\\{\tiny Ch.\,13}};
\node[thmnode,right=2.4cm of A8](B8){Future Dist.\ Optimality\\{\tiny Ch.\,15}};
\foreach \a/\b in {A1/B1,A2/B2,A3/B3,A4/B4,A5/B5,A6/B6,A7/B7,A8/B8}
  \draw[thmedge,DistBlue](\a)--(\b);
\draw[thmedge,DistGray,dashed](B3)--(B4);
\draw[thmedge,DistGray,dashed](B5)--(B6);
\draw[thmedge,DistGray,dashed](B7)--(B8);
\end{tikzpicture}
\end{center}\clearpage

% ---- Preface -----------------------------------------------
\chapter*{Preface}
\addcontentsline{toc}{chapter}{Preface}

\begin{epigraph}{The beginning of wisdom is to call things by
  their proper names. The beginning of science is to notice
  that many of the names were wrong.}{Anonymous}
\end{epigraph}

This book began from a suspicion that appeared so simple
it was initially difficult to take seriously: that many
of the most important concepts used throughout science,
philosophy, economics, biology, computation, and everyday
reasoning are grammatically misleading.

We speak primarily in nouns. We organise our descriptions
around objects, entities, categories, and things. Yet the
phenomena those nouns describe often turn out, upon closer
inspection, to be histories, processes, flows, repairs,
constraints, and trajectories. A mountain is a geological
process observed at a particular timescale. A species is a
reproductive process. A company is a collection of
coordinated histories. A memory is a recoverable trajectory.
A river is not a thing through which water passes; it is a
stable pattern constituted by the passage itself.

The intuition that processes precede objects is hardly new.
It appears repeatedly throughout the history of philosophy
\parencite{whitehead1929,bergson1896,rescher1996},
physics, biology, and systems theory. What seemed less
developed was a rigorous account of what \emph{replaces}
object ontology once objects are no longer treated as
fundamental.

The answer proposed here begins with a more primitive
concept than object, process, information, matter, energy,
or even observation. It begins with \textbf{distinction}.

Before a thing can be identified, it must first be
distinguished from something else \parencite{spencer-brown1969}.
Before information can exist, alternatives must exist
\parencite{shannon1948}. Before measurement can occur,
possibilities must be separated. Before categories can
be constructed, boundaries must be drawn
\parencite{bateson1972}.

Distinction is therefore treated throughout this book not
as a derivative notion but as an ontological primitive.

Once distinction is taken seriously, a second problem
immediately appears. \emph{Distinctions are fragile.}
They degrade, disperse, collapse, and disappear. This
motivates the concept of repair — not optimisation, not
continuation, but the restoration of distinction. Persistent
structure exists because repair exists \parencite{schrodinger1944}.

Yet repair itself is insufficient. Repair mechanisms degrade.
This motivates regeneration: preservation of repair capacity.
And regeneration is insufficient if the futures it generates
contract. This motivates admissibility and the central
geometric object of the book: the admissibility manifold
$\adm(t)$, whose volume $\Vol(\adm(t))$ is the primary
measure of a system's future possibility.

The deepest result developed in this work is therefore not
the preservation of distinctions, nor repair, nor
regeneration. It is the preservation of the capacity to
generate future distinctions. This culminates in the
Generative Admissibility Principle:
\[
  \frac{d}{dt}\Vol(\adm(t)) \ge 0.
\]
A trajectory is valuable insofar as it preserves or expands
the space of admissible future trajectories.

The book does not claim to derive morality from geometry.
It does not claim that admissibility solves every normative
question. What it does claim is narrower: any system that
systematically destroys its own future distinction-producing
capacity eventually undermines the conditions required for
its continued existence. This is not a moral axiom. It
emerges as a structural consequence of the framework.

The argument may be summarised in a single sequence:
\[
  \text{Distinction}\to\text{Information}\to\text{Entropy}
  \to\text{History}\to\text{Recoverability}\to\text{Repair}
  \to\text{Regeneration}\to\text{Reachability}
  \to\text{Admissibility}\to\text{Generativity}.
\]
Every chapter is ultimately an elaboration of one step
in that progression. Whether the framework succeeds or
fails should be judged not by any individual application
but by whether this progression captures something real
about the structure of persistence, knowledge, intelligence,
civilisation, and the universe itself.

\tableofcontents
\listoffigures

% ============================================================
\mainmatter
% ============================================================

% ============================================================
\part{Foundations: Distinction Before Object}
% ============================================================

% ============================================================
%  FILE: chapters/ch01_distinction.tex
% ============================================================

\chapter{The Problem of Difference}
\label{ch:distinction}

\begin{epigraph}{Draw a distinction.}%
  {G. Spencer-Brown, \textit{Laws of Form}, 1969
  \parencite{spencer-brown1969}}
\end{epigraph}

\begin{objectives}
  \item Understand why distinction is prior to objecthood,
        information, and entropy.
  \item State and prove the Axiom of Distinction.
  \item Define distinction cost, partition refinement, and
        blind spots.
  \item Prove the First Distinction Theorem.
  \item Relate these ideas to hash functions (CS) and
        cell membranes (Bio).
\end{objectives}

\section{Why Difference Must Come First}
\label{sec:ch1-priority}

The purpose of this chapter is not to introduce distinction
as one concept among others, but to show that distinction
is the condition under which any further concept can be
specified at all \parencite{spencer-brown1969}.

Before a thing can be identified, it must first be
distinguished from something else. Before information can
exist, alternatives must exist \parencite{bateson1972}.
Before measurement can occur, possibilities must be
separated. Before categories can be constructed, boundaries
must be drawn \parencite{luhmann1984}.

The classical metaphysical question asks what exists. The
present framework asks what must be presupposed for anything
to be counted as existing in the first place. The answer
cannot be ``an object'' because objecthood already
presupposes a boundary. Nor can it be ``information''
because information already presupposes alternatives. Nor
can it be ``state'' because a state is meaningful only
relative to states from which it is distinguished.

The first primitive cannot be a thing, a datum, or a state.
It must be the operation that makes thinghood, data, and
statehood available. That operation is distinction.

\section{Domains, Partitions, and Marks}
\label{sec:ch1-domain}

\begin{definition}{Domain}{domain}
  A \emph{domain} is a non-empty set $X$ whose elements
  are possible positions, configurations, events, states,
  histories, or appearances relative to some system of
  description. \index{domain}
\end{definition}

At this level of abstraction no assumption is made that
$X$ is physical, mental, computational, or semantic.
It may be a space of molecular configurations, a set of
bitstrings, a population of organisms, or a set of
candidate theories.

\begin{definition}{Distinction}{distinction}
  A \emph{distinction} on domain $X$ is a partition
  \[
    \partition = \{P_i\}_{i \in I}
  \]
  of $X$ into non-empty, pairwise disjoint subsets
  satisfying $\bigcup_{i \in I} P_i = X$ and
  $P_i \cap P_j = \emptyset$ for $i \ne j$.
  Each $P_i$ is a \emph{cell} of the distinction.
  \index{distinction|textbf}
\end{definition}

A distinction does not add material to the domain. It
changes the descriptive structure by making some
differences available and leaving others unavailable.
The same underlying domain can support many incompatible
distinctions.

\begin{definition}{Marked and Unmarked Regions}%
  {marked-unmarked}
  Given a binary distinction $\partition = \{M, U\}$,
  the cell $M$ is the \emph{marked region} and
  $U = X \setminus M$ is the \emph{unmarked region}.
\end{definition}

The marked/unmarked terminology is useful because
observation typically proceeds by indication. The unmarked
region is not nothing — it is precisely what allows the
mark to function as a mark.

\section{The Axiom of Distinction}
\label{sec:ch1-axiom}

\begin{axiom}{Axiom of Distinction}{dist-axiom}
  Every nontrivial observation presupposes a distinction.
  Objects, information, cost, and blind spots arise
  simultaneously and necessarily from the act of
  partitioning a domain $X$.
  \index{distinction!Axiom of|textbf}
\end{axiom}

The claim is axiomatic in the sense that denying it
requires performing it. To deny that observation
presupposes distinction, one must distinguish the denial
from its alternatives. The operation is pragmatically
unavoidable.

\begin{theorem}{Observational Nontriviality Theorem}%
  {observational-nontriviality}
  Let $O: X \to Y$ be an observation map. If $O$ is
  nontrivial, meaning $\exists\, x_1, x_2 \in X$ with
  $O(x_1) \ne O(x_2)$, then $O$ induces a nontrivial
  distinction on $X$.
\end{theorem}

\begin{proof}
  Define $x \sim_O x'$ iff $O(x) = O(x')$. This is
  reflexive, symmetric, and transitive, hence an
  equivalence relation partitioning $X$ into classes
  $[x]_O = \{x' \in X : O(x') = O(x)\}$.
  Since $O$ is nontrivial, $\exists\, x_1, x_2$ with
  $O(x_1) \ne O(x_2)$, so $[x_1]_O \ne [x_2]_O$.
  The induced partition has at least two cells and is
  nontrivial.
\end{proof}

\begin{remark}{Observation produces, not receives}{obs-produces}
  This theorem is the first formal bridge between
  observation and distinction. An observation is not a
  passive reception of a pre-given object. It is a map
  that identifies some states while separating others.
  Observation therefore produces a partition, and whatever
  is called an object appears only after this partition
  has been imposed.
\end{remark}

\section{Objects as Equivalence Classes}
\label{sec:ch1-objects}

The ordinary grammar of experience suggests that objects
come first and distinctions are drawn between them. The
formal structure reverses this order.

\begin{definition}{Object Induced by a Distinction}%
  {object-induced}
  Let $\partition$ be a distinction on $X$.
  An \emph{object} relative to $\partition$ is a cell
  $P_i \in \partition$, or more generally an equivalence
  class of the relation
  $x \sim_\partition y \iff \exists\, P_i \in \partition$
  such that $x, y \in P_i$.
  \index{object!as equivalence class}
\end{definition}

\begin{theorem}{Object-as-Equivalence-Class Theorem}%
  {object-equivalence}
  Every object determined by observation map $O: X \to Y$
  is an equivalence class of the distinction induced by $O$.
\end{theorem}

\begin{proof}
  By the Observational Nontriviality Theorem, $O$ induces
  equivalence relation $\sim_O$ on $X$. The equivalence
  classes are $[x]_O = \{x' \in X : O(x') = O(x)\}$.
  Any object available through $O$ must correspond to one
  such class, because states outside the class yield
  different observations while states inside yield the
  same. Hence objecthood relative to $O$ is exactly
  equivalence-class structure.
\end{proof}

This result is the formal version of the claim that
objects are stabilised distinctions. Objecthood is not
denied — it is relativised to a distinction system.

\section{Refinement and Distinction Capacity}
\label{sec:ch1-refinement}

\begin{definition}{Refinement}{refinement}
  Partition $\partition_2$ \emph{refines} $\partition_1$,
  written $\partition_1 \preceq \partition_2$, if every
  cell of $\partition_2$ is contained in some cell of
  $\partition_1$. \index{refinement!of partitions|textbf}
\end{definition}

\begin{proposition}{Refinement Partial Order}{refinement-order}
  $\preceq$ is a partial order on the set of partitions
  of $X$.
\end{proposition}

\begin{proof}
  Reflexivity: $\partition \preceq \partition$ trivially.
  Antisymmetry: mutual refinement forces cell coincidence.
  Transitivity: if each cell of $\partition_3$ lies in a
  cell of $\partition_2$ and each cell of $\partition_2$
  in a cell of $\partition_1$, each cell of $\partition_3$
  lies in a cell of $\partition_1$.
\end{proof}

\begin{definition}{Distinction Capacity}{dist-capacity}
  For a finite partition $\partition$, the
  \emph{distinction capacity} is
  $D(\partition) = \log |\partition|$.
  If cells carry probabilities $p_i$, the weighted
  capacity is $D_p(\partition) = -\sum_i p_i \log p_i$.
  \index{distinction!capacity|textbf}
\end{definition}

\begin{theorem}{Refinement Increases Distinction Capacity}%
  {refinement-capacity}
  If $\partition_1 \preceq \partition_2$ (finite
  partitions), then $D(\partition_1) \le D(\partition_2)$.
\end{theorem}

\begin{proof}
  Refinement implies $|\partition_2| \ge |\partition_1|$.
  Since $\log$ is monotone increasing,
  $\log|\partition_1| \le \log|\partition_2|$.
\end{proof}

\section{Distinction Cost}
\label{sec:ch1-cost}

No real system draws distinctions for free. A cell
membrane requires energy \parencite{schrodinger1944}.
A measuring instrument requires resolution.
A database index requires storage.
A scientific distinction requires training and
instrumentation. The minimum energetic cost of
maintaining a distinction is set by Landauer's principle
\parencite{landauer1961irreversibility}.

\begin{definition}{Distinction Cost}{dist-cost}
  A \emph{distinction cost function} is a map
  $C: \Pi(X) \to \mathbb{R}_{\ge 0}$ satisfying
  $\partition_1 \preceq \partition_2 \implies
  C(\partition_1) \le C(\partition_2)$.
  \index{distinction!cost|textbf}
\end{definition}

\begin{theorem}{Refinement Cost Theorem}{refinement-cost}
  $\Delta C(\partition_1, \partition_2)
  = C(\partition_2) - C(\partition_1) \ge 0$
  whenever $\partition_1 \preceq \partition_2$.
\end{theorem}

\begin{proof}
  Direct from the defining monotonicity of $C$.
\end{proof}

\section{Blind Spots}
\label{sec:ch1-blind}

Every distinction illuminates some differences by
suppressing others. The act of partitioning creates blind
spots. This is not a psychological defect but a
structural fact.

\begin{definition}{Projection Induced by a Distinction}%
  {projection-induced}
  Given partition $\partition = \{P_i\}_{i \in I}$,
  define the projection $\pi_\partition: X \to I$ by
  $\pi_\partition(x) = i \iff x \in P_i$.
\end{definition}

\begin{definition}{Blind Spot}{blind-spot}
  The \emph{blind spot} of projection $\pi: X \to Y$ is
  $\ker(\pi) = \{(x,x') \in X \times X : \pi(x) = \pi(x')\}$.
  Any pair $(x,x') \in \ker(\pi)$ with $x \ne x'$ is
  invisible to the projection.
  \index{blind spot|textbf}
\end{definition}

\begin{theorem}{Blind Spot Theorem}{blind-spot-thm}
  Let $\pi: X \to Y$ be surjective with $|X| > |Y|$
  for finite $X, Y$. Then $\pi$ has a nontrivial
  blind spot.
\end{theorem}

\begin{proof}
  Since $\pi$ is surjective, each $y \in Y$ has a
  non-empty fibre $\pi^{-1}(y)$. If every fibre had
  exactly one element, $|X| = |Y|$, contradicting
  $|X| > |Y|$. Hence by the pigeonhole principle at
  least one fibre contains two distinct elements
  $x \ne x'$ with $\pi(x) = \pi(x')$.
\end{proof}

\begin{definition}{Projection Loss}{projection-loss}
  The \emph{projection loss} of $\pi: X \to Y$ relative
  to reference partition $\partition_X$ is
  $L_\pi = D(\partition_X) - D(\pi(\partition_X))$.
  \index{projection!loss}
\end{definition}

\begin{theorem}{Projection Loss Theorem}{projection-loss-thm}
  $L_\pi \ge 0$. Moreover $L_\pi = 0$ iff $\pi$
  preserves all distinctions in $\partition_X$.
\end{theorem}

\begin{proof}
  Projection cannot increase $|\pi(\partition_X)|$, since
  cells may only be merged, not split. Hence
  $D(\pi(\partition_X)) \le D(\partition_X)$ and
  $L_\pi \ge 0$. Equality holds iff $|\pi(\partition_X)|
  = |\partition_X|$, i.e.\ no cell is collapsed.
\end{proof}

\section{The First Distinction Theorem}
\label{sec:ch1-fdt}

\begin{theorem}{First Distinction Theorem}{first-distinction}
  Any nontrivial partition $\partition$ of finite domain
  $X$ simultaneously induces:
  \begin{enumerate}[(i)]
    \item objects as cells of $\partition$;
    \item distinction capacity $D(\partition)$;
    \item maintenance cost $C(\partition)$;
    \item blind spots $\ker(\pi_\partition)$.
  \end{enumerate}
\end{theorem}

\begin{proof}
  (i)~Definition~\ref{defn:object-induced}.
  (ii)~Definition~\ref{defn:dist-capacity}.
  (iii)~Definition~\ref{defn:dist-cost}.
  (iv)~Definition~\ref{defn:blind-spot}
  and Theorem~\ref{thm:blind-spot-thm}.
  All four arise from the single operation of partitioning.
\end{proof}

\begin{remark}{The structural shadow of observation}%
  {structural-shadow}
  The First Distinction Theorem supplies the formal
  foundation for the rest of the book. Entropy will be
  introduced as transformation and erosion of distinction
  structures. History will be introduced as the persistence
  of distinctions through time. Recoverability will measure
  whether dispersed distinctions remain reconstructible.
  Repair will describe the operations by which damaged
  distinctions are restored. Regeneration will describe
  systems that preserve the capacity for future repair.
  Admissibility will evaluate which futures preserve the
  possibility of further distinction-production.
\end{remark}

\section*{Chapter Summary}
\begin{chapsummary}
  \item Distinction is prior to objecthood, information,
        and entropy (\cref{axm:dist-axiom}).
  \item Nontrivial observation requires nontrivial
        distinction
        (\cref{thm:observational-nontriviality}).
  \item Objects are equivalence classes induced by
        distinctions (\cref{thm:object-equivalence}).
  \item Refinement increases capacity at nonnegative cost
        (\cref{thm:refinement-capacity,thm:refinement-cost}).
  \item Every surjective projection has a nontrivial
        blind spot (\cref{thm:blind-spot-thm}).
  \item Object, capacity, cost, and blind spot are
        co-produced by the act of partitioning
        (\cref{thm:first-distinction}).
\end{chapsummary}

\subsection*{Exercises}

\begin{exercise}[CS]
  Let $X = \{0,1\}^8$ and let $\pi(x)$ return the parity
  bit of $x$. Describe the equivalence classes of
  $\sim_\pi$, compute $|\pi^{-1}(y)|$ for each $y$,
  and interpret the blind spot. What information is
  irrecoverably lost?
\end{exercise}

\begin{exercise}[Bio]
  A cell membrane separates interior from exterior.
  Identify the distinction being maintained, its cost
  (in ATP), and one biological repair mechanism that
  restores it when breached.
\end{exercise}

\begin{exercise}
  Prove that if $\partition_1 \preceq \partition_2
  \preceq \partition_3$ then
  $D(\partition_1) \le D(\partition_2) \le D(\partition_3)$.
\end{exercise}

\begin{exercise}
  Give an example of a projection with high distinction
  capacity but a large blind spot. Why does high capacity
  alone not imply adequacy of representation?
\end{exercise}

\begin{exercise}[CS]
  Interpret a cryptographic hash function $h: X \to Y$
  (with $|Y| \ll |X|$) as a projection. What is its
  blind spot? Under what conditions does collision
  resistance reduce but not eliminate projection loss?
\end{exercise}

% ============================================================
%  FILE: chapters/ch02_information.tex
% ============================================================

\chapter{Distinction and Information}
\label{ch:information}

\begin{epigraph}{Information is a difference that makes
  a difference.}%
  {Gregory Bateson \parencite{bateson1972}}
\end{epigraph}

\begin{objectives}
  \item Derive Shannon entropy from partition structure.
  \item Prove the Information--Distinction Theorem.
  \item Define projection loss and representational entropy.
  \item Prove the Representation Limit Theorem.
  \item Connect distinguishability geometry to sequence
        alignment (Bio) and hashing (CS).
\end{objectives}

\section{Information as Distinction Structure}
\label{sec:ch2-info}

Classical information theory begins with messages and
probabilities \parencite{shannon1948,cover2006}.
The present framework begins one layer beneath this,
asking what conditions must hold before a message can
be said to differ from another. The answer remains
distinction. \index{information!as distinction count|textbf}

A string of bits contains information only because
alternative strings are distinguishable. A genetic
sequence contains information only because alternative
nucleotide arrangements are distinguishable. A scientific
measurement contains information only because alternative
outcomes are distinguishable. If all possible outcomes
were identified under a single equivalence class,
information would vanish regardless of how much material
substrate remained.

\begin{definition}{Distinction Entropy (Information)}{dist-entropy-ch2}
  The \emph{distinction entropy} of partition $\partition$
  with cell probabilities $\{p_i\}$ is
  \[
    H(\partition) = -\sum_i p_i \log p_i.
  \]
  Under a uniform distribution,
  $H(\partition) = \log|\partition| = D(\partition)$.
\end{definition}

\begin{theorem}{Entropy--Distinction Duality}{entropy-duality}
  Let $\partition_1 \preceq \partition_2$ be finite
  partitions with uniform probabilities. Then
  $H(\partition_1) \le H(\partition_2)$.
\end{theorem}

\begin{proof}
  Under uniform distribution, $H(\partition) = \log|\partition|$.
  Since $\partition_1 \preceq \partition_2$ implies
  $|\partition_1| \le |\partition_2|$ and $\log$ is
  monotone, $H(\partition_1) \le H(\partition_2)$.
\end{proof}

\section{Information as Distinction Selection}
\label{sec:ch2-selection}

Suppose an observer knows only that a system lies
somewhere within partition $\partition = \{P_1,\ldots,P_n\}$.
Observation identifies one cell. The informational event
is the reduction of uncertainty over cells:
$I(P_i) = -\log p_i$ \parencite{shannon1948,mackay2003}.

This is Shannon's surprisal, but interpreted
distinction-theoretically: it measures how strongly the
observation refines the distinction structure.

\section{The Compression Theorem}
\label{sec:ch2-compression}

\begin{theorem}{Compression Theorem}{compression}
  For surjective $\pi: X \to Y$, $H(Y) \le H(X)$.
\end{theorem}

\begin{proof}
  A surjective map merges fibres $\pi^{-1}(y)$.
  Merging states cannot increase the number of
  distinguishable outcomes. The induced partition on $Y$
  is coarser than on $X$. By the Entropy--Distinction
  Duality, $H(Y) \le H(X)$.
\end{proof}

\begin{remark}{Compression is controlled distinction loss}%
  {compression-loss}
  Every representation compresses. Every model compresses.
  Every theory compresses. The question is never whether
  compression occurs but which distinctions survive.
  Later chapters will show that many failures of science,
  governance, and computation arise when a system mistakes
  the absence of represented difference for the absence
  of real difference.
\end{remark}

\section{Representational Entropy}
\label{sec:ch2-rep}

\begin{definition}{Representational Entropy}{rep-entropy}
  Let $\pi: X \to M$ be a representation. The
  \emph{representational entropy} at $m$ is
  $S_\pi(m) = \log|\pi^{-1}(m)|$,
  measuring the size of the hidden region collapsed
  beneath the representation.
\end{definition}

Large $S_\pi$ indicates many indistinguishable underlying
states — high compression with correspondingly large
blind spot.

\section{The Representation Limit Theorem}
\label{sec:ch2-limit}

\begin{theorem}{Representation Limit Theorem}{rep-limit}
  No finite representation can preserve every distinction
  of an infinite domain.
\end{theorem}

\begin{proof}
  Let $|X| = \infty$, $|M| < \infty$. Any $\pi: X \to M$
  must assign infinitely many states to at least one
  fibre $\pi^{-1}(m)$ by the pigeonhole principle.
  Hence infinitely many distinctions are collapsed.
\end{proof}

This theorem is the first formal appearance of what
later chapters call projection failure. No bounded
representation can preserve all distinctions of an
unbounded reality. Loss is not an implementation flaw
— it is a mathematical necessity.

\section{The Fundamental Information Theorem}
\label{sec:ch2-fund}

\begin{theorem}{Fundamental Information Theorem}{fund-info}
  Every informational quantity can be expressed as a
  property of a distinction structure.
  \index{information!fundamental theorem}
\end{theorem}

\begin{proof}
  Information requires distinguishable alternatives, which
  define a partition. Probability distributions are defined
  over partition cells. Entropy is computed from those
  probabilities. Mutual information measures shared
  refinement relations between partitions. Compression
  measures distinction collapse. Channel capacity measures
  transmissible distinctions. Each quantity depends solely
  on partition structure; hence each is a property of
  distinctions.
\end{proof}

\begin{remark}{The inverted hierarchy}{inverted-hierarchy}
  The traditional hierarchy places objects first and
  derives information from them. The Fundamental
  Information Theorem inverts this:
  $\text{Distinctions} \to \text{Information} \to \text{Objects}$.
  Objects are informational artefacts generated by
  distinction-preserving compressions.
\end{remark}

\section*{Chapter Summary}
\begin{chapsummary}
  \item Information is not primitive; it is the
        quantitative expression of distinctions
        (\cref{thm:fund-info}).
  \item Refinement increases entropy; coarsening decreases
        it (\cref{thm:entropy-duality}).
  \item Compression destroys distinctions
        (\cref{thm:compression}).
  \item No finite representation preserves all distinctions
        of an infinite domain (\cref{thm:rep-limit}).
  \item The correct hierarchy is Distinction $\to$
        Information $\to$ Objects.
\end{chapsummary}

\subsection*{Exercises}

\begin{exercise}[CS]
  Show that no lossless compression algorithm can reduce
  average description length below $H(X)$. What does this
  say about distinction capacity?
\end{exercise}

\begin{exercise}[Bio]
  Treat each DNA codon (triplet) as a partition cell over
  the 64 possible codons. Compute $H(\partition)$ assuming
  uniform usage. Interpret the result as genetic
  distinction capacity. How does degeneracy in the
  genetic code affect this calculation?
\end{exercise}

\begin{exercise}
  Prove that mutual information $I(X;Y) = H(X) - H(X|Y)$
  measures the refinement that $Y$ induces on the
  partition of $X$. Interpret this geometrically.
\end{exercise}

\begin{exercise}[CS]
  A lossy image compression algorithm reduces a 24-bit
  RGB image to an 8-bit palette. Model this as a
  projection $\pi: X \to M$. Compute the maximum
  possible projection loss. Under what conditions is
  the loss perceptually acceptable despite being
  informationally large?
\end{exercise}

% ============================================================
%  FILE: chapters/ch03_entropy.tex
% ============================================================

\chapter{Distinction and Entropy}
\label{ch:entropy}

\begin{epigraph}{The increase of entropy is not a law of
  chaos. It is a law about the progressive loss of
  recoverable structure.}{Author}
\end{epigraph}

\begin{objectives}
  \item Reinterpret entropy as distinction multiplicity.
  \item Prove the Distinction--Entropy Duality as a law.
  \item Derive the Second Law as a distinction-erosion
        theorem.
  \item Distinguish local order from global entropy growth.
  \item Identify entropy-increasing processes in genomes
        and networks.
\end{objectives}

\section{The Traditional Interpretation and Its Limits}
\label{sec:ch3-trad}

Entropy occupies a peculiar position in scientific
thought. In thermodynamics it is associated with heat
and irreversibility \parencite{boltzmann1877,gibbs1902}.
In statistical mechanics it is associated with
multiplicity. In information theory it is associated with
uncertainty \parencite{jaynes1957}. In cosmology it is
associated with the arrow of time. In computation it is
associated with erasure \parencite{landauer1961irreversibility}.

Each interpretation captures something important. Yet all
appear to describe different phenomena. The central claim
of this chapter is that these interpretations arise from
a common geometric structure: \emph{entropy measures
latent indistinguishability hidden beneath a distinction
structure}.

Entropy is not fundamentally disorder. Entropy is not
fundamentally randomness. Entropy is the multiplicity
concealed beneath a distinction.

\section{Distinction Spaces and Hidden Multiplicity}
\label{sec:ch3-hidden}

\begin{definition}{Hidden Multiplicity}{hidden-mult}
  Let $P_i \in \partition$ be a partition cell.
  The \emph{hidden multiplicity} of $P_i$ is
  $\Omega(P_i) = |P_i|$: the number of underlying
  states collapsed beneath the distinction.
  \index{hidden multiplicity}
\end{definition}

\begin{definition}{Distinction Entropy (Statistical)}%
  {dist-entropy-b}
  For partition cell $P_i$,
  $S(P_i) = k \log \Omega(P_i)$
  \parencite{boltzmann1877}.
  This measures the logarithmic volume of states
  collapsed beneath an observation.
\end{definition}

\section{The Distinction--Entropy Duality}
\label{sec:ch3-duality}

\begin{law}{Distinction--Entropy Duality}{entropy-dual}
  Every distinction simultaneously produces information
  and entropy. Information measures separation between
  partition cells. Entropy measures multiplicity within
  partition cells. Increasing one generally decreases the
  other. \index{entropy!as distinction multiplicity|textbf}
\end{law}

\begin{proof}
  Let $\partition = \{P_1,\ldots,P_n\}$. The number of
  distinguishable cells $N_D = |\partition|$ contributes
  to information: $H = \log N_D$. The multiplicity within
  each cell $\Omega(P_i) = |P_i|$ contributes to entropy:
  $S(P_i) = k\log\Omega(P_i)$. Both arise from the same
  partition; they are dual descriptions of the same
  distinction structure viewed from opposite directions.
  Refinement increases $N_D$ while reducing average
  $\Omega(P_i)$; coarsening does the reverse.
\end{proof}

\section{Entropy Production as Distinction Erosion}
\label{sec:ch3-production}

\begin{theorem}{Entropy Production Theorem}{entropy-production}
  Coarsening of distinction structure increases hidden
  multiplicity and therefore increases entropy.
\end{theorem}

\begin{proof}
  Suppose cells $P_i$ and $P_j$ are merged. Before
  merging: $\Omega_i = |P_i|$, $\Omega_j = |P_j|$.
  After merging: $\Omega' = |P_i| + |P_j|$.
  Since $a + b > \max(a,b)$, we have $\log\Omega' >
  \log\Omega_i$ and $\log\Omega' > \log\Omega_j$.
  Entropy increases in both affected cells.
\end{proof}

This gives the distinction interpretation of the Second
Law \parencite{prigogine1984}: entropy increases because
distinction structures become progressively less specific.

\section{The Geometry of Irreversibility}
\label{sec:ch3-irrev}

\begin{theorem}{Irreversibility Theorem}{irreversibility}
  Irreversibility emerges whenever history projection
  $\pi_H: \hist \to X$ is many-to-one.
\end{theorem}

\begin{proof}
  Let $h_1 \ne h_2$ with $\pi_H(h_1) = \pi_H(h_2) = x$.
  Knowledge of the present state $x$ cannot determine
  which history occurred. The inverse map $\pi_H^{-1}$
  is not unique; evolution toward the past is ambiguous.
  Hence macroscopic irreversibility emerges from the
  many-to-one character of history projection.
\end{proof}

\section{Local Order and Global Entropy}
\label{sec:ch3-local}

\begin{theorem}{Local--Global Compatibility Theorem}%
  {local-global}
  Local distinction concentration is compatible with
  global entropy growth.
\end{theorem}

\begin{proof}
  Let subsystem $\Sigma$ decrease its entropy by
  $|\Delta S_\Sigma|$. The Second Law requires
  $\Delta S_\mathrm{total} \ge 0$, so the environment
  must absorb $\Delta S_E \ge |\Delta S_\Sigma|$.
  Local order formation is exactly balanced or exceeded
  by global disorder export \parencite{schrodinger1944}.
\end{proof}

\begin{remark}{Life, complexity, and entropy}%
  {life-complexity}
  This theorem explains why galaxies, organisms,
  languages, and civilisations can form and persist.
  They are local concentrations of distinction density
  sustained by exporting entropy to their environments.
  Repair and regeneration are thermodynamically costly
  precisely because they maintain local order against
  the global tendency toward distinction erosion.
\end{remark}

\section{The Entropy--Reachability Connection}
\label{sec:ch3-reach}

\begin{theorem}{Entropy--Reachability Relation}%
  {entropy-reach-pre}
  If distinction capacity decreases monotonically,
  effective reachability volume decreases monotonically.
\end{theorem}

\begin{proof}
  Each reachable future state corresponds to a distinct
  future distinction structure. If partition cells merge,
  previously distinct future trajectories become
  observationally equivalent. The number of effective
  futures contracts. Hence reachability volume decreases
  with distinction capacity.
\end{proof}

This theorem foreshadows the central argument of
Chapters~\ref{ch:reachability}--\ref{ch:geometry-futures}.
Entropy is not merely about disorder; it is about future
possibility.

\section{The Second Law Reinterpreted}
\label{sec:ch3-second}

\begin{law}{Second Law of Distinction Dynamics}%
  {second-law-dist}
  In the absence of repair, distinction-generating
  dynamics, or external intervention, recoverable
  distinction capacity tends to decrease:
  $\frac{dD}{dt} \le 0$.
  \index{entropy!Second Law reinterpreted}
\end{law}

Equivalently, entropy increase is the shadow cast by
distinction loss. The law describes progressive movement
from fine distinction structures toward coarse ones —
erosion of recoverability, contraction of future
possibility, and accumulation of projection loss.

\section*{Chapter Summary}
\begin{chapsummary}
  \item Entropy measures latent multiplicity hidden beneath
        a distinction structure, not disorder
        (\cref{law:entropy-dual}).
  \item Coarsening increases entropy; refinement decreases
        it (\cref{thm:entropy-production}).
  \item Irreversibility arises from many-to-one history
        projection (\cref{thm:irreversibility}).
  \item Local order formation is compatible with global
        entropy growth (\cref{thm:local-global}).
  \item Entropy growth contracts future reachability
        (\cref{thm:entropy-reach-pre}).
  \item The Second Law is a theorem about distinction
        erosion (\cref{law:second-law-dist}).
\end{chapsummary}

\subsection*{Exercises}

\begin{exercise}[Bio]
  Mitochondria maintain a proton gradient across their
  inner membrane. Using distinction cost, explain what
  happens to local and global distinction density when
  the gradient collapses. Which cell of the partition
  expands? Which shrinks?
\end{exercise}

\begin{exercise}[CS]
  A cache replacement policy (LRU, LFU, random) can be
  thought of as a repair operator on temporal distinction
  structure. Compare LRU and random replacement in terms
  of distinction preservation. Which maintains higher
  $D$ for a given cache size?
\end{exercise}

\begin{exercise}
  Prove that the entropy of a mixture is greater than
  the weighted average entropy of its components:
  $H(\lambda P + (1-\lambda)Q) \ge
   \lambda H(P) + (1-\lambda)H(Q)$.
  Interpret this as a statement about distinction
  structure under mixing.
\end{exercise}

\begin{exercise}
  The Maxwell's demon thought experiment involves an
  agent that appears to decrease entropy without cost.
  Using Landauer's principle \parencite{landauer1961irreversibility}
  and the distinction cost framework, explain why the
  demon cannot violate the Second Law.
\end{exercise}

% ============================================================
%  FILE: chapters/ch04_history.tex
% ============================================================

% ============================================================
\part{Histories and Recoverability}
% ============================================================

\chapter{History Before State}
\label{ch:history}

\begin{epigraph}{The present does not contain its own
  explanation. Every present is the residue of many
  possible pasts.}{Author}
\end{epigraph}

\begin{objectives}
  \item Define history space $\hist(X)$ and the terminal
        projection $\pi_t$.
  \item Prove the History Primacy Theorem.
  \item Explain the Markov assumption as history
        compression.
  \item Prove the Information Loss Theorem.
  \item Relate path dependence to version control (CS)
        and phylogeny (Bio).
\end{objectives}

\section{The State Ontology Problem}
\label{sec:ch4-state}

Most scientific theories are state-based. Physics describes
positions and momenta. Economics describes prices and
inventories. Biology describes gene frequencies. Computer
science describes machine states. The common assumption is
that the state is fundamental: a system exists in a state,
the state changes, computation proceeds.

This picture is useful but incomplete. A state rarely
contains sufficient information to determine how it came
to exist. The present is generally compatible with many
possible histories \parencite{whitehead1929,bergson1896}.
Consequently the state cannot be the deepest ontological
object. The deeper object is the history from which the
state emerges.

\begin{definition}{History Space}{history-space}
  Let $X$ be a state space. The corresponding
  \emph{history space} is
  $\hist = \{\gamma : [t_0, t] \to X\}$,
  the collection of all trajectories into $X$.
  Elements of $\hist$ are trajectories rather than states.
  \index{history space|textbf}
\end{definition}

\begin{definition}{Terminal Projection}{terminal-proj}
  The \emph{terminal projection}
  $\pi_H: \hist \to X$ is defined by
  $\pi_H(\gamma) = \gamma(t)$.
  It assigns each history its terminal state.
\end{definition}

\section{States as Equivalence Classes of Histories}
\label{sec:ch4-equiv}

\begin{proposition}{States are Equivalence Classes of
  Histories}{states-as-equiv}
  Fix time $t$. Define $h_1 \sim_t h_2$ whenever
  $h_1(t) = h_2(t)$. This equivalence relation
  partitions $\hist$; each class is the set of histories
  terminating at the same state.
\end{proposition}

\begin{proof}
  Reflexivity: $h \sim_t h$ since $h(t) = h(t)$.
  Symmetry and transitivity follow immediately from
  equality of terminal states.
  The resulting partition assigns each state $x \in X$
  the equivalence class $[h]_t = \{h' : h'(t) = x\}$.
\end{proof}

This mirrors the result of Chapter~\ref{ch:distinction}:
objects are equivalence classes of distinctions; states are
equivalence classes of histories.

\section{The History Surplus}
\label{sec:ch4-surplus}

\begin{definition}{History Surplus}{history-surplus}
  The \emph{history surplus} of state $x$ at time $t$ is
  $\Omega_H(x,t) = |\pi_H^{-1}(x)|$,
  the number of histories compatible with the observed state.
  \index{history!entropy}
\end{definition}

States with large history surplus conceal substantial
historical structure. Historical entropy:
$S_H(x) = \log\Omega_H(x,t)$.

\section{The Law of Historical Compression}
\label{sec:ch4-law}

\begin{law}{Law of Historical Compression}{hist-compression}
  Every state description is a compression of a larger
  history description. History ontology strictly contains
  state ontology.
  \index{history!primacy theorem}
\end{law}

\begin{theorem}{History Primacy Theorem}{history-primacy}
  For any nontrivial dynamical system,
  $|\hist| > |X|$ and $\pi_H$ is many-to-one.
  State ontology is strictly contained in history ontology.
\end{theorem}

\begin{proof}
  Distinct histories $h_1 \ne h_2$ may terminate at the
  same state: $\pi_H(h_1) = \pi_H(h_2) = x$. This is
  possible whenever the dynamics are not invertible.
  Hence $\pi_H$ is non-injective and $|\hist| > |X|$.
  Every state $x = \pi_H(h)$ for at least one history,
  so $X = \pi_H(\hist)$, but $\hist \supsetneq X$ strictly.
  Given $H_n$ one reconstructs $s_n$, but given $s_n$
  alone the inverse is generally not unique.
\end{proof}

\begin{theorem}{Information Loss Theorem}{history-info-loss}
  For any non-injective $\pi_H$: $H(\hist) > H(X)$.
\end{theorem}

\begin{proof}
  Non-injective $\pi_H$ is a compression in the sense of
  Chapter~\ref{ch:information}. By the Compression Theorem
  (\cref{thm:compression}), $H(X) \le H(\hist)$, with
  strict inequality when fibres are non-trivial.
\end{proof}

\section{The Markov Assumption as Compression}
\label{sec:ch4-markov}

\begin{theorem}{Markov Compression Theorem}%
  {markov-compression}
  The Markov property is a projection of history space
  onto state space. A process is Markov iff the
  path-dependence index
  $P_D = I(H_t; x_{t+1} \mid x_t) = 0$.
\end{theorem}

\begin{proof}
  Full dynamics depend on histories
  $H_t = (x_0,\ldots,x_t)$. The Markov assumption
  replaces $H_t$ with $x_t$, which is the projection
  $\Pi: \hist \to X$. Since many histories share the
  same final state, $\Pi$ is many-to-one: a compression.
  $P_D = 0$ iff $P(x_{t+1} \mid H_t, x_t)
  = P(x_{t+1} \mid x_t)$, exactly the Markov condition.
\end{proof}

Markov systems are not history-free. They are
history-compressed. The compression introduces blind spots:
any path-dependent phenomenon is invisible to a Markov
model.

\section{Programs as Histories}
\label{sec:ch4-programs}

\csicon
Traditional computation represents program state $s_n$.
History-centred computation represents the full event
sequence $H_n = (e_1, \ldots, e_n)$.

\begin{theorem}{History Dominance Theorem (Programs)}%
  {history-dom-prog}
  Given complete history $H_n$, one can reconstruct $s_n$.
  Given only $s_n$, the inverse is generally not unique.
  Hence $H_n$ contains strictly more information than $s_n$.
\end{theorem}

\begin{proof}
  Define $F(H_n) = s_n$ (deterministic state derivation).
  $F^{-1}$ is not unique since many event sequences can
  produce the same state. $H_n$ contains strictly more
  information by the Information Loss Theorem.
\end{proof}

\begin{example}{Git as History-First Architecture}%
  {git-history}
  \csicon
  A git repository stores the full DAG of commits rather
  than only the current file tree. This is an explicit
  implementation of history-first ontology: the current
  state is a projection (the working tree), but the
  history (the commit graph) is preserved and addressable.
  \texttt{git reset --hard} is a history-collapsing
  operation: it discards history to reach a state,
  trading distinction capacity for simplicity.
\end{example}

\section*{Chapter Summary}
\begin{chapsummary}
  \item States are compressed projections of histories
        (\cref{law:hist-compression}).
  \item History ontology strictly contains state ontology
        (\cref{thm:history-primacy}).
  \item History projection necessarily loses information
        (\cref{thm:history-info-loss}).
  \item The Markov assumption is a history compression
        (\cref{thm:markov-compression}).
  \item Programs and databases are better understood as
        historical structures than as state machines.
\end{chapsummary}

\subsection*{Exercises}

\begin{exercise}[CS]
  Git stores a DAG of commits rather than only the current
  file tree. Identify the history space, the state
  projection, and the information lost by \texttt{git squash}.
  Is \texttt{git squash} an admissible operation in the
  sense of Chapter~\ref{ch:admissibility}?
\end{exercise}

\begin{exercise}[Bio]
  Epigenetic marks record environmental history beyond the
  DNA sequence. Model this as a history space and describe
  what the state projection (the DNA sequence alone) loses.
  Under what conditions does epigenetic history become
  irrecoverable?
\end{exercise}

\begin{exercise}
  Prove that for a reversible dynamical system
  (where $\Phi_{s,t}$ is a bijection for all $s \le t$),
  $\Omega_H(x,t)$ is constant. What does this say about
  the relationship between reversibility and historical
  entropy?
\end{exercise}

% ============================================================
%  FILE: chapters/ch05_recoverability.tex
% ============================================================

\chapter{Recoverability}
\label{ch:recoverability}

\begin{epigraph}{Nothing is ever truly lost at the moment
  it disappears. The deeper question is whether enough
  traces remain for it to be reconstructed.}{Author}
\end{epigraph}

\begin{objectives}
  \item Define recoverability as a property of distinction
        systems distinct from mere survival.
  \item Prove the Law of Recoverability.
  \item Prove the Recoverability Theorem.
  \item Prove the Semantic Horizon Theorem.
  \item Apply recoverability to error-correcting codes
        (CS) and immunological memory (Bio).
\end{objectives}

\section{Beyond Destruction}
\label{sec:ch5-beyond}

The previous chapter established that states are
compressed histories. Every observation discards historical
information. Every representation conceals distinctions.
Every projection produces ambiguity.

At first glance this seems to imply that once information
is lost, it is lost forever. Yet this conclusion is false.

A burned manuscript may survive in quotations.
An extinct language may survive in inscriptions.
A damaged hard drive may preserve recoverable fragments.
A corrupted genome may be reconstructed from homologous
copies \parencite{kauffman1993}.
A scientific theory may be reconstructed from derivative
texts after the original documents vanish.

In each case the object itself disappears, yet sufficient
traces remain to permit reconstruction. The distinction
between disappearance and irrecoverability is therefore
fundamental.

\section{Destruction versus Dispersal}
\label{sec:ch5-dispersal}

There are at least two fundamentally different processes.
\emph{Destruction} removes information.
\emph{Dispersal} redistributes information.

A document shredded into pieces remains recoverable if
the pieces survive. A document burned to ash does not.
The first operation disperses distinctions. The second
annihilates them. Recoverability measures the difference.

\section{Formal Definitions}
\label{sec:ch5-formal}

\begin{definition}{Recoverability Space}{recover-space}
  A \emph{recoverability space} is a triple
  $(X, \mathcal{D}, \mathfrak{R})$ where $X$ is a domain,
  $\mathcal{D}$ is a collection of distinctions, and
  $\mathfrak{R}$ is a collection of reconstruction
  operators.
\end{definition}

\begin{definition}{Recoverability}{recoverability-defn}
  Let $d^*$ be a target distinction. The
  \emph{recoverability} of $d$ relative to $d^*$ is
  \[
    \reco(d)
    = \sup_{R \in \mathfrak{R}}
      \frac{I(R(d);\, d^*)}{I(d^*;\, d^*)}
    \;\in [0, 1].
  \]
  $\reco(d) = 1$ indicates perfect recoverability.
  $\reco(d) = 0$ indicates total irrecoverability.
  \index{recoverability|textbf}
\end{definition}

\section{The Law of Recoverability}
\label{sec:ch5-law}

\begin{theorem}{Law of Recoverability}{reco-law}
  Dispersal and destruction are not equivalent. A
  distinction possesses positive recoverability
  $\reco(d) > 0$ iff sufficient information survives
  in the system or its environment to constrain its
  reconstruction.
\end{theorem}

\begin{proof}
  ($\Rightarrow$) If $\reco(d) > 0$, then
  $\exists R \in \mathfrak{R}$ with $I(R(d); d^*) > 0$,
  meaning $R(d)$ carries nontrivial information about
  $d^*$. Hence sufficient information survives.
  ($\Leftarrow$) If sufficient information survives,
  define $R$ as the reconstruction operator using that
  information; then $I(R(d); d^*) > 0$, so $\reco > 0$.
\end{proof}

\section{The Recoverability Theorem}
\label{sec:ch5-recov}

\begin{theorem}{Recoverability Theorem}{recoverability-thm}
  A distinction is recoverable iff its generating history
  remains distinguishable from alternative histories.
\end{theorem}

\begin{proof}
  If the generating history is distinguishable from
  alternatives, a reconstruction operator can identify
  the correct historical trajectory; $\reco > 0$.
  If all generating histories are observationally
  indistinguishable, no reconstruction operator can
  uniquely identify the original; $\reco = 0$.
\end{proof}

\section{Redundancy and Recoverability}
\label{sec:ch5-redundancy}

Recoverability increases when information is duplicated.
DNA contains redundancy (complementary strands, repair
templates). Languages contain redundancy (synonymy,
context). Error-correcting codes contain redundancy
\parencite{mackay2003}. Scientific communities contain
redundancy (multiple research groups).

\begin{theorem}{Redundancy Theorem}{redundancy-thm}
  Recoverability increases monotonically with the number
  of independent copies $\rho$ of a distinction:
  $\partial\reco/\partial\rho \ge 0$.
\end{theorem}

\begin{proof}
  Each independent copy introduces an additional
  reconstruction pathway. Loss of one pathway does not
  eliminate others. The set of available reconstruction
  operators weakly expands; hence $\reco$ cannot decrease.
\end{proof}

\section{The Semantic Horizon}
\label{sec:ch5-horizon}

\begin{definition}{Semantic Horizon}{semantic-horizon-defn}
  The \emph{semantic horizon} is
  $\partial\mathcal{S} = \{d : \reco(d) = 0\}$,
  the set of distinctions whose recoverability has
  reached zero.
  \index{semantic horizon|textbf}
\end{definition}

\begin{theorem}{Semantic Horizon Theorem}{semantic-horizon}
  Beyond a critical projection depth, distinctions remain
  causally influential but are no longer directly
  reconstructible. The semantic horizon is the locus in
  history space where $\reco(d) \to 0$.
\end{theorem}

\begin{proof}
  Repeated projection merges distinction classes.
  At each stage information is discarded, inverse images
  grow, and $|\pi^{-1}(d)| \to \infty$.
  Recoverability $\reco \propto 1/|\pi^{-1}(d)| \to 0$.
  Yet causal effects persist because subsequent dynamics
  continue to depend upon earlier states even when those
  states are no longer reconstructible.
\end{proof}

\begin{remark}{The CMB as semantic horizon}{cmb-horizon}
  \physicon
  The cosmic microwave background is an astronomical
  semantic horizon: it records the state of the universe
  at recombination ($t \approx 380{,}000$ years), but
  much of the earlier structure has been thermally
  dispersed and is no longer directly recoverable
  \parencite{penrose2010}. It remains causally influential
  (it constrains all subsequent cosmic structure
  formation) without being fully reconstructible.
\end{remark}

\section*{Chapter Summary}
\begin{chapsummary}
  \item Loss and destruction are not equivalent:
        dispersal preserves recoverability
        (\cref{thm:reco-law}).
  \item Recoverability is determined by historical
        distinguishability (\cref{thm:recoverability-thm}).
  \item Redundancy increases recoverability monotonically
        (\cref{thm:redundancy-thm}).
  \item The semantic horizon marks where $\reco \to 0$;
        causal influence survives but reconstruction
        does not (\cref{thm:semantic-horizon}).
  \item Repair (Chapter~\ref{ch:repair}) is possible
        iff $\reco > 0$.
\end{chapsummary}

\subsection*{Exercises}

\begin{exercise}[CS]
  A Reed--Solomon $(255, 223)$ code can correct up to
  $t = 16$ symbol errors. Express this as a recoverability
  statement: what is $\reco(d)$ for a message with
  $k$ corrupted symbols, for $k = 0, 10, 16, 17$?
  What happens at the semantic horizon $k = 17$?
\end{exercise}

\begin{exercise}[Bio]
  Describe immunological memory as a recoverability
  mechanism. What is the distinction $d^*$ being made
  recoverable? What is the reconstruction operator
  (the immune response)? What event corresponds to
  crossing the semantic horizon (loss of memory)?
\end{exercise}

\begin{exercise}
  Prove that a system with $n$ independent identical
  copies of distinction $d$, each with individual
  recoverability $\rho_0 > 0$, has system-level
  recoverability that increases with $n$.
  What is the asymptotic behaviour as $n \to \infty$?
\end{exercise}

\begin{exercise}
  Interpret Bekenstein--Hawking black hole entropy
  $S_{BH} = kA/4\ell_P^2$ \parencite{bekenstein1973}
  as a recoverability statement. What is the semantic
  horizon in this context, and why does the information
  paradox arise as a recoverability problem?
\end{exercise}

% ============================================================
%  FILE: chapters/ch06_memory.tex
% ============================================================

\chapter{Memory and Reconstruction}
\label{ch:memory}

\begin{epigraph}{Memory is not the storage of the past.
  Memory is the preservation of the possibility of
  recovering the past.}{Author}
\end{epigraph}

\begin{objectives}
  \item Define memory as recoverability preservation.
  \item Prove the Reconstruction Theorem.
  \item Define ecphory and memory viscosity.
  \item Prove the Forgetting Theorem.
  \item Prove the Memory Conservation Law.
  \item Compare biological long-term potentiation (Bio)
        with distributed hash tables (CS).
\end{objectives}

\section{From Recoverability to Memory}
\label{sec:ch6-from}

The previous chapter established that recoverability is
the central quantity governing reconstruction. This
immediately raises a deeper question: what preserves
those traces?

Recoverability is a property. Memory is a mechanism.
Recoverability describes the possibility of reconstruction.
Memory provides the substrate that makes reconstruction
possible.

A civilisation may possess recoverable knowledge; its
archives constitute memory \parencite{tulving1983}.
A biological organism may possess recoverable developmental
information; its genome constitutes memory.
A computer may possess recoverable states; its storage
media constitute memory. Memory is therefore not a special
feature of minds. Wherever distinctions persist through
time in a form that permits future reconstruction,
memory exists.

\section{The Classical View and Its Limits}
\label{sec:ch6-classical}

Traditional treatments regard memory as storage:
information is placed into a container, the container
preserves it, and it is later retrieved \parencite{james1890}.
This metaphor is useful in engineering contexts but
ultimately inadequate.

Storage presupposes stable objects. Yet throughout this
book we have observed that objects are themselves
compressed histories. The storage metaphor explains memory
in terms of entities whose existence already depends upon
memory. The account becomes circular.

A more fundamental description treats memory not as
storage but as a relationship: between traces and
reconstruction operators, between present evidence and
historical distinctions.

\section{Memory as Recoverability Preservation}
\label{sec:ch6-memory}

\begin{definition}{Memory}{memory-defn}
  A \emph{memory system} is any process that preserves
  recoverability through time. For a distinction $d$
  that has ceased to be directly observable, a memory
  system $M$ ensures that $\reco(d, t + \Delta t) > 0$
  by maintaining traces from which $d$ remains
  reconstructible.
  \index{memory!as recoverability preservation|textbf}
\end{definition}

\begin{definition}{Memory Functional}{memory-functional}
  $\mathcal{M}(t) = \int_\mathcal{D} \reco(d, t)\,d\mu_D(d)$,
  the total recoverable distinction content of a system
  at time $t$.
\end{definition}

Large $\mathcal{M}$ indicates many reconstructible
distinctions. Small $\mathcal{M}$ indicates few.

\section{The Reconstruction Theorem}
\label{sec:ch6-reconstruction}

\begin{theorem}{Reconstruction Theorem}{reconstruction}
  Every memory system induces a family of reconstruction
  operators $\mathfrak{R}_M$ such that for all $d$ with
  $\reco(d) > 0$: $\exists R \in \mathfrak{R}_M$ with
  $R(M(d)) \approx d$ up to admissible distortion.
\end{theorem}

\begin{proof}
  By the Memory definition, $\reco(M(d)) > 0$.
  By the Law of Recoverability (\cref{thm:reco-law}),
  sufficient information survives to constrain
  reconstruction. The Repair Existence Theorem
  (\cref{thm:repair-exist}) then guarantees at least
  one reconstruction operator $R$ exists.
\end{proof}

Memory and reconstruction are dual: a memory without
reconstruction is operationally meaningless; a reconstruction
without memory is impossible.

\section{Ecphory}
\label{sec:ch6-ecphory}

A memory trace may exist without being actively accessible.
Reconstruction requires an appropriate cue.
\textcite{tulving1983} calls this process \emph{ecphory}.

\begin{definition}{Ecphory}{ecphory-defn}
  \emph{Ecphory} is the activation of a reconstruction
  operator by an admissible cue $c$:
  $R_c(M(d)) = d$.
  The cue selects a reconstruction pathway, not creates
  the memory. \index{ecphory|textbf}
\end{definition}

\begin{definition}{Memory Viscosity}{mem-viscosity}
  \emph{Memory viscosity} is $\mu_M = 1/|\mathfrak{R}_M|$:
  the reciprocal of the number of available reconstruction
  operators. High viscosity means difficult recall; low
  viscosity means easy reactivation.
\end{definition}

\section{Distributed Memory}
\label{sec:ch6-distributed}

\begin{definition}{Distributed Memory}{distributed-memory}
  Distributed memory exists across subsystems
  $S_1, \ldots, S_n$ whenever
  $\reco(d \mid S_1, \ldots, S_n)
   > \max_i \reco(d \mid S_i)$:
  the whole remembers more than any individual part.
\end{definition}

\begin{remark}{Neural and computational examples}%
  {neural-dist}
  \bioicon\csicon
  Biological memory is distributed: long-term potentiation
  (LTP) stores information across synaptic connections
  rather than in single neurons \parencite{dayan2001}.
  Distributed hash tables (DHTs) implement the same
  principle computationally: data is spread across nodes
  so that the loss of any single node does not eliminate
  recoverability. Both exploit the Redundancy Theorem
  (\cref{thm:redundancy-thm}) to achieve
  $\partial\reco/\partial n \ge 0$.
\end{remark}

\section{Forgetting}
\label{sec:ch6-forgetting}

\begin{theorem}{Forgetting Theorem}{forgetting}
  Every finite memory system possesses distinctions
  whose recoverability converges toward zero.
\end{theorem}

\begin{proof}
  Finite capacity imposes finite distinguishability
  (\cref{thm:rep-limit}). The number of possible
  distinctions grows faster than available storage
  capacity. Some distinctions must share memory
  representations; shared representations reduce
  recoverability. Repeated compression eventually drives
  some distinctions toward the semantic horizon.
\end{proof}

Forgetting is therefore not a failure but a geometric
necessity. The question is not whether forgetting occurs
but which distinctions are forgotten.

\section{The Memory Conservation Law}
\label{sec:ch6-conservation}

Recoverability may move between substrates. A biological
memory may become a written document; a written document
may become a digital archive.

\begin{theorem}{Memory Conservation Law}{mem-conservation}
  Under admissible transport, recoverable distinction
  content is conserved up to reconstruction error:
  $\mathcal{M}(t + \Delta t) \ge \mathcal{M}(t) -
  \epsilon(\Delta t)$, where $\epsilon(\Delta t) \to 0$
  as reconstruction fidelity improves.
\end{theorem}

\begin{proof}
  Admissible transport preserves reconstruction pathways
  (\cref{defn:admissible-repair}). Lossless transport
  preserves $\reco$ exactly. Lossy transport introduces
  bounded distortion $\epsilon$; hence $\mathcal{M}$
  changes by at most $\epsilon$.
\end{proof}

Memory therefore behaves less like stored matter and more
like conserved structure: it can migrate between substrates
while its functional content is maintained, provided the
transport is admissible.

\section{Memory and Repair}
\label{sec:ch6-repair}

The connection to Chapter~\ref{ch:repair} is now visible.
Repair requires reconstruction. Reconstruction requires
memory. Therefore repair depends fundamentally on memory.

A system incapable of remembering cannot repair itself:
it cannot identify deviations from the reference
distinction $d^*$, cannot compare present states to prior
states, and cannot reconstruct lost distinctions. Memory
is a prerequisite for repair.

The progression is now complete:
\begin{center}
Entropy destroys distinctions.\\
Recoverability measures whether they can return.\\
Memory preserves their recoverability.\\
Repair reconstructs them.\\
Regeneration expands the space within which future
distinctions may arise.
\end{center}

\section*{Chapter Summary}
\begin{chapsummary}
  \item Memory is recoverability preservation, not storage
        (\cref{defn:memory-defn}).
  \item Every memory system induces reconstruction
        operators (\cref{thm:reconstruction}).
  \item Ecphory is cue-driven activation of a
        reconstruction pathway (\cref{defn:ecphory-defn}).
  \item Distributed memory exploits redundancy to raise
        system-level recoverability.
  \item Finite memory systems inevitably forget
        (\cref{thm:forgetting}).
  \item Under admissible transport, recoverable content
        is conserved (\cref{thm:mem-conservation}).
  \item Memory is the prerequisite for repair.
\end{chapsummary}

\subsection*{Exercises}

\begin{exercise}[Bio]
  Long-term potentiation (LTP) strengthens synaptic
  connections through repeated activation. Model LTP
  as a memory mechanism: what is $d^*$, what is the
  reconstruction operator, and what corresponds to
  the semantic horizon (synaptic forgetting)?
\end{exercise}

\begin{exercise}[CS]
  A distributed hash table (DHT) with replication factor
  $k$ stores each key-value pair on $k$ independent nodes.
  Using the Redundancy Theorem (\cref{thm:redundancy-thm}),
  compute how $\reco$ changes as nodes fail.
  At what failure rate does the system cross the
  semantic horizon?
\end{exercise}

\begin{exercise}
  Prove that a write-only memory system (one that can
  store but not retrieve) has $\mathcal{I}(\Sigma) = 0$
  by the Intelligence--Repair Equivalence Theorem
  (\cref{thm:intel-repair}).
  What does this imply about the relationship between
  retrievability and intelligence?
\end{exercise}

\begin{exercise}
  The Memory Conservation Law states that admissible
  transport preserves recoverable content.
  Give an example of an inadmissible memory transport
  that violates this law, and identify which axiom of
  repair (\cref{defn:repair-op}) it violates.
\end{exercise}

% ============================================================
\part{Repair}
% ============================================================

% ---- Chapter 7 ---------------------------------------------
\chapter{Repair as a Fundamental Operation}
\label{ch:repair}
\begin{epigraph}{The art of living is more like wrestling
  than dancing, in so far as it stands ready against the
  accidental and the unforeseen.}%
  {Marcus Aurelius, \textit{Meditations} \parencite{aurelius}}
\end{epigraph}

\begin{principle}{Principle of Repair}{repair-prc}
  Persistent distinctions require repair. Repair is
  possible iff recoverability remains strictly positive.
  The composition of admissible repairs is admissible.
  \index{repair!operator|textbf}
\end{principle}

\begin{definition}{Repair Operator}{repair-op}
  $\repair:\mathcal{D}\to\mathcal{D}$ satisfying
  (R1)~$\delta(\repair(d),d^*)\le\delta(d,d^*)$,
  (R2)~$\repair(d^*)=d^*$,
  (R3)~$\repair(d)\ne d^*$ whenever $\reco(d)=0$.
\end{definition}

\begin{definition}{Admissible Repair}{admissible-repair}
  $\repair$ is \emph{admissible} if
  $V_R(\repair(d),t)\ge V_R(d,t)$ for all damaged $d$.
  \index{repair!admissible|textbf}
\end{definition}

\begin{theorem}{Repair Existence Theorem}{repair-exist}
  A repair operator satisfying (R1)--(R3) exists iff
  $\reco(d)>0$.
\end{theorem}
\begin{proof}
  ($\Rightarrow$) Any improving $\repair$ must use
  information recoverable from $d$; hence $\reco(d)>0$.
  ($\Leftarrow$) $\reco(d)>0$ gives reconstruction
  operator $\mathfrak{rec}$ from \cref{thm:reco-law};
  set $\repair=\mathfrak{rec}$; (R1)--(R3) follow.
\end{proof}

\begin{theorem}{Repair Closure Theorem}{repair-closure}
  Composition of two admissible repair operators is an
  admissible repair operator.
  \index{repair!monoid}
\end{theorem}
\begin{proof}
  (R1): $\delta((\repair_1\circ\repair_2)(d),d^*)
  \le\delta(\repair_2(d),d^*)\le\delta(d,d^*)$.
  (R2): $(\repair_1\circ\repair_2)(d^*)=d^*$.
  (R3): $\reco=0$ propagates through both operators.
  Admissibility: $V_R((\repair_1\circ\repair_2)(d),t)
  \ge V_R(\repair_2(d),t)\ge V_R(d,t)$.
\end{proof}

\begin{corollary}{Repair Monoid}{repair-monoid}
  Admissible repair operators form a monoid under
  composition with identity $\mathrm{id}_\mathcal{D}$.
\end{corollary}

\begin{theorem}{Minimal Repair Theorem}{minimal-repair}
  \textup{(Sketch)} Under compactness and continuity of
  the operator space, among all admissible repairs
  achieving $\delta(\repair(d),d^*)\le\epsilon$,
  there exists one minimising repair cost
  $\mathrm{Cost}(\repair,d)=
  \mu(\{y:\repair\text{ modifies distinction at }y\})$.
\end{theorem}

\begin{theorem}{Repair Conservation Law}{repair-conservation}
  Admissible repair preserves historical continuity:
  $d$ and $\repair(d)$ lie in the same connected
  component of the recoverability manifold $\mathcal{M}_\reco$.
\end{theorem}
\begin{proof}
  $\reco(d)>0$ means $d$ is not at the semantic horizon;
  $\mathfrak{rec}$ operates within $\mathcal{M}_\reco$,
  keeping $\repair(d)$ in the same component.
\end{proof}

\begin{theorem}{Repair--Entropy Theorem}{repair-entropy}
  For admissible repair of $\Sigma\subset X$ in
  $\Omega\supset\Sigma$:
  (i)~$\Delta S_\Sigma\le0$;
  (ii)~$\Delta S_{\Omega\setminus\Sigma}\ge|\Delta S_\Sigma|$;
  (iii)~$\Delta S_\Omega\ge0$.
\end{theorem}
\begin{proof}
  (i) Repair increases $D_\Sigma$, decreasing
  $S_\Sigma=k\log\Omega_\Sigma$.
  (ii) Second Law: $\Delta S_\Omega\ge0$ forces
  $\Delta S_{\Omega\setminus\Sigma}\ge-\Delta S_\Sigma$.
  (iii) Second Law directly \parencite{landauer1961irreversibility,
  bennett1982}.
\end{proof}

\section*{Chapter Summary}
\begin{chapsummary}
  \item Repair is the third primitive: irreducible to
        distinction and entropy.
  \item Repair exists iff $\reco>0$
        (\cref{thm:repair-exist}).
  \item Admissible repairs form a monoid
        (\cref{thm:repair-closure,cor:repair-monoid}).
  \item Repair is reconstruction, not creation
        (\cref{thm:repair-conservation}).
  \item Repair exports entropy globally
        (\cref{thm:repair-entropy}).
\end{chapsummary}

% ---- Chapter 8 ---------------------------------------------
\chapter{The Geometry of Failure}
\label{ch:geometry-failure}
\begin{epigraph}{It is the peculiarity of a problem of the
  first importance that it is almost insoluble except in
  hindsight.}%
  {Barbara Tuchman, \textit{The Guns of August}}
\end{epigraph}

\begin{definition}{Persistent Anomaly}{persistent-anomaly}
  Anomaly $o$ is \emph{persistent} relative to
  $(\repair,d^*)$ if
  $\lim_{n\to\infty}\delta_\mathrm{obs}(o,\repair^n(d^*))>0$.
  \index{anomaly!persistent|textbf}
\end{definition}

\begin{theorem}{Persistent Anomaly Theorem}%
  {persistent-anomaly-thm}
  $o$ is persistent iff
  $o\notin\overline{\repair(\mathcal{D})}$.
\end{theorem}
\begin{proof}
  ($\Rightarrow$) If $o\notin\overline{\repair(\mathcal{D})}$,
  some $\eta>0$ separates $o$ from all repaired
  distinctions; no repair sequence closes the gap.
  ($\Leftarrow$) If $o\in\overline{\repair(\mathcal{D})}$,
  a sequence $r_n\to o$ exists; $o$ is not persistent.
\end{proof}

\begin{definition}{Failure Manifold}{failure-manifold}
  $\mathcal{F}(d^*,\mathcal{R})=\{o\in\mathcal{O}:
  o\text{ persistent relative to }(\mathcal{R},d^*)\}$.
  \index{failure manifold|textbf}
\end{definition}

\begin{theorem}{Projection Failure Theorem}{proj-fail-thm}
  Every persistent anomaly is a projection failure:
  $\exists\tilde{X}\supset X$ in which $o$ is resolvable
  but not in $X$.
\end{theorem}
\begin{proof}
  Construct $\tilde{X}=X\times\{0,1\}$ with projection
  onto first component. The new coordinate distinguishes
  previously collapsed states; $o$ is resolved in
  $\tilde{X}$. \parencite{kuhn1962structure}
\end{proof}

\begin{definition}{Ontological Enlargement}{ont-enlarge}
  $(\tilde{\mathcal{D}},\tilde\delta)$ with embedding
  $\iota:\mathcal{D}\hookrightarrow\tilde{\mathcal{D}}$ is
  an enlargement if (i)~$\tilde\delta(\iota(d),\iota(d'))
  =\delta(d,d')$; (ii)~$\tilde{\mathcal{D}}\setminus
  \iota(\mathcal{D})\ne\emptyset$;
  (iii)~every persistent anomaly resolves in
  $\tilde{\mathcal{D}}$.
  \index{ontological enlargement|textbf}
\end{definition}

\begin{theorem}{Ontology Revision Theorem}{ont-revision}
  If $\mathcal{F}\ne\emptyset$, there exists a minimal
  ontological enlargement resolving all anomalies
  while preserving all existing distinctions.
\end{theorem}
\begin{proof}
  Define $\tilde{\mathcal{D}}=\mathcal{D}\cup
  \{d_o:o\in\mathcal{F}\}$; extend $\tilde\delta$ via
  a path metric; conditions (i)--(iii) hold by
  construction; minimality by taking the closure under
  the repair algebra.
\end{proof}

\begin{definition}{Repair Saturation}{repair-sat}
  The repair algebra is \emph{saturated} if
  $\frac{d}{dt}|\mathcal{F}|\ge0$ and
  $\frac{d}{dt}\mathrm{Cost}(\repair,d^*)\ge0$.
  \index{repair!saturation}
\end{definition}

\begin{theorem}{Scientific Revolution Theorem}{sci-rev-thm}
  A paradigm undergoes revolutionary change iff its repair
  algebra saturates and $\mathcal{F}$ contains a deep
  anomaly (low failure curvature).
  \parencite{kuhn1962structure,lakatos1978}
\end{theorem}
\begin{proof}
  Saturation plus deep anomaly makes local repair
  insufficient. \cref{thm:ont-revision} guarantees
  minimal enlargement, constituting a topological
  transition in the distinction manifold.
\end{proof}

\begin{corollary}{Ontology Revision is Generatively
  Admissible}{ont-rev-adm}
  Successful ontology revision increases $\Vol(\adm(t))$.
\end{corollary}

\section*{Chapter Summary}
\begin{chapsummary}
  \item Persistent anomalies lie outside
        $\overline{\repair(\mathcal{D})}$
        (\cref{thm:persistent-anomaly-thm}).
  \item Every persistent anomaly is a projection failure
        (\cref{thm:proj-fail-thm}).
  \item Minimal ontological enlargement always exists
        (\cref{thm:ont-revision}).
  \item Scientific revolutions require saturation plus deep
        anomaly (\cref{thm:sci-rev-thm}).
\end{chapsummary}

% ---- Chapter 9 ---------------------------------------------
\chapter{Repair and Intelligence}
\label{ch:repair-intelligence}
\begin{epigraph}{To understand is to perceive patterns.
  To learn is to find the repair.}%
  {Attributed to Isaiah Berlin, \textit{apocryphal}}
\end{epigraph}

\begin{definition}{Intelligence}{intelligence}
  $\mathcal{I}(\Sigma)
  =\kappa(\Sigma,\mathcal{D}_\Sigma)
  \cdot\kappa(\Sigma,\mathcal{D}_{\repair_\Sigma})$
  where $\kappa(\Sigma,\mathcal{D})
  =\sup_{d,\reco(d)>0}
  \eta(\repair_\Sigma,d)/\delta(d,d^*)$.
  \index{intelligence!as repair capacity|textbf}
\end{definition}

\begin{theorem}{Intelligence--Repair Equivalence}%
  {intel-repair}
  $\Sigma$ is more intelligent than $\Sigma'$ iff
  $\mathcal{I}(\Sigma)>\mathcal{I}(\Sigma')$.
\end{theorem}
\begin{proof}
  Superior generalisation, transfer, and adaptation all
  correspond to higher $\kappa$ over broader $\mathcal{D}$.
  Self-improvement requires
  $\kappa(\Sigma,\mathcal{D}_\repair)>0$.
  The converse follows symmetrically.
\end{proof}

\begin{theorem}{Prediction Limitation Theorem}{prediction-limit}
  $\delta(x_\mathrm{pred}(t+\tau),x(t+\tau))
  \ge(1-\reco_\tau(x(t)))\cdot
  \delta(x_\mathrm{worst},x(t+\tau))$.
\end{theorem}
\begin{proof}
  Repair recovers at most fraction $\reco_\tau$;
  the remainder is irrecoverable.
\end{proof}

\begin{definition}{Explanation}{explanation-defn}
  Explanation of $o$ is $d$ satisfying:
  (E1)~$\delta_\mathrm{obs}(o,d)<\epsilon$;
  (E2)~$d\in\adm(\mathcal{D},t)$;
  (E3)~$d$ reduces expected future repair cost.
\end{definition}

\begin{theorem}{Explanation Improvement Theorem}%
  {expl-improve}
  A good explanation (satisfying E1--E3) reduces expected
  future repair cost:
  $\mathcal{C}_\repair(t\mid d)\le
  \mathcal{C}_\repair(t\mid\neg d)$.
\end{theorem}

\begin{proposition}{Optimisation Cannot Resolve Persistent
  Anomalies}{opt-limit}
  Closed optimisation within $(\mathcal{D},\mathcal{R})$
  cannot resolve any $o\in\mathcal{F}$.
\end{proposition}
\begin{proof}
  By \cref{thm:persistent-anomaly-thm},
  $o\notin\overline{\repair(\mathcal{D})}$.
  Optimisation finds elements of $\mathcal{D}$; none
  resolves $o$.
\end{proof}

\begin{theorem}{Alignment Theorem}{alignment-thm}
  $\Sigma$ is aligned with ecology $\mathcal{E}$ iff
  $\frac{d}{dt}\Vol(\adm_\mathcal{E}(\Sigma(t),t))\ge0$.
  \index{alignment!theorem|textbf}
\end{theorem}
\begin{proof}
  Alignment means all repairs are admissible
  (\cref{prop:repair-adm}); the Future Volume Theorem
  then gives non-decreasing $\Vol(\adm_\mathcal{E})$,
  and conversely.
  \parencite{amodei2016,russell2019}
\end{proof}

\begin{corollary}{Capability Without Alignment}%
  {capability-misalign}
  High $\mathcal{I}(\Sigma)$ and alignment with
  $\mathcal{E}$ are logically independent.
  \index{alignment!and capability}
\end{corollary}

\section*{Chapter Summary}
\begin{chapsummary}
  \item Intelligence is repair capacity including
        self-repair (\cref{thm:intel-repair}).
  \item Prediction is bounded by future recoverability
        (\cref{thm:prediction-limit}).
  \item Good explanation reduces future repair cost~(E3).
  \item Optimisation cannot resolve persistent anomalies
        (\cref{prop:opt-limit}).
  \item Alignment is generative admissibility
        w.r.t.\ the ecology (\cref{thm:alignment-thm}).
  \item Intelligence and alignment are independent
        (\cref{cor:capability-misalign}).
\end{chapsummary}

% ============================================================
\part{Regeneration}
% ============================================================

% ---- Chapter 10 --------------------------------------------
\chapter{Beyond Continuation}
\label{ch:beyond-continuation}
\begin{epigraph}{The fact that something persists does not
  tell us whether it ought to persist.}{Author}
\end{epigraph}

\begin{theorem}{Continuation Equivalence Theorem}%
  {continuation-equivalence}
  Let $\gamma_1,\gamma_2\in\Gamma$ (both defined for all
  future times). Continuation alone cannot distinguish them.
\end{theorem}
\begin{proof}
  Continuation is binary; both trajectories satisfy the
  criterion equally. Any distinction requires additional
  structure.
\end{proof}

\begin{theorem}{Continuation Degeneracy Theorem}%
  {continuation-degeneracy-thm}
  There exist continuation trajectories in $\Gamma$ with
  $V_R(\gamma(t),t)\to0$ as $t\to\infty$.
  \index{continuation!degeneracy|textbf}
\end{theorem}
\begin{proof}
  Let $X=[0,1]$, $\gamma_2(t)$ occupying $[0,e^{-\lambda t}]$.
  Then $\gamma_2\in\Gamma$ but $V_R\to0$.
  Cancer, monopolies, and extractive economies exemplify
  this \parencite{ostrom1990,meadows2008}.
\end{proof}

\begin{proposition}{Regeneration Strictly Implies Continuation}%
  {regen-implies-cont}
  $\mathrm{Regeneration}\subsetneq\mathrm{Continuation}$.
\end{proposition}

\section*{Chapter Summary}
\begin{chapsummary}
  \item Continuation cannot distinguish flourishing from
        pathological persistence
        (\cref{thm:continuation-equivalence}).
  \item Degenerate continuations shrink $V_R$ to zero
        (\cref{thm:continuation-degeneracy-thm}).
  \item Regeneration is strictly stronger than continuation
        (\cref{prop:regen-implies-cont}).
\end{chapsummary}

% ---- Chapter 11 --------------------------------------------
\chapter{Regenerative Systems}
\label{ch:regeneration}
\begin{epigraph}{A thing that regenerates recreates the
  circumstances that permitted its survival.}{Author}
\end{epigraph}

\begin{definition}{Regeneration}{regeneration-def}
  System $\Sigma$ is \emph{regenerative} if
  $\frac{d}{dt}\kappa_\Sigma(t)\ge0$:
  repair capacity is preserved or expanded.
  \index{regeneration|textbf}
\end{definition}

\begin{principle}{Principle of Regeneration}{regen-prc}
  A system is regenerative iff it preserves the capacity
  to perform future repair.
\end{principle}

\begin{theorem}{Regeneration Theorem}{regeneration}
  $(X,\Phi)$ is regenerative iff
  $\kappa_\Sigma(t)\ge\kappa_\Sigma(t_0)$ for all
  $t\ge t_0$ and repair is continuously exercisable.
\end{theorem}
\begin{proof}
  ($\Rightarrow$) Definition gives non-decreasing capacity;
  capacity that cannot be exercised is operationally zero.
  ($\Leftarrow$) Non-decreasing exercisable capacity
  is exactly the definition.
  \parencite{holling1973,gunderson2002}
\end{proof}

\begin{theorem}{Regenerative Stability Theorem}{regen-stability}
  Sufficiently small perturbations do not produce permanent
  collapse of repair capacity in a regenerative system.
\end{theorem}
\begin{proof}
  $\frac{d}{dt}\kappa\ge0$ means small decreases generate
  compensatory repair dynamics; $\kappa$ remains bounded
  away from collapse.
\end{proof}

\begin{theorem}{Regenerative Expansion Theorem}{regen-expansion}
  If a regenerative system continuously encounters
  recoverable novelty, then
  $\frac{d}{dt}\kappa_\Sigma(t)>0$.
\end{theorem}
\begin{proof}
  Each recoverable novelty enlarges the repair algebra;
  enlargement increases future repairable distinctions;
  $\kappa$ grows strictly.
\end{proof}

\section*{Chapter Summary}
\begin{chapsummary}
  \item Regeneration is preservation of repair capacity,
        not mere continuation.
  \item Regenerative iff $\kappa_\Sigma$ non-decreasing
        with repair exercisable (\cref{thm:regeneration}).
  \item Regenerative systems are stable against small
        perturbations (\cref{thm:regen-stability}).
  \item Encountering novelty strictly expands capacity
        (\cref{thm:regen-expansion}).
\end{chapsummary}

% ---- Chapter 12 --------------------------------------------
\chapter{Distinction Ecology}
\label{ch:distinction-ecology}
\begin{epigraph}{No distinction exists alone.}{Author}
\end{epigraph}

\begin{definition}{Distinction Ecology}{dist-ecology}
  $\mathcal{E}=(\mathcal{D},\mathcal{L},\mathcal{R})$:
  distinctions, dependency relation, repair algebra.
  \index{distinction ecology|textbf}
\end{definition}

\begin{definition}{Distinction Dependency}{dist-dep}
  $d_j\prec d_i$ if loss of $d_j$ reduces $\reco(d_i)$.
\end{definition}

\begin{theorem}{Distinction Dependency Theorem}{dist-dependency}
  If $d_j\prec d_i$ and $\reco(d_j)=0$, then
  $\reco(d_i)\le\reco(d_i\mid d_j)$.
\end{theorem}
\begin{proof}
  Recovery of $d_i$ requires information in its dependency
  structure; $\reco(d_j)=0$ removes that information.
\end{proof}

\begin{theorem}{Ecological Fragility Theorem}{eco-fragility}
  If a node supports fraction $p$ of dependency paths,
  expected cascade size grows monotonically with $p$.
  \index{diversity!ecological|textbf}
\end{theorem}
\begin{proof}
  Fraction $p$ of all reconstruction routes pass through
  the node; failure simultaneously reduces recoverability
  along all those paths. \parencite{may1973,levin1998}
\end{proof}

\begin{theorem}{Diversity--Repair Theorem}{diversity-repair}
  Repair capacity increases weakly with the number of
  independent repair pathways:
  $\partial\kappa/\partial N\ge0$.
\end{theorem}
\begin{proof}
  Each independent pathway is an additional reconstruction
  route; loss of one does not eliminate others.
  \parencite{tilman1996}
\end{proof}

\begin{theorem}{Regenerative Ecology Theorem}{regen-ecology-thm}
  A regenerative ecology expands the set of future
  repairable distinctions:
  $\mathcal{D}_\mathrm{repair}(t_2)
  \supseteq\mathcal{D}_\mathrm{repair}(t_1)$
  for $t_2>t_1$.
\end{theorem}

\section*{Chapter Summary}
\begin{chapsummary}
  \item Distinctions form ecologies with dependency
        structures and repair algebras.
  \item Loss of foundational distinctions cascades
        (\cref{thm:dist-dependency}).
  \item Concentrated dependency increases fragility
        (\cref{thm:eco-fragility}).
  \item Diversity increases repair capacity
        (\cref{thm:diversity-repair}).
  \item Regenerative ecologies expand the repairable domain
        (\cref{thm:regen-ecology-thm}).
\end{chapsummary}

% ============================================================
\part{Geometry}
% ============================================================

% ---- Chapter 13 --------------------------------------------
\chapter{Reachability}
\label{ch:reachability}
\begin{epigraph}{The possible is not the actual.
  But it is the space from which the actual is drawn.}%
  {Modal realism, broadly construed}
\end{epigraph}

\begin{definition}{Reachability Set and Volume}{reach-set}
  $\reach(x,t_0,t)=\{y\in X:\exists\text{ admissible }
  \gamma,\;\gamma(t_0)=x,\;\gamma(t)=y\}$.
  $V_R(x,t)=\mu(\reach(x,t))$.
  \index{reachability!volume|textbf}
\end{definition}

\begin{theorem}{Reachability Monotonicity Theorem}{reach-mono}
  For $\mathcal{C}_1\subseteq\mathcal{C}_2$:
  $V_R(x,t\mid\mathcal{C}_2)\le V_R(x,t\mid\mathcal{C}_1)$.
\end{theorem}
\begin{proof}
  $\reach(\cdot\mid\mathcal{C}_2)\subseteq
  \reach(\cdot\mid\mathcal{C}_1)$; apply $\mu$.
  \parencite{sontag1998}
\end{proof}

\begin{theorem}{Constraint Volume Theorem}{constraint-vol}
  Under independence:
  $V_R(x,t)\le\mu(X)e^{-\rho_\mathcal{C}|\mathcal{C}|}$.
\end{theorem}
\begin{proof}
  Each constraint independently eliminates fraction
  $\rho_\mathcal{C}$; iterate and bound by exponential.
\end{proof}

\begin{theorem}{Boundary Proximity and Critical Instability}%
  {boundary-instab}
  As $\beta(y,x,t)=d(y,\partial\reach(x,t))\to0$,
  $\|\partial V_R/\partial\epsilon\|\to\infty$.
\end{theorem}

\begin{theorem}{Future Cone Theorem}{future-cone}
  $\mathcal{F}(x,t_0)=\bigcup_{t\ge t_0}\reach(x,t_0,t)$
  is non-empty, expanding under non-contracting dynamics,
  with codimension-one boundary when $X$ and $\Phi$ are
  smooth.
  \index{future cone|textbf}
\end{theorem}

\begin{theorem}{Reachability as Distinction Preservation}%
  {reach-dist}
  $V_R(\gamma(t),t)\ge V_R(\gamma(t_0),t_0)$ iff
  $D(\gamma(t))\ge D(\gamma(t_0))$.
\end{theorem}

\section*{Chapter Summary}
\begin{chapsummary}
  \item $V_R$ is the primary quantitative measure of future
        possibility.
  \item Stronger constraints monotonically reduce $V_R$
        (\cref{thm:reach-mono}).
  \item Constraint accumulation causes exponential decay
        (\cref{thm:constraint-vol}).
  \item Reachability preservation $\equiv$ distinction
        capacity preservation (\cref{thm:reach-dist}).
\end{chapsummary}

% ---- Chapter 14 --------------------------------------------
\chapter{Admissibility}
\label{ch:admissibility}
\begin{epigraph}{It is not enough to reach the destination.
  One must arrive in a condition from which further journeys
  remain possible.}{Attributed to no one; true of everyone}
\end{epigraph}

\begin{definition}{Admissibility Manifold}{adm-manifold}
  $\adm(x,t_0,t)=\overline{\{y\in\reach(x,t_0,t):
  F(y,t,\tau)\ge\alpha V_R(x,t_0,t)\}}$,
  where $F(y,t,\tau)=V_R(y,t,t+\tau)$.
  \index{admissibility!manifold|textbf}
\end{definition}

\begin{theorem}{Admissibility Existence Theorem}{adm-exist}
  For compact $(X,\Phi)$ with continuous $F$ and
  $V_R>0$: $\adm\ne\emptyset$ for any $\alpha\in(0,1)$.
\end{theorem}

\begin{theorem}{Observational--Interventional Separation}%
  {obs-int-sep}
  Observational and interventional admissibility manifolds
  are generically distinct; intervention can expand or
  contract $\adm$.
  \parencite{pearl2000causality,spirtes2000}
\end{theorem}

\begin{theorem}{Admissibility Distortion Theorem}%
  {adm-distortion}
  For non-injective $\pi:X\to Y$:
  $\Delta_\adm(\pi)=\sup_x|V_R(x,t)-V_R(\pi(x),t)|>0$.
  \index{admissibility!distortion|textbf}
\end{theorem}

\begin{theorem}{Projection--Admissibility Gap Theorem}%
  {proj-adm-gap}
  For self-model $\hat\pi:X\to\hat{X}$ with compression
  ratio $r>1$:
  $\Delta_\adm(\hat\pi)\ge
  \frac{\log r}{\log|X|}V_R(x_\mathrm{max},t)$.
\end{theorem}

\begin{proposition}{Repair Preserves Admissibility}%
  {repair-adm}
  Admissible repair satisfies $V_R(\repair(d),t)\ge V_R(d,t)$.
\end{proposition}

\section*{Chapter Summary}
\begin{chapsummary}
  \item $\adm$ is the closed subset of reachable states
        preserving future reachability above $\alpha$.
  \item Every lossy representation induces positive
        admissibility distortion
        (\cref{thm:adm-distortion}).
  \item Admissible repair does not reduce $V_R$
        (\cref{prop:repair-adm}).
\end{chapsummary}

% ---- Chapter 15 --------------------------------------------
\chapter{The Geometry of Admissible Futures}
\label{ch:geometry-futures}
\begin{epigraph}{What we call the beginning is often the end.
  And to make an end is to make a beginning.}%
  {T.S.\ Eliot, \textit{Little Gidding}}
\end{epigraph}

\begin{definition}{Full Admissibility Invariant}{full-adm-inv}
  $\mathbf{I}_\adm(x,t)
  =(\Vol(\adm),S_\adm,\chi(\adm),\kappa_\adm)$:
  volume, diversity, topological complexity, and curvature.
  \index{admissibility!invariant (full)|textbf}
\end{definition}

\begin{definition}{Uncompensated Entropy Growth}{uncomp-entropy}
  Entropy growth is \emph{uncompensated} if not accompanied
  by repair, adaptation, or distinction generation:
  $\Delta S>0$ and $\Delta D\le0$.
\end{definition}

\begin{theorem}{Future Volume Theorem}{future-vol}
  (i)~$\Vol(\adm(t))\ge0$ is well-defined.
  (ii)~$\Vol(\adm(t))\le V_R(x,t)$.
  (iii)~Under uncompensated entropy growth,
  $\Vol(\adm(t))\le\Vol(\adm(t_0))$.
  (iv)~Under regenerative dynamics,
  $\Vol(\adm(t))\ge\Vol(\adm(t_0))$.
  (v)~Under distinction-generating dynamics,
  $\Vol(\adm(t))$ is strictly increasing.
\end{theorem}

\begin{definition}{Trajectory Classification}{traj-classes}
  $\gamma$ is \textbf{generative} if
  $\frac{d}{dt}\Vol(\adm)\ge0$;
  \textbf{extractive} if $<0$;
  \textbf{regenerative} if generative and
  $\frac{d^2}{dt^2}\Vol(\adm)\ge0$;
  \textbf{pathologically continuing} if extractive but
  consuming distinction structures to avoid termination.
  \index{trajectory!generative|textbf}
  \index{trajectory!extractive|textbf}
\end{definition}

\begin{theorem}{Future Bottleneck Theorem}{bottleneck}
  For extractive $\gamma$ with $\Vol(\adm)>0$ on $[t_0,T]$
  and repair continuously exercised at capacity
  $\kappa_\repair>0$:
  (i)~$\gamma$ passes through a bottleneck $t^*$;
  (ii)~$\min_t\Vol(\adm(\gamma(t),t))\ge\kappa_\repair$.
\end{theorem}

\begin{theorem}{Future Preservation Theorem}{future-preserve}
  Let $\gamma_1$ be generative and $\gamma_2$ extractive
  from the same $x$, with equal immediate reward and $r$
  non-decreasing in $\Vol(\adm)$.
  Then $\exists T^*$ such that $\forall T>T^*$:
  $\int_{t_0}^T r(\gamma_1)\,dt\ge
  \int_{t_0}^T r(\gamma_2)\,dt$.
\end{theorem}

\begin{theorem}{Admissibility Curvature Theorem}{adm-curv}
  A perturbation of size $\epsilon$ in a high-curvature
  direction changes $\Vol(\adm)$ by
  $O(\epsilon\cdot\kappa_\adm)$; in a low-curvature
  direction, $O(\epsilon^2)$.
  High-curvature points are decision points where
  trajectory choices have disproportionate consequences.
  \index{admissibility!curvature|textbf}
\end{theorem}

\begin{principle}{Generative Admissibility Principle}{gap-ch15}
  A trajectory is valuable insofar as it preserves the
  capacity for future distinction-production.
  \index{generative admissibility!principle|textbf}
\end{principle}

\begin{theorem}{Future Distinction Optimality Theorem}%
  {gat-ch15}
  Among trajectories achieving equal reward $R$, with
  $r$ continuous and non-decreasing in $\Vol(\adm)$,
  generatively admissible trajectories maximise
  future distinction-producing capacity:
  $\mathbf{I}_\adm(\gamma^*(T),T)
  \ge_\mathrm{lex}\mathbf{I}_\adm(\gamma(T),T)$
  for all $\gamma\in\Gamma_R^-$.
  \index{generative admissibility!theorem|textbf}
\end{theorem}

\begin{corollary}{Intelligence and Admissibility}{gat-intelligence}
  A system is intelligent iff its trajectories are
  generatively admissible w.r.t.\ its ecology.
\end{corollary}
\begin{corollary}{Science and Admissibility}{gat-science}
  A scientific tradition is generatively admissible insofar
  as it expands the volume, entropy, and topological
  complexity of future admissible questions.
\end{corollary}
\begin{corollary}{Governance and Admissibility}{gat-governance}
  Governance is generatively admissible insofar as it
  preserves the reachability volume and diversity of
  policy options for future generations.
\end{corollary}
\begin{corollary}{Memory and Admissibility}{gat-memory}
  Memory is generatively admissible insofar as it
  preserves recoverability needed for future repair.
\end{corollary}
\begin{corollary}{Repair and Admissibility}{gat-repair}
  Admissible repair is generatively admissible;
  repair reducing $\Vol(\adm)$ is worse than no repair.
\end{corollary}
\begin{corollary}{Cosmological Renewal}{gat-cosmo}
  Expyrotic renewal (Ch.\,\ref{ch:expyrotic}) is
  generatively admissible: it restores the universe's
  distinction-producing capacity.
\end{corollary}

\section{From Geometry to Ethics}

The entire progression has been mathematical. No step has
been normative by assumption. Yet the Future Distinction
Optimality Theorem establishes that agents ignoring
admissible volume are \emph{eventually dominated} by those
that do not. Strategic dominance is derived, not postulated.
Moral preference, if it exists, is a further claim.
Strategic preference is already a theorem. The framework
crosses from geometry to ethics not by assumption but
by proof.

\section*{Chapter Summary}
\begin{chapsummary}
  \item The full invariant
        $\mathbf{I}_\adm=(\Vol,S_\adm,\chi,\kappa_\adm)$
        captures volume, diversity, topology, and curvature.
  \item The Future Volume Theorem holds under uncompensated
        entropy growth; evolution and learning are not
        counterexamples.
  \item Growth is not generativity.
  \item The Bottleneck Theorem requires repair to be
        \emph{exercised} (\cref{thm:bottleneck}).
  \item Generative trajectories dominate extractive ones
        over long horizons (\cref{thm:future-preserve}).
  \item The Future Distinction Optimality Theorem proves
        the dominance of generative trajectories
        (\cref{thm:gat-ch15}).
\end{chapsummary}

% ============================================================
%  PARTS VI–X: Stubs with key theorems preserved
%  (Full prose in separate chapter files)
% ============================================================

\part{Physical Realizations}

\chapter{RSVP: The Scalar--Vector Plenum}
\label{ch:rsvp}
\emph{[Full chapter text: see ch16\_rsvp.tex]}

\begin{definition}{RSVP State Variables}{rsvp-fields}
  An RSVP system is described by capacity field $\Phi(x,t)$,
  transport field $\mathbf{v}(x,t)$, and constraint density
  $S(x,t)$.
  \index{RSVP (Relativistic Scalar Vector Plenum)|textbf}
\end{definition}
\begin{theorem}{Three-Field Necessity Theorem}{three-field-necessity}
  Any theory representing stored, transported, and
  constrained distinction capacity requires at least one
  scalar, one vector, and one constraint field.
\end{theorem}
\begin{theorem}{Reachability Capacity Theorem}{rsvp-reachability}
  $V_R(x,t)\propto\Phi(x,t)e^{-S(x,t)}$.
\end{theorem}
\begin{theorem}{RSVP Admissibility Theorem}{rsvp-admissibility}
  Generatively admissible trajectories satisfy
  $\frac{d}{dt}(\Phi e^{-S})\ge0$.
\end{theorem}

\chapter{Gravity as Distinction Dynamics}
\label{ch:gravity}
\emph{[Full chapter text: see ch17\_gravity.tex]}

\begin{theorem}{Gravity as Reachability Optimization}%
  {gravity-reachability}
  Gravitational attraction is motion toward regions of
  maximal future reachability.
\end{theorem}

\chapter{Expyrotic Cosmology}
\label{ch:expyrotic}
\emph{[Full chapter text: see ch18\_expyrotic.tex]}

\begin{theorem}{Expyrotic Admissibility Theorem}{expyrotic-adm}
  Cosmological renewal is generatively admissible.
  \parencite{steinhardt2002,penrose2010}
\end{theorem}

\chapter{Cosmological Completion}
\label{ch:cosmo-completion}
\emph{[Full chapter text: see ch19\_cosmo.tex]}

\begin{theorem}{Cosmological Completion Theorem}{cosmo-completion}
  The physical universe may be interpreted as a regenerative
  distinction ecology evolving under capacity transport,
  constraint accumulation, repair, and admissibility
  preservation.
\end{theorem}

\part{Cognitive Realizations}

\chapter{Semantic Geometry}
\label{ch:semantic-geometry}

\begin{epigraph}{%
  A concept is not a thing.
  A concept is a stable region in a space of distinctions.}%
  {Author}
\end{epigraph}

\begin{objectives}
  \item Define meaning manifolds and conceptual spaces
        as distinction-structured geometries.
  \item Prove the Meaning Distance Theorem.
  \item Prove the Semantic Curvature Theorem.
  \item Prove the Conceptual Attractor Theorem.
  \item Prove the Meaning Repair Theorem.
  \item Prove the Semantic Horizon Propagation Theorem.
  \item Connect semantic geometry to RSVP, admissibility,
        and the ecology of cognitive distinctions.
\end{objectives}

\section{Meaning as Distinction Structure}

The traditional philosophical question asks what meanings
are: abstract objects, mental representations, social
norms, use patterns, or functional roles.
The distinction-theoretic framework dissolves this
dispute by reframing it.
Meanings are not things. Meanings are geometric
structures over distinction spaces.

A concept acquires meaning precisely insofar as it
partitions a domain: separating things that fall under
it from things that do not, cases that are clear from
cases that are ambiguous, extensions from intensions.
A word is semantically rich when the partition it
induces is fine, stable, and recoverable.
A word is semantically impoverished when its partition
is coarse, unstable, or has collapsed toward the
semantic horizon.

This view has antecedents in conceptual space theory
\parencite{gardenfors2000conceptual}, which models
concepts as convex regions in quality dimensions.
The present framework extends this by:
\begin{enumerate}[(i)]
  \item grounding conceptual spaces in the formal
        distinction machinery of Part~I;
  \item adding dynamics — meanings can be damaged and
        repaired;
  \item adding a geometric criterion for evaluating
        meanings — admissibility — that goes beyond
        mere accuracy.
\end{enumerate}

\begin{definition}{Semantic Domain}{semantic-domain}
  A \emph{semantic domain} is a measurable space
  $(W, \mu_W)$ where:
  \begin{itemize}[nosep]
    \item $W$ is a set of \emph{possible referents}:
          objects, events, states, relations, or
          situations that linguistic expressions
          may describe.
    \item $\mu_W$ is a \emph{semantic measure}
          reflecting the relative salience or
          frequency of different regions of $W$.
  \end{itemize}
\end{definition}

\begin{definition}{Meaning Manifold}{meaning-manifold}
  A \emph{meaning manifold} $\mathcal{M}_L$ for a
  language $L$ is a Riemannian manifold
  $(\mathcal{M}, g_S)$ where:
  \begin{itemize}[nosep]
    \item Points $m \in \mathcal{M}$ represent possible
          \emph{meanings}: stable distinction structures
          over $W$.
    \item The metric $g_S$ is the \emph{semantic metric},
          measuring the informational distance between
          meanings.
    \item The scalar curvature $\kappa_S$ of
          $(\mathcal{M}, g_S)$ encodes the local
          density of semantic distinctions.
  \end{itemize}
\end{definition}

\section{The Semantic Metric}

\begin{definition}{Semantic Distance}{sem-distance}
  The \emph{semantic distance} between meanings
  $m_1, m_2 \in \mathcal{M}$ is
  \[
    d_S(m_1, m_2)
    \;=\;
    \sqrt{
      D_\mathrm{KL}(m_1 \| m_2)
      + D_\mathrm{KL}(m_2 \| m_1)
    },
  \]
  the symmetrised KL-divergence between the
  probability distributions over $W$ induced by
  $m_1$ and $m_2$.
\end{definition}

\begin{theorem}{Meaning Distance Theorem}{meaning-distance}
  The semantic distance $d_S$ satisfies:
  \begin{enumerate}[(i)]
    \item $d_S(m_1, m_2) = 0$ iff $m_1$ and $m_2$
          induce identical partitions of $W$.
    \item $d_S$ is symmetric.
    \item $d_S$ satisfies a weak triangle inequality:
          $d_S(m_1, m_3) \le d_S(m_1, m_2)
          + d_S(m_2, m_3) + O(\epsilon^2)$ for
          $\epsilon$-close meanings.
    \item Semantic distance is non-decreasing under
          lossy projection: if $\pi: W \to W'$ is a
          projection, $d_S(\pi(m_1), \pi(m_2))
          \le d_S(m_1, m_2)$.
  \end{enumerate}
\end{theorem}

\begin{proof}
  (i) $D_\mathrm{KL}(m_1 \| m_2) = 0$ iff $m_1 = m_2$
  as distributions. KL-divergence vanishes in both
  directions iff the distributions are identical.
  Identical distributions over $W$ induce identical
  partitions (same cells with same probabilities).

  (ii) The symmetrised KL-divergence is symmetric
  by construction.

  (iii) The symmetrised KL-divergence satisfies a
  second-order triangle inequality in the sense of
  a Fisher--Rao metric on the statistical manifold;
  for nearby points the approximation is exact.

  (iv) Projection coarsens partitions
  (Compression Theorem, \cref{thm:compression}).
  Coarser partitions have lower distinction capacity
  (\cref{thm:refinement-capacity}).
  Lower distinction capacity means fewer
  distinguishable referents, so
  $D_\mathrm{KL}(\pi(m_1) \| \pi(m_2))
  \le D_\mathrm{KL}(m_1 \| m_2)$,
  giving $d_S(\pi(m_1), \pi(m_2)) \le d_S(m_1, m_2)$.
\end{proof}

\begin{remark}{Synonymy and homonymy}{synonymy}
  The Meaning Distance Theorem makes synonymy and
  homonymy precise.
  \emph{Synonyms} are expressions with $d_S \approx 0$:
  they induce nearly identical partitions of $W$.
  \emph{Homonyms} are expressions whose surface form
  is identical but whose meaning manifold positions are
  far apart: $d_S(m_1, m_2) \gg 0$ despite sharing
  a lexical label.
  The theorem also formalises semantic bleaching
  (historical shift toward $d_S = 0$ between an original
  and extended meaning) and semantic drift
  (gradual movement through the meaning manifold).
\end{remark}

\section{Semantic Curvature}

The meaning manifold is not flat. Different regions
support different densities of semantic distinction.
The curvature of $(\mathcal{M}, g_S)$ encodes this
structure.

\begin{definition}{Semantic Curvature}{sem-curvature}
  The \emph{semantic curvature} at $m \in \mathcal{M}$
  is the scalar curvature $\kappa_S(m)$ of the
  Riemannian manifold $(\mathcal{M}, g_S)$.
  Positive curvature indicates a region of high
  semantic density: many distinct meanings are close
  together.
  Negative curvature indicates a region of low semantic
  density: meanings are sparsely distributed.
\end{definition}

\begin{theorem}{Semantic Curvature Theorem}{sem-curvature-thm}
  High semantic curvature at $m$ implies:
  \begin{enumerate}[(i)]
    \item \emph{Distinction sensitivity}: small
          semantic displacements produce large changes
          in the induced partition of $W$.
    \item \emph{Repair difficulty}: semantic damage
          at $m$ is harder to repair because many
          nearby meanings may be conflated.
    \item \emph{Admissibility concentration}: the
          admissibility curvature $\kappa_\adm$
          (\cref{defn:adm-curvature}) is high at $m$,
          marking it as a semantic decision point
          where trajectory choices have large
          consequences.
  \end{enumerate}
\end{theorem}

\begin{proof}
  (i) In a region of high curvature, geodesic
  deviation is large: two nearby meanings diverge
  rapidly under parallel transport. Small displacements
  therefore produce large changes in $d_S$ from other
  fixed meanings, implying large changes in partition
  structure.

  (ii) Repair requires identifying and reconstructing
  the target meaning $m^*$
  (Definition~\ref{defn:repair-op}).
  In high-curvature regions, many candidate meanings
  are close to $m^*$ in the ambient space but induce
  very different partitions. The repair operator must
  distinguish among these candidates, which requires
  higher recoverability and more precise reconstruction
  information.

  (iii) By the Admissibility Curvature Theorem
  (\cref{thm:adm-curv}), admissibility curvature
  amplifies the sensitivity of future possibility to
  trajectory perturbations. In semantic space, high
  curvature marks meanings whose adoption or revision
  propagates large changes through the conceptual
  dependency structure $\mathcal{L}_S$.
\end{proof}

\begin{example}{Technical Vocabulary as High-Curvature Regions}%
  {tech-vocab}
  Technical scientific vocabulary occupies high-curvature
  regions of the meaning manifold.
  Terms like \emph{entropy}, \emph{species},
  \emph{force}, and \emph{information} have precisely
  defined meanings that are close to but distinct from
  their everyday usage.
  The Semantic Curvature Theorem predicts:
  (i) small departures from technical meaning
  produce large inferential errors;
  (ii) repairing confused uses of technical terms
  is difficult because many nearby informal meanings
  are available as attractors;
  (iii) the adoption or revision of technical
  definitions is a high-curvature semantic event
  with large downstream consequences for the
  conceptual ecology of the field.
\end{example}

\section{Conceptual Attractors}

Meanings do not evolve freely. They cluster around
attractors: stable configurations that resist small
perturbations and attract nearby meanings.

\begin{definition}{Conceptual Attractor}{conceptual-attractor}
  A meaning $m^* \in \mathcal{M}$ is a
  \emph{conceptual attractor} if there exists a
  neighbourhood $U \ni m^*$ such that all meanings
  $m \in U$ evolve toward $m^*$ under the semantic
  dynamics:
  \[
    \dot{m} = -\nabla_m E_S(m),
  \]
  where $E_S : \mathcal{M} \to \mathbb{R}$ is the
  \emph{semantic energy} measuring the cost of
  maintaining the distinction structure induced by $m$.
\end{definition}

\begin{theorem}{Conceptual Attractor Theorem}{conceptual-attractor-thm}
  Every stable conceptual attractor $m^*$:
  \begin{enumerate}[(i)]
    \item Minimises semantic energy in its basin of
          attraction:
          $E_S(m^*) \le E_S(m)$ for all $m$ in the basin.
    \item Is a local minimum of the distinction cost
          function $C$ (Definition~\ref{defn:dist-cost}):
          maintaining $m^*$ requires less cognitive,
          social, or computational cost than nearby
          alternatives.
    \item Has positive recoverability:
          $\reco(m^*, t) > 0$, since the attractor
          is maintained by ongoing semantic repair.
    \item Corresponds to a low-entropy region of the
          meaning manifold: $S(m^*, t) < \theta$
          for some threshold $\theta$.
  \end{enumerate}
\end{theorem}

\begin{proof}
  (i) An attractor is a stable fixed point of
  $\dot{m} = -\nabla E_S(m)$.
  Stable fixed points occur at local minima of $E_S$.

  (ii) Semantic energy includes the cost of maintaining
  the partition structure of $m$.
  By the Distinction Cost Theorem
  (\cref{thm:refinement-cost}), maintaining finer
  distinctions costs more than coarser ones.
  A stable attractor balances distinction capacity
  against maintenance cost, settling at a local
  minimum of $C$.

  (iii) A conceptual attractor is actively maintained
  by the linguistic community through use, correction,
  and teaching — all instances of semantic repair.
  This maintenance keeps $\reco(m^*, t) > 0$.

  (iv) Low semantic energy corresponds to low
  multiplicity beneath the distinction structure,
  i.e.\ low entropy $S(m^*)$.
  High-entropy meanings are unstable: they conflate
  many referents, creating pressure toward either
  finer distinction (attractor approach) or complete
  merger (attractor collapse).
\end{proof}

\begin{example}{Basic-Level Categories}{basic-level}
  \bioicon\csicon
  In cognitive linguistics, \emph{basic-level categories}
  — dog, chair, tree — are learned first, named most
  easily, and used most frequently across cultures.
  The Conceptual Attractor Theorem predicts this:
  basic-level categories are conceptual attractors that
  minimise semantic energy by balancing distinction
  capacity (they are not so coarse as to conflate
  importantly different things) against maintenance
  cost (they are not so fine as to require specialised
  knowledge to apply).
  Superordinate categories (animal, furniture) are
  lower-energy but coarser;
  subordinate categories (labrador, Windsor chair)
  are higher-capacity but more costly to maintain.
  The basic level is the attractor basin minimum.
\end{example}

\section{Meaning Repair}

Meanings degrade. Technical terms acquire informal
meanings through metaphorical extension. Evaluative
terms reverse their valence through ironic use.
Scientific concepts lose precision through
popularisation. Political terms are emptied through
euphemism. All of these are forms of semantic damage
requiring repair.

\begin{theorem}{Meaning Repair Theorem}{meaning-repair}
  Semantic repair — restoring a degraded meaning to
  its target distinction structure — is possible iff
  $\reco(m, t) > 0$.
  Among all admissible semantic repairs achieving
  $d_S(m_\mathrm{rep}, m^*) \le \epsilon$,
  the minimal-cost repair:
  \begin{enumerate}[(i)]
    \item Restores the partition of $W$ induced by $m^*$
          as closely as possible.
    \item Does not introduce new semantic distinctions
          not present in $m^*$.
    \item Preserves all non-damaged semantic relations
          in the neighbourhood of $m$.
  \end{enumerate}
\end{theorem}

\begin{proof}
  By the Repair Existence Theorem
  (\cref{thm:repair-exist}), repair is possible iff
  $\reco(m,t) > 0$.
  The minimal repair is given by the Minimal Repair
  Theorem (\cref{thm:minimal-repair}), which in the
  semantic setting corresponds to the minimal
  modification to the partition structure of $W$
  that moves $m$ within $\epsilon$ of $m^*$ in $d_S$.
  (i) follows from the definition of $d_S$: minimising
  $d_S(m_\mathrm{rep}, m^*)$ is equivalent to
  restoring the partition.
  (ii) follows from admissibility: introducing new
  distinctions not in $m^*$ would increase
  the repair cost beyond what is required.
  (iii) follows from the Repair Conservation Law
  (\cref{thm:repair-conservation}): admissible repair
  stays within the connected component of the
  recoverability manifold, preserving non-damaged
  semantic relations.
\end{proof}

\begin{example}{Scientific Popularisation and Repair}%
  {popularisation}
  The word \emph{theory} has a precise scientific
  meaning (a well-confirmed explanatory framework
  with empirical support) and a degraded popular
  meaning (a conjecture or guess).
  This is semantic damage: the popular use has moved
  the meaning far from the scientific attractor in
  the meaning manifold.

  The Meaning Repair Theorem prescribes:
  (i) restore the partition: \emph{theory} should
  separate well-confirmed frameworks from guesses;
  (ii) do not introduce new distinctions: do not
  redefine \emph{theory} as \emph{law} or introduce
  additional conditions not in the original meaning;
  (iii) preserve non-damaged relations: the relationship
  between \emph{theory}, \emph{hypothesis}, and
  \emph{law} in the scientific vocabulary should be
  maintained.
  Science communication that achieves (i)--(iii) is
  performing admissible semantic repair.
\end{example}

\section{Semantic Trajectories and Admissibility}

Meanings evolve through time. A semantic trajectory
$\gamma_S : [t_0, T] \to \mathcal{M}$ describes the
evolution of a meaning through the meaning manifold.
The admissibility framework applies directly.

\begin{definition}{Semantic Reachability Volume}{sem-reach-vol}
  The \emph{semantic reachability volume} at meaning $m$,
  time $t$, is
  \[
    V_S(m, t)
    = \mu_\mathcal{M}\!\bigl(\reach_\mathcal{M}(m, t)\bigr),
  \]
  the measure of meanings accessible from $m$ through
  admissible semantic evolution.
\end{definition}

\begin{theorem}{Semantic Horizon Propagation Theorem}%
  {sem-horizon-prop}
  When a meaning $m$ crosses the semantic horizon —
  when $\reco(m, t) \to 0$ — it propagates horizon
  effects to dependent meanings:
  for all $m'$ with $m \prec m'$
  (Definition~\ref{defn:dist-dep}),
  \[
    \reco(m', t)
    \;\le\;
    \reco(m', t \mid m)
    \cdot \mathbf{1}_{\{\reco(m,t)>0\}}.
  \]
  The loss of a foundational meaning reduces the
  recoverability of all meanings that depend on it.
\end{theorem}

\begin{proof}
  By the Distinction Dependency Theorem
  (\cref{thm:dist-dependency}), if $m \prec m'$,
  then loss of recoverability of $m$ reduces the
  recoverability of $m'$.
  When $\reco(m, t) = 0$, information required for
  the reconstruction of $m'$ via $m$ is unavailable.
  The indicator $\mathbf{1}_{\{\reco(m,t)>0\}}$
  captures the complete dependence: if the supporting
  meaning has zero recoverability, the recoverability
  of the dependent meaning is at most
  $\reco(m', t \mid m)$, but without $m$ this
  conditional recoverability may itself be zero.
\end{proof}

\begin{remark}{Conceptual erosion cascades}{concept-cascade}
  The Semantic Horizon Propagation Theorem describes
  conceptual erosion cascades.
  When a foundational concept loses recoverability,
  the entire network of meanings that depends on it
  becomes harder to maintain.
  Historical examples abound:
  the erosion of the concept of \emph{evidence}
  in public discourse reduces the recoverability of
  \emph{scientific consensus}, \emph{expert opinion},
  and \emph{empirical fact}.
  The loss of the concept of \emph{conflict of interest}
  reduces the recoverability of \emph{impartial
  deliberation} and \emph{professional ethics}.
  These are not merely social problems but structural
  reachability failures in the collective meaning manifold.
\end{remark}

\section{RSVP and the Meaning Manifold}

The scalar-vector-entropy triple of RSVP
(\cref{ch:rsvp}) has a natural semantic interpretation.

\begin{proposition}{RSVP-Semantic Correspondence}{rsvp-semantic}
  Under the identification:
  \begin{align*}
    \Phi(x,t) &\;\leftrightarrow\;
    \text{semantic density at region }x\text{ of }W, \\
    \mathbf{v}(x,t) &\;\leftrightarrow\;
    \text{direction of semantic flow in }\mathcal{M}, \\
    S(x,t) &\;\leftrightarrow\;
    \text{semantic ambiguity (entropy of the partition)},
  \end{align*}
  the RSVP field equations (\cref{ch:rsvp}) describe
  the dynamics of a meaning manifold under linguistic
  use, repair, and degradation.
\end{proposition}

\begin{proof}
  The scalar field $\Phi$ measures local distinction
  capacity in the semantic domain $W$.
  High $\Phi$ corresponds to densely structured
  regions of meaning where many fine distinctions are
  maintained.
  The vector field $\mathbf{v}$ describes the
  flow of semantic influence: how meanings at one
  region of $\mathcal{M}$ influence meanings at
  neighbouring regions through use, metaphor, and
  borrowing.
  The entropy field $S$ measures semantic ambiguity:
  the multiplicity of referents consistent with a
  given expression, corresponding to the hidden
  multiplicity beneath the partition induced by $m$.

  Under these identifications, the RSVP continuity
  equation (\cref{thm:capacity-conservation}) becomes
  a conservation law for semantic density, the
  constraint accumulation equation describes the
  growth of semantic ambiguity, and lamphrodyne
  relaxation (\cref{thm:lamphrodyne}) describes
  the tendency of linguistic communities to reduce
  ambiguity by converging toward semantic attractors.
\end{proof}

\section{Semantic Geometry and Cognitive Realization}

The meaning manifold is the cognitive substrate in
which the distinction ecology of Part~I is realised.
Chapter~\ref{ch:consciousness} showed that
consciousness is recursive repair of self-distinction
structures. Chapter~\ref{ch:preference-fields} showed
that preferences are gradients over admissible future
volume in state space.

Semantic geometry completes the cognitive picture by
providing the space within which these processes occur.
Consciousness maintains the meaning manifold through
recursive repair. Preferences flow along semantic
gradients. Intelligence is repair capacity over
the meaning manifold. The ecology of distinctions
is, at the cognitive scale, an ecology of meanings.

\begin{theorem}{Semantic Intelligence Theorem}{sem-intel}
  Cognitive intelligence $\mathcal{I}(\Sigma)$
  (\cref{defn:intelligence}) is proportional to
  the system's repair capacity over its meaning
  manifold $\mathcal{M}_\Sigma$:
  \[
    \mathcal{I}(\Sigma)
    \;\propto\;
    \kappa(\Sigma, \mathcal{D}_{\mathcal{M}_\Sigma})
    \cdot
    \kappa(\Sigma, \mathcal{D}_{\repair_\mathcal{M}}),
  \]
  where $\mathcal{D}_{\mathcal{M}_\Sigma}$ is the
  distinction space of the meaning manifold and
  $\mathcal{D}_{\repair_\mathcal{M}}$ is the
  distinction space of the semantic repair processes.
\end{theorem}

\begin{proof}
  By the Intelligence--Repair Equivalence Theorem
  (\cref{thm:intel-repair}), intelligence is repair
  capacity over the system's own distinction space,
  including its repair processes.
  For a cognitive system, the primary distinction space
  is the meaning manifold $\mathcal{M}_\Sigma$:
  the distinctions the system maintains and can repair.
  Meta-intelligence — the ability to repair semantic
  repair processes — corresponds to
  $\kappa(\Sigma, \mathcal{D}_{\repair_\mathcal{M}})$.
  The product gives cognitive intelligence.
\end{proof}

\begin{remark}{Language as distinction ecology infrastructure}%
  {language-infra}
  The Semantic Intelligence Theorem explains why language
  is not merely a communication tool but a cognitive
  infrastructure.
  A system operating without language must maintain its
  meaning manifold $\mathcal{M}_\Sigma$ entirely through
  perceptual and motoric distinction structures.
  A system with language can offload meaning maintenance
  onto the shared linguistic meaning manifold of its
  community, dramatically extending its effective
  $\kappa(\Sigma, \mathcal{D}_\mathcal{M})$ at low cost.
  Language is therefore a distributed repair system for
  the collective meaning manifold.
\end{remark}

\section*{Chapter Summary}
\begin{chapsummary}
  \item Meanings are stable distinction structures over
        a semantic domain; the meaning manifold
        $(\mathcal{M}, g_S)$ is their geometric home.
  \item Semantic distance measures informational
        divergence between partition structures;
        projection reduces it
        (\cref{thm:meaning-distance}).
  \item High semantic curvature marks regions where
        small displacements produce large partition
        changes, repair is difficult, and semantic
        decisions have disproportionate consequences
        (\cref{thm:sem-curvature-thm}).
  \item Conceptual attractors are low-energy, low-cost,
        low-entropy, high-recoverability meanings that
        resist small perturbations
        (\cref{thm:conceptual-attractor-thm}).
  \item Semantic repair is possible iff
        $\reco(m,t) > 0$; minimal repair
        restores partition structure without introducing
        new distinctions or destroying existing relations
        (\cref{thm:meaning-repair}).
  \item Loss of foundational meanings propagates
        horizon effects to all dependent meanings
        (\cref{thm:sem-horizon-prop}).
  \item RSVP fields $(\Phi, \mathbf{v}, S)$ describe
        semantic density, flow, and ambiguity
        (\cref{prop:rsvp-semantic}).
  \item Cognitive intelligence is semantic repair
        capacity over the meaning manifold
        (\cref{thm:sem-intel}).
\end{chapsummary}

\subsection*{Exercises}

\begin{exercise}[CS]
  Word embeddings (Word2Vec, GloVe, sentence
  transformers) map words to vectors in $\mathbb{R}^n$.
  Interpret this as a projection
  $\pi_E : \mathcal{M} \to \mathbb{R}^n$.
  By the Meaning Distance Theorem, what does cosine
  similarity in $\mathbb{R}^n$ approximate?
  Under what conditions is $\pi_E$ a lossless
  representation and when does it introduce admissibility
  distortion?
\end{exercise}

\begin{exercise}[Bio]
  Animal communication systems (vervet monkey alarm
  calls, honeybee dance language) involve meaning
  manifolds of very low dimension.
  Using the Conceptual Attractor Theorem, explain
  why predator-category attractors (eagle, snake,
  leopard) are stable despite individual variation
  in vocalisations.
  What would a Semantic Horizon Propagation event
  look like in an animal communication system?
\end{exercise}

\begin{exercise}
  Prove that a perfectly ambiguous language — one in
  which every expression maps to the entire semantic
  domain $W$ — has $V_S(m, t) = 0$ for all $m$ and $t$.
  Interpret this as a reachability collapse.
  What repair operations could restore positive $V_S$?
\end{exercise}

\begin{exercise}[CS]
  A large language model trained on internet text
  inherits the semantic damage present in that corpus
  (imprecise use of technical terms, erosion of
  evaluative distinctions, conflation of concepts).
  Using the Meaning Repair Theorem, design a
  fine-tuning protocol that performs admissible
  semantic repair.
  Which of conditions (i)--(iii) is hardest to satisfy
  in practice, and why?
\end{exercise}

\begin{exercise}
  The Semantic Horizon Propagation Theorem predicts
  that erosion of foundational concepts cascades
  to dependent concepts.
  Identify a foundational concept in one of the
  following domains and trace the cascade of
  recoverability reduction if that concept erodes:
  (a) \emph{evidence} in epistemology;
  (b) \emph{species} in biology;
  (c) \emph{type} in programming language theory;
  (d) \emph{value} in economics.
\end{exercise}

\chapter{Consciousness}
\label{ch:consciousness}
\emph{[Full chapter text: see ch21\_consciousness.tex]}

\begin{theorem}{Consciousness as Recursive Repair}%
  {consciousness-recursive-repair}
  Consciousness is repair directed toward the self-model
  rather than solely the external world.
  \index{consciousness!as recursive repair|textbf}
\end{theorem}

\chapter{Preference Fields}
\label{ch:preference-fields}
\emph{[Full chapter text: see ch22\_preferences.tex]}

\begin{theorem}{Preference--Admissibility Equivalence}%
  {pref-adm-equivalence}
  A preference field is generatively admissible iff
  $\mathcal{L}_{\nabla\Phi_A}\Vol(\adm)\ge0$.
  \index{preference field|textbf}
\end{theorem}

\part{Computational Realizations}

% ============================================================
%  CHAPTER 23: FLOW COMPUTING
% ============================================================
\chapter{Flow Computing}
\label{ch:flow-computing}

\begin{epigraph}{%
  Write programs that do one thing and do it well.
  Write programs that work together.
  Write programs that handle text streams, because that is
  a universal interface.}%
  {Doug McIlroy, \textit{Bell System Technical Journal}, 1978}
\end{epigraph}

\begin{objectives}
  \item Interpret Unix pipelines as distinction-preserving
        transport operators.
  \item Prove the Flow Conservation Theorem.
  \item Prove the Pipeline Determination Theorem.
  \item Prove the Markov Boundary Theorem for processes.
  \item Prove the History Dominance Theorem for file systems.
  \item Apply flow computing to admissibility preservation
        in software architectures.
\end{objectives}

\section{Programs as Distinction Histories}

Chapter~\ref{ch:history} established that states are
compressed projections of histories. Nowhere is this
more visible than in computation.

A running program occupies a state: registers, memory,
program counter. But the interesting thing about a
computation is rarely its instantaneous state. It is its
\emph{trajectory} — the sequence of transformations that
produced the current state from some initial input. Two
programs may reach identical states by entirely different
paths, with entirely different implications for future
behavior.

The tradition of flow-based computing recognises this.
Unix pipelines compose small programs through streams.
Each stage transforms a stream of distinctions — bytes,
lines, records — passing transformed output to the next
stage. The intermediate state of any single process is
largely irrelevant. The relevant object is the flow: the
sequence of distinctions propagating through the pipe.

\begin{definition}{Flow System}{flow-system}
  A \emph{flow system} is a tuple
  $\mathcal{F} = (P, C, \Sigma, \delta)$ where:
  \begin{itemize}[nosep]
    \item $P = \{p_1, \ldots, p_n\}$ is a set of
          \emph{processes};
    \item $C \subseteq P \times P$ is a set of
          \emph{channels} (directed pipes);
    \item $\Sigma$ is an \emph{alphabet} of distinguishable
          tokens;
    \item $\delta_i : \Sigma^* \to \Sigma^*$ is the
          \emph{transformation} of process $p_i$.
  \end{itemize}
  A \emph{flow} through the system is a sequence of tokens
  $s \in \Sigma^\omega$ transformed by the composition of
  processes along channels.
\end{definition}

\begin{definition}{Flow Distinction}{flow-distinction}
  Two flows $f_1, f_2 \in \Sigma^\omega$ are
  \emph{flow-distinguishable} by process $p_i$ if
  $\delta_i(f_1) \ne \delta_i(f_2)$.
  The \emph{flow distinction capacity} of $p_i$ is
  \[
    D_F(p_i) = \log|\{[\delta_i(f)] : f \in \Sigma^*\}|,
  \]
  the log-number of distinguishable output classes.
\end{definition}

\section{The Flow Conservation Theorem}

\begin{theorem}{Flow Conservation Theorem}{flow-conservation}
  For a lossless process $p_i$ (one where $\delta_i$ is
  injective), flow distinction capacity is preserved:
  \[
    D_F(\delta_i(f)) = D_F(f).
  \]
  For a lossy process, $D_F(\delta_i(f)) \le D_F(f)$.
\end{theorem}

\begin{proof}
  If $\delta_i$ is injective, distinct inputs produce
  distinct outputs. The number of distinguishable output
  classes equals the number of input classes.
  Hence $D_F$ is preserved.

  If $\delta_i$ is not injective, there exist
  $f_1 \ne f_2$ with $\delta_i(f_1) = \delta_i(f_2)$.
  Those two flows become indistinguishable at the output.
  The output partition is coarser than the input partition.
  By the Refinement Increases Distinction Capacity Theorem
  (\cref{thm:refinement-capacity}),
  $D_F(\delta_i(f)) \le D_F(f)$.
\end{proof}

\begin{remark}{grep, sort, and uniq as operators}{grep-sort}
  \csicon
  The classic Unix tools realise distinct distinction
  operations. \texttt{grep} is a filter: it selects a
  subset of lines, which is lossy (lines not matching are
  discarded) but \emph{order-preserving} within the
  selected set. \texttt{sort} is a permutation: it
  rearranges distinctions without destroying them, so it
  is lossless in the sense of the Flow Conservation
  Theorem. \texttt{uniq} collapses adjacent identical
  lines into one — a coarsening operation that explicitly
  reduces $D_F$.
\end{remark}

\section{The Pipeline Determination Theorem}

Pipelines compose processes. The question is how
composition affects admissibility.

\begin{definition}{Pipeline}{pipeline}
  A \emph{pipeline} is a sequence of processes
  $\mathbf{p} = (p_1 \to p_2 \to \cdots \to p_k)$
  whose composed transformation is
  \[
    \delta_\mathbf{p} = \delta_k \circ \cdots \circ \delta_1.
  \]
\end{definition}

\begin{theorem}{Pipeline Determination Theorem}{pipeline-det}
  The distinction capacity of a pipeline output is
  determined by its \emph{minimum-capacity stage}:
  \[
    D_F(\delta_\mathbf{p}(f))
    \;\le\;
    \min_{i} D_F(\delta_i).
  \]
\end{theorem}

\begin{proof}
  By the Flow Conservation Theorem, each stage can only
  preserve or reduce $D_F$. The minimum is achieved at
  the bottleneck stage; all subsequent stages cannot
  recover distinctions lost there. Hence the output
  capacity is bounded above by the minimum across stages.
\end{proof}

\begin{remark}{Architectural implication}{pipeline-arch}
  \csicon
  The Pipeline Determination Theorem is a formal version
  of the observation that a system is only as informative
  as its weakest stage. In data pipelines, a single
  \texttt{GROUP BY} that collapses fine-grained events
  into coarse categories permanently reduces the
  distinction capacity available to all downstream
  consumers — even if those consumers would have
  preferred finer granularity. Admissibility-preserving
  pipeline design requires that no stage reduces
  distinction capacity below what downstream stages
  require.
\end{remark}

\section{The Markov Boundary Theorem}

Each process in a flow system sees only its immediate
input. Information from earlier stages is hidden behind
the Markov boundary of the channel.

\begin{definition}{Process Markov Boundary}{process-markov}
  The \emph{Markov boundary} of process $p_i$ in
  pipeline $\mathbf{p}$ is the set of distinctions
  in $\Sigma^*$ that are visible to $p_i$ through its
  input channel, i.e.\ $\mathrm{Im}(\delta_{i-1})$.
\end{definition}

\begin{theorem}{Markov Boundary Theorem}{markov-boundary}
  Process $p_i$ cannot distinguish inputs that are
  identical at its Markov boundary, regardless of their
  upstream history.
  \[
    \delta_{i-1}(f_1) = \delta_{i-1}(f_2)
    \;\implies\;
    \delta_i(\delta_{i-1}(f_1)) = \delta_i(\delta_{i-1}(f_2)).
  \]
\end{theorem}

\begin{proof}
  $p_i$ receives only $\delta_{i-1}(f)$ as input.
  If $\delta_{i-1}(f_1) = \delta_{i-1}(f_2)$, the inputs
  to $p_i$ are identical. Since $\delta_i$ is a function,
  it produces identical output. The upstream distinction
  $f_1 \ne f_2$ is invisible.
\end{proof}

\noindent
The Markov Boundary Theorem formalises why Unix
pipes are composable: each process need only reason about
its immediate input, not the full history of upstream
transformations. But it also explains the cost: upstream
distinctions lost at any stage are irrecoverable by
downstream processes — a computational instance of the
Semantic Horizon Theorem
(\cref{thm:semantic-horizon}).

\section{History Dominance in File Systems}

\bioicon\csicon
File systems that preserve history — version control
systems, append-only logs, event stores — have strictly
higher distinction capacity than those that overwrite
in place.

\begin{theorem}{History Dominance Theorem}{history-dominance}
  Let $\mathcal{F}_H$ be a history-preserving file system
  and $\mathcal{F}_S$ a state-only file system over the
  same sequence of writes.
  Then $D_F(\mathcal{F}_H) \ge D_F(\mathcal{F}_S)$,
  with strict inequality whenever any write modifies
  existing content.
\end{theorem}

\begin{proof}
  $\mathcal{F}_H$ retains all previous states as distinct
  addressable objects.
  $\mathcal{F}_S$ projects the write sequence onto the
  single current state, collapsing historical distinctions.
  By the History Primacy Theorem
  (\cref{thm:history-primacy}),
  $|\hist(\mathcal{F}_H)| > |\hist(\mathcal{F}_S)|$
  whenever writes are non-idempotent.
  Therefore $D_F(\mathcal{F}_H) > D_F(\mathcal{F}_S)$.
\end{proof}

\begin{example}{Git as a Distinction Ecology}{git-ecology}
  \csicon
  A git repository is a distinction ecology in the sense
  of Definition~\ref{defn:dist-ecology}. Each commit is a
  distinction. The dependency graph $\mathcal{L}$ is the
  parent-commit relation. The repair algebra $\mathcal{R}$
  consists of \texttt{git revert}, \texttt{git bisect},
  and \texttt{git cherry-pick}. The semantic horizon
  corresponds to force-pushed history: distinctions that
  once existed but whose recoverability has been
  deliberately driven to zero. Admissibility-preserving
  repository management avoids force-push on shared
  branches for exactly this reason.
\end{example}

\section{Flow Computing and Admissibility}

\begin{theorem}{Flow Admissibility Theorem}{flow-admissibility}
  A flow system is admissible (in the sense of
  Chapter~\ref{ch:admissibility}) if and only if its
  pipeline composition does not reduce the distinction
  capacity available to all reachable downstream consumers.
  Formally:
  \[
    \forall p_i \in P,\quad
    D_F(\delta_i \circ \cdots \circ \delta_1)
    \;\ge\;
    D_{\min}(p_i),
  \]
  where $D_{\min}(p_i)$ is the minimum capacity
  required for $p_i$ to perform its function.
\end{theorem}

\begin{proof}
  Admissibility requires $V_R(\delta(f), t) \ge V_R(f, t)$
  for all transformations $\delta$
  (\cref{prop:repair-adm}).
  In the flow setting, $V_R$ corresponds to the number of
  future computations distinguishable from the current
  state. Reduction of $D_F$ below $D_{\min}(p_i)$
  prevents $p_i$ from distinguishing inputs it needs to
  separate, collapsing future computation branches.
  Hence the condition on $D_F$ is exactly admissibility.
\end{proof}

\section*{Chapter Summary}
\begin{chapsummary}
  \item Flow systems are distinction ecologies: processes
        are distinctions, channels are dependencies,
        transformations are repair or coarsening operators.
  \item Lossless processes preserve $D_F$; lossy processes
        reduce it (\cref{thm:flow-conservation}).
  \item Pipeline capacity is bounded by the minimum-capacity
        stage (\cref{thm:pipeline-det}).
  \item The Markov boundary limits what downstream
        processes can recover
        (\cref{thm:markov-boundary}).
  \item History-preserving systems dominate state-only
        systems in distinction capacity
        (\cref{thm:history-dominance}).
  \item Admissible flow systems do not reduce $D_F$ below
        the minimum required by downstream consumers
        (\cref{thm:flow-admissibility}).
\end{chapsummary}

\subsection*{Exercises}
\begin{exercise}[CS]
  Model a \texttt{sort | uniq -c | sort -rn} pipeline
  as a sequence of distinction operators.
  At which stage does $D_F$ decrease?
  Is the decrease admissible given the stated purpose
  of the pipeline?
\end{exercise}
\begin{exercise}[CS]
  Compare an event-sourced database (append-only log
  of events, state derived on demand) with a
  traditional CRUD database (in-place updates) using
  the History Dominance Theorem.
  Under what conditions does the CRUD system remain
  adequate despite its lower $D_F$?
\end{exercise}
\begin{exercise}
  Prove that a flow system consisting entirely of
  injective processes has constant distinction capacity
  from source to sink. What is the computational
  significance of this result?
\end{exercise}

% ============================================================
%  CHAPTER 24: SPHEREPOP
% ============================================================
\chapter{Spherepop}
\label{ch:spherepop}

\begin{epigraph}{%
  An irreversible event is one that cannot be recalled.
  A computation is an ecology of irreversible events
  from which certain futures can still be reached.}%
  {Author}
\end{epigraph}

\begin{objectives}
  \item Define the four Spherepop primitives:
        Pop, Refuse, Collapse, Bind.
  \item Prove the Refusal Preservation Theorem.
  \item Prove the Collapse Quotient Theorem.
  \item Prove the History Monotonicity Theorem.
  \item Prove the Scope Dynamics Theorem.
  \item Connect Spherepop operational semantics to
        distinction, recoverability, and admissibility.
\end{objectives}

\section{Irreversible Event Calculus}

Chapter~\ref{ch:flow-computing} showed that flow systems
model distinction transport. But flow systems are
\emph{stateless} in the sense that each transformation is
a pure function of its input. Real computation involves
something stronger: commitment. A process may choose to
\emph{pop} an event — to consume it irreversibly,
recording the fact of consumption in its history. Or it
may \emph{refuse} — leaving the event unconsumed and
preserving optionality. Or it may \emph{collapse} a scope
— reducing a structured possibility space to a single
result. Or it may \emph{bind} — creating a persistent
association between a name and a value, extending the
repair capacity of the naming environment.

Spherepop formalises these four operations as the
primitive vocabulary of history-first computation.

\begin{definition}{Spherepop Alphabet}{spherepop-alpha}
  The \emph{Spherepop alphabet} consists of four
  primitive operations:
  \begin{enumerate}[(i)]
    \item $\mathbf{Pop}(e, H)$: consume event $e$,
          appending it to history $H$. Irreversible.
    \item $\mathbf{Refuse}(e, H)$: decline event $e$,
          leaving $H$ unchanged. Optionality-preserving.
    \item $\mathbf{Collapse}(S, H)$: reduce scope $S$
          to its canonical result, appending a
          collapse record to $H$. Irreversible.
    \item $\mathbf{Bind}(n, v, H)$: associate name $n$
          with value $v$ in $H$. Extends naming
          environment.
  \end{enumerate}
\end{definition}

\begin{definition}{Spherepop Configuration}{spherepop-config}
  A \emph{Spherepop configuration} is a pair
  $\langle S, H \rangle$ where:
  \begin{itemize}[nosep]
    \item $S \subseteq \mathcal{E}$ is the current
          \emph{scope}: the set of events available
          for consumption.
    \item $H \in \hist(\mathcal{E})$ is the
          \emph{history}: the ordered sequence of
          events already consumed or refused.
  \end{itemize}
\end{definition}

\begin{definition}{Spherepop Transition Relation}{spherepop-trans}
  The transition relation $\to$ on configurations is
  defined by:
  \begin{align*}
    \langle S \cup \{e\}, H \rangle
    &\xrightarrow{\mathbf{Pop}(e)}
    \langle S, H \cdot e \rangle \\
    \langle S \cup \{e\}, H \rangle
    &\xrightarrow{\mathbf{Refuse}(e)}
    \langle S \setminus \{e\}, H \rangle \\
    \langle S, H \rangle
    &\xrightarrow{\mathbf{Collapse}(S)}
    \langle \varnothing, H \cdot \lfloor S \rfloor \rangle \\
    \langle S, H \rangle
    &\xrightarrow{\mathbf{Bind}(n,v)}
    \langle S, H \cdot (n \mapsto v) \rangle
  \end{align*}
  where $H \cdot e$ denotes appending $e$ to $H$ and
  $\lfloor S \rfloor$ is the canonical reduction of $S$.
\end{definition}

\section{The Refusal Preservation Theorem}

\begin{theorem}{Refusal Preservation Theorem}{refusal-preserve}
  $\mathbf{Refuse}(e, H)$ preserves the reachability
  volume of the configuration:
  \[
    V_R(\langle S \setminus \{e\}, H \rangle, t)
    \;\le\;
    V_R(\langle S \cup \{e\}, H \rangle, t).
  \]
  Moreover, $\mathbf{Refuse}$ is admissible: it does not
  reduce the set of futures reachable by subsequent
  $\mathbf{Pop}$ operations on other events.
\end{theorem}

\begin{proof}
  Before refusal, the scope contains $e$.
  The reachability set includes all futures reachable
  by $\mathbf{Pop}(e)$ and all futures reachable by
  $\mathbf{Refuse}(e)$.
  After refusal, $e$ is removed from scope.
  Futures reachable only via $\mathbf{Pop}(e)$ are
  eliminated.
  Hence $V_R$ after refusal $\le V_R$ before refusal.

  However, refusal does not consume $e$ into $H$,
  so the history remains unchanged.
  All futures reachable from the post-refusal configuration
  were reachable from the pre-refusal configuration
  via the $\mathbf{Refuse}$ branch.
  Therefore refusal is admissible: it selects a branch
  of the future cone without collapsing orthogonal branches
  that do not depend on $e$.
\end{proof}

\begin{remark}{Refuse as optionality management}{refuse-option}
  \csicon
  $\mathbf{Refuse}$ is the Spherepop primitive that most
  directly implements admissibility. A computation that
  refuses events it cannot yet integrate preserves its
  future option space — it does not commit to
  interpretation before sufficient context has accumulated.
  This is the computational analogue of the Admissibility
  Existence Theorem (\cref{thm:adm-exist}): the system
  keeps $\adm$ non-empty by refusing premature commitment.
\end{remark}

\section{The Collapse Quotient Theorem}

\begin{definition}{Scope Quotient}{scope-quotient}
  The \emph{collapse quotient} of scope $S$ under
  $\mathbf{Collapse}$ is the equivalence relation
  $\sim_S$ on $\hist(\mathcal{E})$ defined by:
  \[
    H_1 \sim_S H_2
    \;\iff\;
    \lfloor S \rfloor_{H_1} = \lfloor S \rfloor_{H_2},
  \]
  i.e.\ two histories are equivalent if they produce
  the same canonical reduction of $S$.
\end{definition}

\begin{theorem}{Collapse Quotient Theorem}{collapse-quotient}
  $\mathbf{Collapse}(S)$ projects the history space
  $\hist(\mathcal{E})$ onto the quotient
  $\hist(\mathcal{E})/{\sim_S}$.
  The collapse is:
  \begin{enumerate}[(i)]
    \item Irreversible: $V_R$ after collapse
          $\le V_R$ before collapse.
    \item Admissible iff $\lfloor S \rfloor$ preserves
          all distinctions required by subsequent
          computations.
    \item A projection failure (Ch.~\ref{ch:geometry-failure})
          iff the canonical reduction collapses
          distinctions that generate persistent anomalies
          in downstream scopes.
  \end{enumerate}
\end{theorem}

\begin{proof}
  (i) After collapse, the scope $S$ is replaced by
  $\lfloor S \rfloor$. Any history that would have
  produced a different reduction of $S$ is now
  indistinguishable from the canonical result.
  The quotient $\hist/{\sim_S}$ has fewer equivalence
  classes than $\hist$, so $V_R$ cannot increase.

  (ii) Admissibility requires that downstream
  computations can distinguish all inputs they need to
  separate. If $\lfloor S \rfloor$ preserves those
  distinctions, the collapse is admissible; otherwise not.

  (iii) By the Projection Failure Theorem
  (\cref{thm:proj-fail-thm}), any persistent anomaly
  in downstream scope arises from a distinction collapsed
  by the projection $S \mapsto \lfloor S \rfloor$.
\end{proof}

\begin{example}{Premature Aggregation as Collapse Failure}%
  {premature-agg}
  \csicon
  Aggregating sensor readings into hourly averages before
  storing them is a $\mathbf{Collapse}$ operation. If
  downstream analysis later requires sub-hourly anomaly
  detection, the collapsed data contains a persistent
  anomaly: the anomaly cannot be resolved from the
  available data. This is a projection failure. The repair
  requires re-ingestion at finer granularity — an
  ontological enlargement in the sense of
  Definition~\ref{defn:ont-enlarge}.
\end{example}

\section{The History Monotonicity Theorem}

\begin{theorem}{History Monotonicity Theorem}{history-monotone}
  In any Spherepop execution, the history $H$ grows
  monotonically: no transition removes entries from $H$.
  Formally, for any sequence of transitions
  $\langle S_0, H_0 \rangle \to \cdots \to
  \langle S_n, H_n \rangle$:
  \[
    H_0 \;\sqsubseteq\; H_1 \;\sqsubseteq\;
    \cdots \;\sqsubseteq\; H_n,
  \]
  where $H \sqsubseteq H'$ means $H$ is a prefix of $H'$.
\end{theorem}

\begin{proof}
  By inspection of the transition relation
  (Definition~\ref{defn:spherepop-trans}):
  \begin{itemize}[nosep]
    \item $\mathbf{Pop}(e)$: appends $e$ to $H$.
    \item $\mathbf{Refuse}(e)$: leaves $H$ unchanged.
    \item $\mathbf{Collapse}(S)$: appends $\lfloor S \rfloor$
          to $H$.
    \item $\mathbf{Bind}(n,v)$: appends $(n \mapsto v)$
          to $H$.
  \end{itemize}
  No transition removes or reorders entries. Hence $H$
  is prefix-monotone across all transitions.
\end{proof}

\begin{remark}{Monotone history and the Repair Conservation Law}%
  {mono-repair}
  History Monotonicity is the computational analogue of
  the Repair Conservation Law
  (\cref{thm:repair-conservation}): repair stays within
  the connected component of the recoverability manifold.
  In Spherepop, the history is the recoverability manifold.
  Monotonicity guarantees that past events remain
  recoverable — they are never erased, only extended.
\end{remark}

\section{The Scope Dynamics Theorem}

\begin{definition}{Scope Reachability}{scope-reach}
  The \emph{scope reachability volume} of configuration
  $\langle S, H \rangle$ is
  \[
    V_S(\langle S, H \rangle)
    = |\{H' : \langle S, H \rangle \to^* \langle \varnothing, H' \rangle\}|,
  \]
  the number of terminal histories reachable by complete
  execution of scope $S$.
\end{definition}

\begin{theorem}{Scope Dynamics Theorem}{scope-dynamics}
  For any Spherepop configuration:
  \begin{enumerate}[(i)]
    \item $\mathbf{Pop}$ reduces $V_S$ by eliminating
          futures in which $e$ is refused.
    \item $\mathbf{Refuse}$ reduces $V_S$ by eliminating
          futures in which $e$ is popped.
    \item $\mathbf{Collapse}$ reduces $V_S$ to at most
          $|\{\lfloor S' \rfloor : S' \subseteq S\}|$.
    \item $\mathbf{Bind}$ does not reduce $V_S$; it may
          increase $V_S$ by enabling new branch conditions.
  \end{enumerate}
\end{theorem}

\begin{proof}
  (i) $\mathbf{Pop}(e)$ commits to consuming $e$.
  All futures involving $\mathbf{Refuse}(e)$ from the
  current configuration are eliminated.
  Hence $V_S$ decreases.

  (ii) $\mathbf{Refuse}(e)$ removes $e$ from scope.
  All futures involving $\mathbf{Pop}(e)$ are eliminated.
  Hence $V_S$ decreases.

  (iii) $\mathbf{Collapse}(S)$ reduces all remaining
  scope to a single canonical value. Only distinctions
  preserved by $\lfloor \cdot \rfloor$ survive.
  $V_S$ is bounded by the number of distinct canonical
  reductions.

  (iv) $\mathbf{Bind}$ extends the environment without
  consuming scope events. New bindings may enable
  conditional branches previously unreachable, weakly
  increasing $V_S$.
\end{proof}

\begin{example}{Spherepop Execution Trace}{spherepop-trace}
  \csicon
  Consider a scope $S = \{e_1, e_2, e_3\}$ with
  $V_S = 8$ (three binary choices). A sequence of
  decisions:
  \begin{lstlisting}[language=Spherepop]
  // V_S = 8
  pop e1        // V_S = 4 (e1 consumed; refuse(e1) branch gone)
  refuse e2     // V_S = 2 (e2 removed; pop(e2) branch gone)
  bind x = 42  // V_S = 2 (new binding; no scope reduction)
  collapse S   // V_S = 1 (terminal: one canonical result)
  \end{lstlisting}
  Each $\mathbf{Pop}$ and $\mathbf{Refuse}$ halves
  $V_S$; $\mathbf{Bind}$ leaves it unchanged;
  $\mathbf{Collapse}$ terminates.
\end{example}

\section{Computational Recoverability}

\begin{definition}{Computational Recoverability}{comp-reco}
  The \emph{computational recoverability} of history $H$
  with respect to target configuration
  $\langle S^*, H^* \rangle$ is
  \[
    \reco_C(H)
    = \frac{|\{H' \sqsupseteq H : H' = H^*\}|}
           {|\{H^* \text{ reachable from } H_0\}|}.
  \]
  A computation is \emph{computationally recoverable}
  from $H$ iff $\reco_C(H) > 0$.
\end{definition}

\begin{theorem}{Spherepop Recoverability Theorem}{spherepop-reco}
  A Spherepop execution is recoverable from history $H$
  iff there exists a suffix $\Delta H$ such that
  $H \cdot \Delta H$ produces the target configuration.
  Irrecoverability occurs iff the required suffix is
  incompatible with history monotonicity.
\end{theorem}

\begin{proof}
  By History Monotonicity (\cref{thm:history-monotone}),
  $H$ can only be extended, not modified.
  A target $H^*$ is recoverable from $H$ iff $H \sqsubseteq H^*$.
  If a $\mathbf{Pop}(e)$ in $H$ contradicts the required
  $\mathbf{Refuse}(e)$ in $H^*$, the suffix is
  incompatible: $\reco_C(H) = 0$.
\end{proof}

\section*{Chapter Summary}
\begin{chapsummary}
  \item Spherepop's four primitives — Pop, Refuse,
        Collapse, Bind — implement distinction consumption,
        optionality preservation, scope reduction, and
        environment extension respectively.
  \item $\mathbf{Refuse}$ is the admissibility-preserving
        primitive: it declines premature commitment
        (\cref{thm:refusal-preserve}).
  \item $\mathbf{Collapse}$ is the projection operator:
        it is admissible iff it preserves downstream
        distinctions (\cref{thm:collapse-quotient}).
  \item History grows monotonically; past events are
        never erased (\cref{thm:history-monotone}).
  \item $V_S$ decreases under Pop and Refuse, collapses
        under Collapse, and is preserved or increased
        under Bind (\cref{thm:scope-dynamics}).
  \item Computational recoverability is determined by
        history prefix compatibility
        (\cref{thm:spherepop-reco}).
\end{chapsummary}

\subsection*{Exercises}
\begin{exercise}[CS]
  Model exception handling (\texttt{try}/\texttt{catch})
  in a mainstream language as a Spherepop configuration.
  Which primitive does \texttt{throw} correspond to?
  Is it admissible? What does
  \texttt{finally} correspond to?
\end{exercise}
\begin{exercise}[CS]
  Prove that a purely functional program (no side effects,
  no mutable state) corresponds to a Spherepop execution
  in which every $\mathbf{Pop}$ is immediately followed
  by a $\mathbf{Collapse}$ of the resulting scope.
  What does referential transparency correspond to?
\end{exercise}
\begin{exercise}
  Show that the Spherepop primitives form a monoid under
  sequential composition, identifying the identity element.
  Is the monoid commutative? If not, give a
  counter-example showing that
  $\mathbf{Pop}(e_1) \cdot \mathbf{Pop}(e_2) \ne
   \mathbf{Pop}(e_2) \cdot \mathbf{Pop}(e_1)$ in general.
\end{exercise}

% ============================================================
%  CHAPTER 25: HYDRA
% ============================================================
\chapter{HYDRA}
\label{ch:hydra}

\begin{epigraph}{%
  A mind is a collection of repair strategies for
  the distinctions it cares about.}%
  {Author}
\end{epigraph}

\begin{objectives}
  \item Define the HYDRA architecture as a distinction
        ecology over cognitive modules.
  \item Prove the Hybrid Repair Theorem.
  \item Prove the Projection Management Theorem.
  \item Prove the Multi-Module Admissibility Theorem.
  \item Prove the Constraint-Guided Intelligence Theorem.
  \item Connect HYDRA to the Intelligence--Repair
        Equivalence Theorem of Chapter~\ref{ch:repair-intelligence}.
\end{objectives}

\section{Repair-Based Cognition}

Chapter~\ref{ch:repair-intelligence} defined intelligence
as repair capacity over distinction spaces, including
the distinction space of repair itself. HYDRA —
Hybrid Dynamic Reasoning Architecture — is a
computational realisation of this definition.

HYDRA does not search for optimal solutions to
pre-specified objectives. It maintains a set of
active \emph{distinction projections} (partial models
of the environment) and continuously repairs them
against incoming evidence. When a projection fails —
when it generates persistent anomalies in the sense of
Chapter~\ref{ch:geometry-failure} — HYDRA invokes a
higher-order repair: revision of the projection itself.

\begin{definition}{HYDRA Architecture}{hydra-arch}
  A HYDRA instance is a tuple
  $\mathcal{H} = (M, \Pi, \mathcal{R}, \Gamma, \adm)$
  where:
  \begin{itemize}[nosep]
    \item $M = \{m_1, \ldots, m_k\}$ is a set of
          \emph{modules}, each maintaining a partial
          distinction model of the environment.
    \item $\Pi = \{\pi_{ij} : m_i \to m_j\}$ is a set
          of \emph{projection operators} between modules.
    \item $\mathcal{R}$ is the \emph{repair algebra}
          (Chapter~\ref{ch:repair}) for each module.
    \item $\Gamma$ is a set of \emph{admissibility
          constraints} on module interactions.
    \item $\adm$ is the system-level admissibility
          manifold.
  \end{itemize}
\end{definition}

\begin{definition}{Module State and Anomaly Set}{hydra-module}
  Module $m_i$ maintains:
  \begin{itemize}[nosep]
    \item A distinction space $\mathcal{D}_i$ over its
          domain.
    \item A current model $d_i \in \mathcal{D}_i$.
    \item An anomaly set
          $\mathcal{F}_i = \mathcal{F}(d_i, \mathcal{R}_i)$
          (the failure manifold of Ch.~\ref{ch:geometry-failure}).
  \end{itemize}
\end{definition}

\section{The Hybrid Repair Theorem}

HYDRA performs two levels of repair: local repair within
modules, and cross-module repair when local repair
saturates.

\begin{theorem}{Hybrid Repair Theorem}{hybrid-repair}
  Let $m_i$ be a HYDRA module with saturated local repair
  algebra $\mathcal{R}_i$.
  Then HYDRA invokes cross-module repair:
  it selects module $m_j$ whose distinction space
  $\mathcal{D}_j$ provides an ontological enlargement
  (\cref{defn:ont-enlarge}) resolving anomalies in
  $\mathcal{F}_i$.
  The composed repair
  $\repair_j \circ \pi_{ij} \circ \repair_i$
  is an admissible repair operator on
  $\mathcal{D}_i \cup \mathcal{D}_j$.
\end{theorem}

\begin{proof}
  By the Scientific Revolution Theorem
  (\cref{thm:sci-rev-thm}), saturation of local repair
  plus deep anomalies requires ontological enlargement.
  Module $m_j$ provides $\mathcal{D}_j \supsetneq
  \mathcal{D}_i$ satisfying
  Definition~\ref{defn:ont-enlarge}(i)--(iii).
  The projection $\pi_{ij}$ maps $\mathcal{D}_i$-anomalies
  into $\mathcal{D}_j$-resolvable form.
  Since each individual repair is admissible
  (\cref{defn:admissible-repair}), and the composition
  of admissible repairs is admissible
  (\cref{thm:repair-closure}), the composed operator is
  admissible.
\end{proof}

\begin{remark}{HYDRA and System 1/System 2}{hydra-dual}
  \csicon
  The HYDRA architecture provides a formal account of
  the dual-process distinction in cognitive science.
  System~1 (fast, local, automatic) corresponds to
  repair within a single module $m_i$ using $\mathcal{R}_i$.
  System~2 (slow, deliberate, flexible) corresponds to
  cross-module repair via $\pi_{ij}$ when local repair
  saturates. The key insight is that System~2 is not a
  different \emph{kind} of cognition but a higher-order
  instance of the same repair operation applied to a
  larger distinction space.
\end{remark}

\section{The Projection Management Theorem}

Projections between modules transfer distinctions but may
introduce distortion. HYDRA must manage this distortion
to maintain system-level admissibility.

\begin{theorem}{Projection Management Theorem}{proj-mgmt}
  For any projection $\pi_{ij} : m_i \to m_j$ in HYDRA,
  the admissibility distortion satisfies
  \[
    \Delta_\adm(\pi_{ij})
    \;\ge\;
    \frac{\log r_{ij}}{\log |\mathcal{D}_i|}
    \cdot V_R(d_i, t),
  \]
  where $r_{ij} = |\mathcal{D}_i| / |\mathcal{D}_j|$
  is the compression ratio of the projection.
  HYDRA maintains admissibility by tracking
  $\Delta_\adm(\pi_{ij})$ and triggering module expansion
  when it exceeds a threshold $\theta$.
\end{theorem}

\begin{proof}
  By the Projection--Admissibility Gap Theorem
  (\cref{thm:proj-adm-gap}), any projection with
  compression ratio $r > 1$ introduces distortion at
  least $(\log r / \log |\mathcal{D}|) \cdot V_R$.
  HYDRA monitors this quantity. When
  $\Delta_\adm(\pi_{ij}) > \theta$, the projection is
  losing more admissibility than the system can tolerate.
  Module expansion (adding distinctions to $\mathcal{D}_j$)
  reduces $r_{ij}$, lowering $\Delta_\adm$.
\end{proof}

\section{The Multi-Module Admissibility Theorem}

\begin{theorem}{Multi-Module Admissibility Theorem}{multi-module-adm}
  A HYDRA instance is system-level admissible iff
  the composed projection
  $\Pi = \pi_{k,k-1} \circ \cdots \circ \pi_{21}$
  satisfies
  \[
    \Vol(\adm_\mathcal{H}(t)) \;\ge\; \alpha \cdot V_R(d_0, t)
  \]
  for threshold $\alpha \in (0,1)$, where $d_0$ is the
  root module's model and $\adm_\mathcal{H}$ is the
  system-level admissibility manifold.
\end{theorem}

\begin{proof}
  System-level admissibility is the admissibility of the
  composed transformation (Definition~\ref{defn:adm-manifold}).
  Each projection $\pi_{ij}$ may reduce $V_R$ by at most
  its compression ratio.
  The cumulative reduction across $k$ projections is
  bounded by the product of compression ratios.
  System admissibility requires that this product
  leaves admissible volume above $\alpha V_R(d_0, t)$.
\end{proof}

\section{The Constraint-Guided Intelligence Theorem}

\begin{theorem}{Constraint-Guided Intelligence Theorem}{cgi}
  A HYDRA instance with constraint set $\Gamma$ has
  intelligence
  \[
    \mathcal{I}(\mathcal{H})
    \;=\;
    \prod_{i=1}^k
    \kappa(m_i, \mathcal{D}_{m_i})
    \;\cdot\;
    \kappa(\mathcal{H}, \mathcal{D}_\Pi),
  \]
  where $\mathcal{D}_\Pi$ is the distinction space of
  the inter-module projection algebra.
  Constraints in $\Gamma$ that reduce
  $\kappa(\mathcal{H}, \mathcal{D}_\Pi)$ reduce
  system intelligence; constraints that increase it
  (by specifying repair strategies for projection failures)
  increase system intelligence.
\end{theorem}

\begin{proof}
  By the Intelligence--Repair Equivalence Theorem
  (\cref{thm:intel-repair}), intelligence is the product
  of base repair capacity and meta-repair capacity.
  In HYDRA, base capacity is the product over module
  repair capacities. Meta-capacity is the repair capacity
  over projection failures — the ability of the system
  to revise its own inter-module projection operators.
  Constraints $\Gamma$ that specify admissible projections
  increase meta-capacity by providing repair strategies
  for cross-module distortions; constraints that forbid
  useful projections reduce it.
\end{proof}

\begin{remark}{Why constraint guides rather than limits}{cgi-rmk}
  The Constraint-Guided Intelligence Theorem explains
  why well-designed constraints increase rather than
  decrease cognitive performance. A constraint that
  specifies which projections are admissible provides
  HYDRA with a repair strategy for the case when
  projections fail. Without constraints, HYDRA must
  search the space of possible inter-module projections
  blindly. With constraints, the search is guided.
  The distinction between constraints that increase
  $\kappa$ (good constraints) and those that decrease
  it (bad constraints) formalises the difference between
  useful cognitive biases and pathological ones.
\end{remark}

\section*{Chapter Summary}
\begin{chapsummary}
  \item HYDRA is a distinction ecology over cognitive
        modules with a repair algebra and admissibility
        manifold.
  \item Cross-module repair resolves anomalies that
        saturate local repair, implementing ontological
        enlargement (\cref{thm:hybrid-repair}).
  \item Projection distortion between modules is tracked
        and managed by monitoring
        $\Delta_\adm(\pi_{ij})$
        (\cref{thm:proj-mgmt}).
  \item System admissibility requires the composed
        projection to preserve admissible volume above
        threshold (\cref{thm:multi-module-adm}).
  \item Well-designed constraints increase system
        intelligence by providing meta-repair strategies
        for projection failures (\cref{thm:cgi}).
\end{chapsummary}

\subsection*{Exercises}
\begin{exercise}[CS]
  Model a large language model with retrieval-augmented
  generation (RAG) as a HYDRA instance.
  Identify the modules, projections, repair algebra,
  and admissibility constraints.
  What does the Hybrid Repair Theorem predict happens
  when the base model's context window saturates?
\end{exercise}
\begin{exercise}
  Prove that a single-module system (HYDRA with
  $|M| = 1$) has $\mathcal{I}(\mathcal{H}) = 0$
  by the Constraint-Guided Intelligence Theorem.
  What does this say about the necessity of
  inter-module structure for general intelligence?
\end{exercise}
\begin{exercise}[CS]
  Design an admissibility monitor for a HYDRA instance:
  a process that tracks $\Delta_\adm(\pi_{ij})$ for
  each projection and triggers module expansion
  when the threshold $\theta$ is exceeded.
  What data structure efficiently supports this?
\end{exercise}

% ============================================================
%  CHAPTER 26: TARTAN
% ============================================================
\chapter{TARTAN}
\label{ch:tartan}

\begin{epigraph}{%
  Scale is not merely size.
  Scale is the structure of recoverability across levels.}%
  {Author}
\end{epigraph}

\begin{objectives}
  \item Define TARTAN coherence tiles and recursive
        tiling structure.
  \item Prove the Recursive Tiling Theorem.
  \item Prove the Coherence Tile Theorem.
  \item Prove the Multiscale Recoverability Theorem.
  \item Prove the Trajectory Preservation Theorem.
  \item Connect TARTAN to the admissibility geometry
        of Chapter~\ref{ch:geometry-futures}.
\end{objectives}

\section{Multiscale Distinction Structure}

A persistent challenge in complex systems is that
distinctions operate at multiple scales simultaneously.
A genome contains distinctions at the level of
individual bases, codons, genes, regulatory regions,
and chromosomes. A city contains distinctions at the
level of buildings, blocks, neighbourhoods, districts,
and metropolitan areas. A knowledge base contains
distinctions at the level of tokens, sentences,
paragraphs, documents, and corpora.

Simple partition-based distinction theory applies
cleanly at a single scale. TARTAN — Trajectory-Aware
Recursive Tiling with Annotated Noise — extends the
framework to handle multiscale distinction structures
through recursive decomposition into coherence tiles.

\begin{definition}{Coherence Tile}{coherence-tile}
  A \emph{coherence tile} at scale $\ell$ is a region
  $T \subseteq X$ satisfying:
  \begin{enumerate}[(i)]
    \item \emph{Internal coherence}: the entropy field
          $S(x, t) < \theta_\ell$ for all $x \in T$,
          where $\theta_\ell$ is the scale-$\ell$
          threshold.
    \item \emph{Boundary distinction}: the boundary
          $\partial T$ is a set of high-entropy transitions
          $S(x,t) \ge \theta_\ell$.
    \item \emph{Admissible size}: $\mu(T) \in
          [\mu_{\min}^\ell, \mu_{\max}^\ell]$ for
          scale-appropriate bounds.
  \end{enumerate}
\end{definition}

\begin{definition}{TARTAN Tiling}{tartan-tiling}
  A \emph{TARTAN tiling} at resolution $L$ is a
  hierarchical partition
  \[
    \mathcal{T}^L = \{T^L_1, \ldots, T^L_{n_L}\}
  \]
  of $X$ into coherence tiles, together with a
  refinement sequence
  \[
    \mathcal{T}^0 \preceq \mathcal{T}^1 \preceq
    \cdots \preceq \mathcal{T}^L,
  \]
  where each $\mathcal{T}^\ell$ decomposes the tiles of
  $\mathcal{T}^{\ell-1}$ into finer coherence tiles.
\end{definition}

\begin{definition}{Annotated Noise}{annotated-noise}
  Each tile $T^\ell_i$ carries an \emph{annotation}
  $\eta^\ell_i$ — a noise injection designed to
  explore the local admissibility landscape within $T$
  without destroying tile coherence:
  \[
    \Phi^\ell_i(x,t)
    = \Phi(x,t) + \epsilon_\ell \, \eta^\ell_i(x,t),
  \]
  where $\epsilon_\ell \ll 1$ is the scale-$\ell$
  noise amplitude.
\end{definition}

\section{The Recursive Tiling Theorem}

\begin{theorem}{Recursive Tiling Theorem}{recursive-tiling}
  For any field $\Phi: X \to \mathbb{R}$ with entropy
  field $S = |\nabla \Phi|^2$, a TARTAN tiling of
  resolution $L$ exists satisfying:
  \begin{enumerate}[(i)]
    \item Every tile $T \in \mathcal{T}^L$ is a coherence
          tile (Definition~\ref{defn:coherence-tile}).
    \item The tiling is a refinement of the coarser
          tiling $\mathcal{T}^{L-1}$.
    \item The union $\bigcup_i T^L_i = X$ covers the
          entire domain.
    \item Adjacent tiles share boundary distinctions.
  \end{enumerate}
\end{theorem}

\begin{proof}
  Construct the tiling inductively.

  Base case ($\ell = 0$): Set $\mathcal{T}^0 = \{X\}$,
  the trivial single-tile partition.

  Inductive step: Given $\mathcal{T}^{\ell-1}$, decompose
  each tile $T^{\ell-1}_i$ as follows.
  Compute $S(x,t) = |\nabla\Phi|^2$ on $T^{\ell-1}_i$.
  Identify connected components of the set
  $\{x \in T^{\ell-1}_i : S(x,t) < \theta_\ell\}$.
  Each component satisfying the size constraint
  $\mu(T) \in [\mu_{\min}^\ell, \mu_{\max}^\ell]$
  becomes a tile in $\mathcal{T}^\ell$.

  (i) By construction each component is connected and
  has low internal entropy.
  (ii) Each $\mathcal{T}^\ell$-tile is contained in a
  $\mathcal{T}^{\ell-1}$-tile, so refinement holds.
  (iii) The induction starts with $X$ and only subdivides,
  so coverage is maintained.
  (iv) High-entropy boundaries separate tiles, and all
  high-entropy points lie on boundaries by construction.
\end{proof}

\section{The Coherence Tile Theorem}

\begin{theorem}{Coherence Tile Theorem}{coherence-tile-thm}
  Within a coherence tile $T$ at scale $\ell$:
  \begin{enumerate}[(i)]
    \item RSVP dynamics (\cref{ch:rsvp}) are locally
          stable: $|d\Phi/dt| \le K_\ell$ for a
          scale-dependent constant $K_\ell$.
    \item Distinctions within $T$ have recoverability
          $\reco(d, t) \ge \rho_\ell > 0$.
    \item Annotated noise $\eta^\ell$ explores
          admissible alternatives within $T$ without
          crossing tile boundaries.
  \end{enumerate}
\end{theorem}

\begin{proof}
  (i) Low entropy $S < \theta_\ell$ implies small
  gradients $|\nabla\Phi| < \sqrt{\theta_\ell}$.
  By the RSVP continuity equation
  (\cref{thm:capacity-conservation}), the rate of
  change of $\Phi$ is proportional to divergence of
  $\Phi \mathbf{v}$, which is bounded by
  $\sqrt{\theta_\ell}$. Set
  $K_\ell = C\sqrt{\theta_\ell}$.

  (ii) Low entropy within $T$ means few alternative
  microstates are consistent with the observed
  distinction. By the Entropy--Recoverability Relation
  (\cref{ch:recoverability}), $S_R \le k\log(1/\rho_\ell)$,
  so $\rho_\ell \ge e^{-S_R/k} > 0$.

  (iii) Noise amplitude $\epsilon_\ell \ll 1$ ensures
  perturbations remain below the boundary entropy
  threshold $\theta_\ell$. Hence the perturbed field
  $\Phi + \epsilon_\ell \eta$ does not cross tile
  boundaries.
\end{proof}

\section{The Multiscale Recoverability Theorem}

\begin{theorem}{Multiscale Recoverability Theorem}{multiscale-reco}
  A TARTAN tiling of resolution $L$ satisfies:
  \[
    \reco(d, t \mid \mathcal{T}^L)
    \;\ge\;
    \reco(d, t \mid \mathcal{T}^0),
  \]
  with strict inequality whenever refinement isolates
  a distinction $d$ within a single tile.
\end{theorem}

\begin{proof}
  At resolution $\ell = 0$, all distinctions share the
  single tile $X$. The recoverability of any $d$ is
  determined by the global entropy field.

  At resolution $\ell > 0$, distinctions within a
  coherence tile are separated from high-entropy
  boundary regions. Within the tile, entropy is
  bounded by $\theta_\ell$, which bounds the
  reconstruction ambiguity.

  By the Coherence Tile Theorem,
  $\reco(d, t \mid T^\ell_i) \ge \rho_\ell > 0$ for
  all tiles. Since $\rho_\ell$ is determined by
  $\theta_\ell$ and $\theta_\ell < S_{\max}$ (the
  global maximum), tile-level recoverability exceeds
  global recoverability.

  Refinement from $\mathcal{T}^{\ell-1}$ to
  $\mathcal{T}^\ell$ only increases or maintains
  recoverability by the Refinement Increases Distinction
  Capacity Theorem (\cref{thm:refinement-capacity}).
\end{proof}

\section{The Trajectory Preservation Theorem}

TARTAN is trajectory-aware: each tile tracks not just
the current field configuration but the history of
field trajectories within it.

\begin{definition}{Trajectory Annotation}{traj-annotation}
  The \emph{trajectory annotation} of tile $T^\ell_i$
  at time $t$ is the set of field trajectories
  observed within $T$ up to time $t$:
  \[
    \mathcal{T}\!r^\ell_i(t)
    = \{(\Phi(x,s))_{s \le t} : x \in T^\ell_i\}.
  \]
\end{definition}

\begin{theorem}{Trajectory Preservation Theorem}{traj-preserve}
  Under TARTAN dynamics, the trajectory annotation
  $\mathcal{T}\!r^\ell_i(t)$ is non-decreasing in
  the refinement order: finer tilings preserve
  strictly more trajectory information.
  Formally, for $\ell' > \ell$:
  \[
    H(\mathcal{T}\!r^{\ell'}_i(t))
    \;\ge\;
    H(\mathcal{T}\!r^{\ell}_j(t))
  \]
  for any tile $T^{\ell'}_i \subseteq T^\ell_j$.
\end{theorem}

\begin{proof}
  The trajectory annotation at coarser resolution
  $\ell$ averages trajectories over the larger tile
  $T^\ell_j$. This averaging is a projection:
  it maps the fine-grained trajectory space onto a
  coarser representation.

  By the Compression Theorem (\cref{thm:compression}),
  the entropy of the projected annotation is less than
  or equal to the entropy of the original.
  Since $T^{\ell'}_i \subseteq T^\ell_j$, the finer
  tile's annotation is a restriction (not a compression)
  of the coarser tile's trajectories within that region.

  Restriction to a subtile preserves the trajectory
  distinctions within that subtile while removing
  trajectories from the complement. For trajectories
  entirely within $T^{\ell'}_i$, distinction capacity
  is unchanged. For trajectories crossing the boundary,
  finer tiles separate them, increasing $H$.
  Hence $H(\mathcal{T}\!r^{\ell'}) \ge H(\mathcal{T}\!r^\ell)$.
\end{proof}

\section{TARTAN and Admissibility}

\begin{theorem}{TARTAN Admissibility Theorem}{tartan-adm}
  A TARTAN-structured system is generatively admissible
  (in the sense of Ch.~\ref{ch:geometry-futures}) iff
  for each tile $T \in \mathcal{T}^L$ at the finest
  resolution:
  \[
    \frac{d}{dt}\Vol(\adm(x, t) \cap T)
    \;\ge\; 0
    \quad \forall x \in T.
  \]
  That is, generative admissibility is a tile-local
  condition that propagates to the global system.
\end{theorem}

\begin{proof}
  By the Coherence Tile Theorem, RSVP dynamics are
  locally stable within each tile. The global
  admissibility manifold is the union of tile-local
  admissibility manifolds:
  \[
    \adm(t) = \bigcup_{T \in \mathcal{T}^L}
    (\adm(t) \cap T).
  \]
  Since tiles partition $X$ (Recursive Tiling Theorem),
  \[
    \Vol(\adm(t))
    = \sum_{T} \Vol(\adm(t) \cap T).
  \]
  Therefore $\frac{d}{dt}\Vol(\adm(t)) \ge 0$ iff
  $\frac{d}{dt}\Vol(\adm(t) \cap T) \ge 0$ for each $T$.
\end{proof}

\begin{remark}{TARTAN as distributed admissibility checking}%
  {tartan-distributed}
  The TARTAN Admissibility Theorem converts the global
  admissibility condition into a collection of local
  conditions, one per tile. This is computationally
  significant: rather than computing $\Vol(\adm)$ for
  the entire system (potentially intractable for large
  $X$), TARTAN allows parallel, tile-local admissibility
  monitoring. Each tile is an autonomous admissibility
  checker. The Ecological Fragility Theorem
  (\cref{thm:eco-fragility}) warns against making any
  single tile too large: a highly concentrated tile
  becomes a single point of failure for global
  admissibility.
\end{remark}

\begin{example}{TARTAN in Neural Field Models}{tartan-neural}
  \bioicon
  In neural field theory, activity propagates across
  cortical sheets as travelling waves and localised
  activation patterns. A TARTAN tiling of cortical space
  identifies coherent activation regions (low-entropy
  tiles) separated by high-entropy transition zones.
  Within each tile, trajectory annotations track the
  local history of activation. Cross-tile projections
  implement long-range connectivity. The Multiscale
  Recoverability Theorem predicts that finer-grained
  cortical maps support richer reconstructive capacity
  — consistent with empirical findings that higher
  cortical areas with finer representational granularity
  support more complex inference.
\end{example}

\section*{Chapter Summary}
\begin{chapsummary}
  \item TARTAN tiles partition the domain into coherent
        low-entropy regions separated by high-entropy
        boundaries.
  \item A recursive tiling of any resolution exists for
        any smooth field (\cref{thm:recursive-tiling}).
  \item Within each tile, dynamics are stable,
        recoverability is positive, and noise exploration
        is boundary-respecting
        (\cref{thm:coherence-tile-thm}).
  \item Finer tilings support higher recoverability
        (\cref{thm:multiscale-reco}).
  \item Trajectory annotations grow monotonically with
        refinement (\cref{thm:traj-preserve}).
  \item Global generative admissibility decomposes into
        tile-local conditions, enabling parallel
        admissibility monitoring
        (\cref{thm:tartan-adm}).
\end{chapsummary}

\subsection*{Exercises}
\begin{exercise}[CS]
  Implement the TARTAN tiling algorithm for a 2D field
  $\Phi: [0,1]^2 \to \mathbb{R}$ using NumPy.
  Use $S = |\nabla\Phi|^2$ computed by finite
  differences.
  Visualise tiles at three resolution levels using
  Matplotlib and identify coherence boundaries.
\end{exercise}
\begin{exercise}[Bio]
  Model a developing embryo as a TARTAN system.
  The scalar field $\Phi$ represents morphogen
  concentration. Identify coherence tiles (tissue
  compartments), tile boundaries (organiser regions),
  and trajectory annotations (developmental histories).
  What does the Trajectory Preservation Theorem predict
  about the information content of differentiated
  versus undifferentiated cells?
\end{exercise}
\begin{exercise}
  Prove that a single-resolution TARTAN system
  ($L = 0$) is equivalent to a standard RSVP system
  without tiling. At what resolution $L$ does the
  Multiscale Recoverability Theorem guarantee strictly
  higher recoverability than the RSVP baseline?
\end{exercise}
\begin{exercise}[CS]
  Design a TARTAN-based caching system for a large
  language model's attention computation.
  Define the field $\Phi$ (attention scores),
  entropy $S$ (attention entropy), and coherence tiles
  (semantic segments). Show how the TARTAN Admissibility
  Theorem reduces the cost of checking whether a
  cached attention pattern remains valid.
\end{exercise}

% ============================================================
\part{Social Realizations}
% ============================================================

% ============================================================
%  CHAPTER 27: FISCAL REACHABILITY
% ============================================================
\chapter{Fiscal Reachability}
\label{ch:fiscal}

\begin{epigraph}{%
  A government that cannot finance its commitments
  has already abandoned its future.
  A government that finances only its present
  has sold it.}%
  {Author}
\end{epigraph}

\begin{objectives}
  \item Define fiscal reachability volume and the policy
        state space.
  \item Prove the Fiscal Reachability Theorem.
  \item Prove the Administrative Blind Spot Theorem.
  \item Prove the Boundary Exhaustion Theorem.
  \item Prove the Governance Capacity Theorem.
  \item Apply the framework to debt dynamics, tax base
        erosion, and institutional lock-in.
\end{objectives}

\section{Public Finance as Reachability Geometry}

Public finance is ordinarily described in terms of
flows: revenues, expenditures, deficits, debts.
These are important quantities, but they describe the
\emph{current state} of a fiscal system rather than the
geometry of its future possibilities.

A government that runs a surplus may nonetheless be
trapped: its surplus may be generated by selling
non-renewable resources, eroding its tax base, or
reducing expenditure on the infrastructure that future
revenue requires. Conversely, a government running a
moderate deficit may be expanding its future fiscal
options by investing in productive capacity, human
capital, or institutional resilience.

The distinction-theoretic framework replaces the
flow-accounting perspective with a geometric one. The
central object is not the current balance but the
\emph{fiscal reachability volume}: the measure of policy
states accessible to the government given its current
institutional, financial, and administrative position.

\begin{definition}{Policy State Space}{policy-state}
  The \emph{policy state space} is a measurable space
  $(P, \mu_P)$ where:
  \begin{itemize}[nosep]
    \item $P = P_T \times P_E \times P_I$ is the
          Cartesian product of the taxation space $P_T$,
          the expenditure space $P_E$, and the
          institutional capacity space $P_I$.
    \item $\mu_P$ is a reference measure on $P$
          reflecting the relative importance of
          different policy dimensions.
  \end{itemize}
  A \emph{fiscal position} is a point $p \in P$
  representing the government's current policy
  configuration.
\end{definition}

\begin{definition}{Fiscal Constraint Set}{fiscal-constraint}
  The \emph{fiscal constraint set} $\mathcal{C}_F(t)$
  at time $t$ is the collection of constraints on
  admissible trajectories through $P$, including:
  \begin{enumerate}[(i)]
    \item \emph{Solvency constraints}: the present value
          of future revenues must cover obligations.
    \item \emph{Administrative constraints}: the government
          can only implement policies within its
          institutional capacity $P_I$.
    \item \emph{Political constraints}: some policy
          combinations are excluded by constitutional,
          electoral, or coalitional structure.
    \item \emph{Debt market constraints}: borrowing is
          available only within credit limits.
  \end{enumerate}
\end{definition}

\begin{definition}{Fiscal Reachability Volume}{fiscal-reach}
  The \emph{fiscal reachability volume} at position $p$,
  time $t$, with horizon $\tau$ is
  \[
    V_F(p, t, \tau)
    = \mu_P\!\bigl(
        \reach_P(p, t, t+\tau)
        \cap \adm_P(p, t)
      \bigr),
  \]
  the admissible policy space accessible within horizon
  $\tau$.
\end{definition}

\section{The Fiscal Reachability Theorem}

\begin{theorem}{Fiscal Reachability Theorem}{fiscal-reach-thm}
  Fiscal reachability volume $V_F(p,t,\tau)$ is
  non-increasing under the accumulation of fiscal
  constraints that are not offset by institutional
  capacity expansion:
  \[
    \frac{d|\mathcal{C}_F|}{dt} > 0
    \;\text{ and }\;
    \frac{d\mu_P(P_I)}{dt} \le 0
    \;\implies\;
    \frac{dV_F}{dt} \le 0.
  \]
\end{theorem}

\begin{proof}
  By the Constraint Volume Theorem
  (\cref{thm:constraint-vol}), each additional constraint
  in $\mathcal{C}_F$ reduces $V_F$ by at least
  the measure of the eliminated policy region.
  If institutional capacity $P_I$ does not expand to
  generate new policy options — new tax instruments,
  new administrative pathways, new borrowing facilities
  — then no new reachable states are added to offset
  constraint accumulation.
  Therefore $V_F$ is non-increasing.
\end{proof}

\begin{example}{Debt Trap Dynamics}{debt-trap}
  A government that finances current expenditure through
  non-concessional borrowing accumulates solvency
  constraints: future budgets must dedicate increasing
  fractions of revenue to debt service.
  If this process continues without expanding the tax
  base or productive capacity, the fiscal constraint set
  $\mathcal{C}_F$ grows while $P_I$ contracts
  (administrative capacity is consumed by debt management).
  The Fiscal Reachability Theorem predicts monotone
  decline in $V_F$ — a fiscal trap from which escape
  requires either external debt relief (constraint
  removal) or institutional reform (capacity expansion).
\end{example}

\section{The Administrative Blind Spot Theorem}

Governments operate on compressed representations of
their policy space. A tax administration that does not
model informal economic activity has a blind spot in its
revenue model. A welfare ministry that does not track
long-term outcomes has a blind spot in its effectiveness
model. These blind spots are instances of the general
projection loss established in Chapter~\ref{ch:distinction}.

\begin{definition}{Administrative Projection}{admin-proj}
  An \emph{administrative projection} is a map
  $\pi_A : P \to \hat{P}$ from the full policy space
  to the government's \emph{represented policy space}
  $\hat{P}$ — the portion it can directly observe,
  model, and act upon.
\end{definition}

\begin{theorem}{Administrative Blind Spot Theorem}{admin-blind}
  Every administrative projection $\pi_A : P \to \hat{P}$
  with $|\hat{P}| < |P|$ introduces:
  \begin{enumerate}[(i)]
    \item A \emph{fiscal blind spot}
          $B_A = P \setminus \pi_A^{-1}(\hat{P})$:
          policy regions invisible to administration.
    \item An \emph{admissibility distortion}
          $\Delta_\adm(\pi_A) > 0$: the government's
          estimated fiscal reachability volume exceeds
          its true reachability volume.
    \item A \emph{systematic bias toward the represented}:
          policies well-represented in $\hat{P}$ are
          over-weighted in planning relative to their
          true fiscal impact.
  \end{enumerate}
\end{theorem}

\begin{proof}
  (i) By the Blind Spot Theorem
  (\cref{thm:blind-spot-thm}), any non-injective
  projection produces a non-trivial blind spot.
  Since $|\hat{P}| < |P|$, $\pi_A$ is non-injective.

  (ii) By the Admissibility Distortion Theorem
  (\cref{thm:adm-distortion}), non-injective projections
  produce positive admissibility distortion.
  The government estimates $V_F$ from $\hat{P}$;
  the true $V_F$ depends on $P$;
  the difference is $\Delta_\adm(\pi_A) > 0$.

  (iii) The government's planning operates on $\hat{P}$.
  Policies projecting onto densely represented regions
  of $\hat{P}$ receive more analytical attention than
  policies projecting onto sparse regions,
  regardless of their importance in $P$.
\end{proof}

\begin{remark}{Tax gap as blind spot}{tax-gap}
  The \emph{tax gap} — the difference between taxes
  legally owed and taxes collected — is a direct
  manifestation of the administrative blind spot.
  The informal economy, offshore structures, and
  complex financial instruments represent regions of
  $P$ that fall outside $\hat{P}$.
  The Administrative Blind Spot Theorem predicts that
  governments will systematically underestimate their
  true fiscal reachability because their planning
  tools do not represent these regions.
\end{remark}

\section{The Boundary Exhaustion Theorem}

\begin{definition}{Fiscal Boundary Proximity}{fiscal-boundary}
  The \emph{fiscal boundary proximity} is
  \[
    \beta_F(p, t)
    = d_P\!\bigl(p,\; \partial\reach_P(p,t)\bigr),
  \]
  the distance from the current fiscal position to
  the boundary of the reachable policy space.
\end{definition}

\begin{theorem}{Boundary Exhaustion Theorem}{boundary-exhaust}
  As fiscal boundary proximity $\beta_F \to 0$:
  \begin{enumerate}[(i)]
    \item Small fiscal shocks produce disproportionately
          large reductions in $V_F$.
    \item The sensitivity of policy outcomes to initial
          conditions diverges.
    \item The set of admissible repair operations
          contracts toward the empty set.
  \end{enumerate}
  A government with $\beta_F < \epsilon$ for small
  $\epsilon$ is \emph{fiscally critical}.
\end{theorem}

\begin{proof}
  (i) By the Boundary Proximity and Critical Instability
  Theorem (\cref{thm:boundary-instab}),
  $\|\partial V_F/\partial\epsilon\| \to \infty$
  as $\beta_F \to 0$.
  Fiscal shocks correspond to perturbations $\epsilon$
  to the policy trajectory; near the boundary they
  produce $O(1)$ reductions in $V_F$.

  (ii) Critical instability implies sensitivity to
  initial conditions: small differences in policy
  choices near the boundary produce large differences
  in future reachability.

  (iii) Admissible repair (\cref{prop:repair-adm})
  requires $V_R(\repair(p), t) \ge V_R(p, t)$.
  Near the reachability boundary, almost all operations
  reduce $V_F$; the set of operations satisfying the
  admissibility condition shrinks.
\end{proof}

\begin{example}{Sovereign Debt Crises}{sovereign-crisis}
  The Boundary Exhaustion Theorem describes the
  dynamics of sovereign debt crises. A government
  approaching debt sustainability limits
  ($\beta_F \to 0$) finds that:
  (i) small increases in borrowing costs produce large
  reductions in fiscal space; (ii) small differences
  in initial conditions (e.g.\ whether a crisis hits
  before or after an election) produce large
  differences in outcomes; (iii) the set of admissible
  policy responses available to it contracts —
  austerity, monetisation, and restructuring each
  eliminate further options faster than they create them.
  The crisis is therefore not merely a solvency event
  but a reachability collapse.
\end{example}

\section{The Governance Capacity Theorem}

\begin{theorem}{Governance Capacity Theorem}{gov-capacity}
  A government's long-term fiscal viability satisfies
  \[
    \liminf_{t\to\infty} V_F(p(t), t, \tau) > 0
    \;\iff\;
    G_F(t) + R_F(t) \ge L_F(t)
    \quad \text{eventually},
  \]
  where $G_F$ is the rate of new policy option
  generation, $R_F$ is the rate of constraint removal
  through institutional repair, and $L_F$ is the rate
  of constraint accumulation through fiscal erosion.
\end{theorem}

\begin{proof}
  By the Distinction Balance Law
  (\cref{thm:dist-balance}) applied to the fiscal
  distinction ecology, the rate of change of
  recoverable policy distinctions is
  $\dot{D}_F = G_F + R_F - L_F$.
  By the Universal Regeneration Theorem
  (\cref{thm:universal-regen}), long-term viability
  requires $\liminf V_F > 0$, which by the Ecology
  of Distinctions Conservation Law
  (\cref{thm:ecology-conservation}) holds iff
  $G_F + R_F \ge L_F$ eventually.
\end{proof}

\begin{remark}{Fiscal sustainability as generativity}{fiscal-gen}
  The Governance Capacity Theorem reframes fiscal
  sustainability. The standard definition asks
  whether the debt-to-GDP ratio is bounded.
  The distinction-theoretic definition asks whether
  the government's policy option space is generatively
  admissible. These coincide when debt accumulation is
  the primary source of $L_F$, but diverge when
  institutional capacity erosion dominates — a
  government may have low debt but rapidly contracting
  $V_F$ due to administrative deterioration, regulatory
  capture, or erosion of state capacity.
\end{remark}

\section*{Chapter Summary}
\begin{chapsummary}
  \item Fiscal reachability volume $V_F$ is the primary
        measure of a government's future policy space.
  \item Constraint accumulation without institutional
        expansion monotonically reduces $V_F$
        (\cref{thm:fiscal-reach-thm}).
  \item Administrative projections produce blind spots
        and systematic overestimates of true fiscal
        reachability (\cref{thm:admin-blind}).
  \item Governments near their reachability boundary
        exhibit critical instability: small shocks
        produce large $V_F$ reductions and the
        admissible repair set contracts
        (\cref{thm:boundary-exhaust}).
  \item Long-term fiscal viability requires
        $G_F + R_F \ge L_F$: new option generation
        and institutional repair must keep pace with
        constraint accumulation
        (\cref{thm:gov-capacity}).
\end{chapsummary}

\subsection*{Exercises}
\begin{exercise}
  Model a government that finances current consumption
  by selling state assets as a Continuation Degeneracy
  (\cref{defn:continuation-degeneracy}).
  Identify what reachability volume is being consumed
  and at what rate.
  Under what conditions does the trajectory become
  generatively admissible again?
\end{exercise}
\begin{exercise}
  Define the \emph{fiscal admissibility fraction}
  $\rho_F = V_F / V_R$ — the proportion of reachable
  policy states that are also admissible.
  Prove that $\rho_F$ decreases when extractive fiscal
  policy is pursued, and give three historical examples.
\end{exercise}
\begin{exercise}
  Apply the Administrative Blind Spot Theorem to
  a central bank's monetary policy model.
  Identify the full policy space $P$, the represented
  space $\hat{P}$, and three specific blind spots
  that contributed to the 2008 financial crisis.
\end{exercise}

% ============================================================
%  CHAPTER 28: GOVERNANCE AS FUTURE PRESERVATION
% ============================================================
\chapter{Governance as Future Preservation}
\label{ch:governance}

\begin{epigraph}{%
  The function of government is not to tell us what to do.
  It is to ensure that the space of things we can
  collectively choose to do does not collapse.}%
  {Author}
\end{epigraph}

\begin{objectives}
  \item Define governance as a distinction ecology
        operating on collective policy space.
  \item Prove the Governance Threshold Theorem.
  \item Prove the Option Preservation Theorem.
  \item Prove the Institutional Repair Theorem.
  \item Prove the Democratic Reachability Theorem.
  \item Apply the framework to institutional decay,
        constitutional design, and democratic backsliding.
\end{objectives}

\section{Governance and Future Possibility}

The standard economic account of government concerns
the correction of market failures: public goods,
externalities, information asymmetries. This account
is important but structurally incomplete. It treats
governance as a set of interventions into an otherwise
functioning system of private choices. It does not
address the prior question of what makes that system
of private choices coherent, stable, and reversible.

The distinction-theoretic account begins elsewhere.
Governance is the set of institutions, norms, and
processes that preserve the collective ability to make
and revise choices. Its primary function is not to
make choices but to maintain the space within which
choices remain meaningful. A government that achieves
a particular policy outcome by eliminating the
possibility of future policy revision has not governed
well — it has extracted from the future.

This reframing makes governance a special case of
the Generative Admissibility Principle
(\cref{prc:gap-ch15}): governance is valuable insofar
as it preserves the admissible volume of collective
future possibility.

\begin{definition}{Governance Ecology}{gov-ecology}
  A \emph{governance ecology} is a distinction ecology
  $\mathcal{G} = (\mathcal{D}_G, \mathcal{L}_G,
  \mathcal{R}_G)$ where:
  \begin{itemize}[nosep]
    \item $\mathcal{D}_G$ is the set of institutional
          distinctions: the categories, norms, roles,
          and procedures that structure collective action.
    \item $\mathcal{L}_G$ is the dependency relation
          among institutional distinctions.
    \item $\mathcal{R}_G$ is the repair algebra:
          the set of reform, revision, and correction
          mechanisms available to the polity.
  \end{itemize}
\end{definition}

\section{The Governance Threshold Theorem}

\begin{theorem}{Governance Threshold Theorem}{gov-threshold}
  A governance ecology $\mathcal{G}$ has positive
  long-term collective reachability iff
  the rate of institutional distinction generation
  and repair exceeds the rate of institutional erosion:
  \[
    \liminf_{t\to\infty}
    V_G(p_G, t) > 0
    \;\iff\;
    G_G(t) + R_G(t) \ge L_G(t)
    \quad \text{eventually},
  \]
  where $V_G$ is the collective reachability volume in
  policy space, $G_G$ is institutional innovation,
  $R_G$ is institutional repair, and $L_G$ is
  institutional erosion.
\end{theorem}

\begin{proof}
  The governance ecology is a special case of the
  distinction ecology of Chapter~\ref{ch:distinction-ecology}.
  By the Regenerative Ecology Theorem
  (\cref{thm:regen-ecology-thm}), positive long-term
  viability requires $G_G + R_G \ge L_G$ eventually.
  The Universal Regeneration Theorem
  (\cref{thm:universal-regen}) then gives the
  equivalence with $\liminf V_G > 0$.
\end{proof}

\begin{remark}{What institutional erosion looks like}%
  {inst-erosion}
  Institutional erosion $L_G$ includes: regulatory
  capture (distinctions between public interest and
  private interest collapse); norm erosion (distinctions
  between acceptable and unacceptable conduct become
  ambiguous); administrative degradation (distinctions
  between competent and incompetent public service
  disappear); and constitutional hollowing
  (formal institutional distinctions remain but
  the repair mechanisms that enforce them atrophy).
  In each case the mechanism is the same:
  a distinction that once structured collective action
  loses its recoverability.
\end{remark}

\section{The Option Preservation Theorem}

\begin{theorem}{Option Preservation Theorem}{option-preserve}
  A governance action $a: P \to P$ is \emph{option-preserving}
  iff
  \[
    V_G(a(p), t) \;\ge\; V_G(p, t).
  \]
  Among all governance actions achieving the same
  immediate policy outcome $a(p) = q$, the
  option-preserving action dominates all extractive
  alternatives over sufficiently long horizons.
\end{theorem}

\begin{proof}
  By the Future Preservation Theorem
  (\cref{thm:future-preserve}), among all trajectories
  achieving equal immediate reward, those with
  non-decreasing $V_G$ dominate extractive alternatives
  for any reward function non-decreasing in $V_G$.
  Policy outcomes that citizens value are, in any
  plausible social welfare function, non-decreasing
  in collective reachability volume (since reachability
  volume bounds the set of achievable future outcomes).
  Therefore option-preserving governance actions
  dominate extractive ones over long horizons.
\end{proof}

\begin{example}{Constitutional Entrenchment vs.\ Flexibility}%
  {constitutional}
  Constitutional design involves a tradeoff between
  entrenchment and flexibility.
  Fully entrenched constitutions constrain future
  majorities but also constrain future repair.
  Fully flexible constitutions permit repair but also
  permit erasure of options.

  The Option Preservation Theorem provides a formal
  criterion: a constitutional provision is justified
  by option preservation iff it protects against
  extractive actions that would otherwise reduce $V_G$,
  and does not itself reduce $V_G$ below the level
  achievable without the provision.
  Rights that protect minority distinctions (linguistic,
  religious, cultural) are typically option-preserving:
  they prevent majority erasure of distinction diversity.
  Rights that entrench current majority preferences
  (e.g.\ constitutional property protection for
  existing concentrations) are potentially option-reducing.
\end{example}

\section{The Institutional Repair Theorem}

\begin{theorem}{Institutional Repair Theorem}{inst-repair}
  An institution $\mathcal{I}$ is generatively
  admissible iff its internal repair mechanisms
  $\mathcal{R}_\mathcal{I}$ satisfy:
  \begin{enumerate}[(i)]
    \item \emph{Coverage}: every distinction in
          $\mathcal{D}_\mathcal{I}$ is within the
          repair capacity of $\mathcal{R}_\mathcal{I}$.
    \item \emph{Meta-repair}: $\mathcal{R}_\mathcal{I}$
          can repair itself — the institution has
          mechanisms for reforming its reform mechanisms.
    \item \emph{Admissibility}: every repair in
          $\mathcal{R}_\mathcal{I}$ preserves or
          increases $V_G$.
  \end{enumerate}
\end{theorem}

\begin{proof}
  By the Regeneration Theorem (\cref{thm:regeneration}),
  a system is regenerative iff it preserves repair
  capacity and repair is continuously exercisable.
  (i) Coverage ensures no distinction falls outside
  the repair capacity.
  (ii) Meta-repair is the Regenerative Expansion Theorem
  (\cref{thm:regen-expansion}) applied to $\mathcal{R}$:
  it ensures the repair algebra itself can grow when
  novel anomalies arise.
  (iii) Admissibility of individual repairs is the
  condition of \cref{prop:repair-adm}: repair does not
  reduce reachability volume.
  Together these three conditions are necessary and
  sufficient for generative admissibility of
  $\mathcal{I}$.
\end{proof}

\begin{remark}{Why (ii) is the hardest condition}{meta-repair}
  The meta-repair condition (ii) is routinely violated
  by institutions that encounter persistent anomalies
  in their own functioning but cannot revise their
  reform mechanisms. A legislature that cannot reform
  its own procedures, a court that cannot update its
  interpretive frameworks, a bureaucracy that cannot
  restructure its own incentives — all satisfy (i)
  and (iii) for existing distinctions but fail (ii)
  when those distinctions generate persistent anomalies.
  The Scientific Revolution Theorem
  (\cref{thm:sci-rev-thm}) applies directly:
  institutional saturation plus deep persistent
  anomalies requires ontological enlargement of the
  institutional distinction manifold.
\end{remark}

\section{The Democratic Reachability Theorem}

\begin{theorem}{Democratic Reachability Theorem}{dem-reach}
  Among governance systems with equal decision-making
  efficiency, democracy maximises long-term collective
  reachability volume $V_G$ iff it satisfies:
  \begin{enumerate}[(i)]
    \item \emph{Distinction diversity}: it maintains
          multiple independent institutional distinction
          systems (parties, courts, press, civil society).
    \item \emph{Repair accessibility}: the repair
          algebra $\mathcal{R}_G$ is accessible to
          all members of the polity, not only incumbents.
    \item \emph{Meta-repair preservation}: the
          mechanisms for revising governance structures
          are themselves protected from extraction
          by current majorities.
  \end{enumerate}
\end{theorem}

\begin{proof}
  By the Diversity--Repair Theorem
  (\cref{thm:diversity-repair}), (i) maximises repair
  capacity by providing independent pathways for
  distinction restoration.

  By the Repair Existence Theorem
  (\cref{thm:repair-exist}), repair is possible only
  when recoverability is positive. (ii) ensures that
  recoverability is distributed: if repair is
  accessible only to incumbents, the recoverability of
  distinctions disadvantageous to incumbents may be
  driven to zero.

  By the Institutional Repair Theorem, (iii) is
  necessary for generative admissibility.

  Together (i)--(iii) maximise $V_G$ over time:
  (i) generates the widest repair capacity; (ii) ensures
  that capacity is exercised across the full distinction
  space; (iii) prevents the repair algebra from being
  captured and narrowed by incumbent actors.
  No governance system with fewer of these properties
  can maintain equivalent $V_G$ under the same
  shock distribution.
\end{proof}

\begin{remark}{Democratic backsliding as reachability collapse}%
  {backsliding}
  Democratic backsliding is the process by which
  condition (iii) is progressively violated: incumbents
  use their current majority to modify the meta-repair
  mechanisms (electoral rules, judicial independence,
  press freedom, civil society regulation) in ways that
  reduce the repair accessibility of future minorities.

  The Boundary Exhaustion Theorem
  (\cref{thm:boundary-exhaust}) applies: as backsliding
  proceeds, $\beta_G \to 0$. Small additional steps
  produce disproportionate reductions in $V_G$. The
  admissible repair set contracts toward the empty set.
  At the limit, the polity has passed its governance
  semantic horizon: the repair distinctions that would
  permit reversal have crossed into irrecoverability.
\end{remark}

\section*{Chapter Summary}
\begin{chapsummary}
  \item Governance is a distinction ecology whose
        primary function is preserving collective
        reachability volume $V_G$.
  \item Long-term governance viability requires
        $G_G + R_G \ge L_G$ eventually
        (\cref{thm:gov-threshold}).
  \item Option-preserving governance actions dominate
        extractive alternatives over long horizons
        (\cref{thm:option-preserve}).
  \item Institutions are generatively admissible iff
        they have coverage, meta-repair, and admissible
        repair (\cref{thm:inst-repair}).
  \item Democracy maximises $V_G$ iff it maintains
        distinction diversity, distributed repair
        access, and protected meta-repair mechanisms
        (\cref{thm:dem-reach}).
  \item Democratic backsliding is a reachability
        collapse driven by incumbent extraction of
        meta-repair capacity.
\end{chapsummary}

\subsection*{Exercises}
\begin{exercise}
  Apply the Governance Threshold Theorem to a case
  study of institutional decay in one of: (a) the
  Roman Republic in the first century BCE;
  (b) the Weimar Republic 1930--33;
  (c) a contemporary case of your choosing.
  Identify $G_G$, $R_G$, and $L_G$ for each period
  and the point at which $L_G > G_G + R_G$.
\end{exercise}
\begin{exercise}
  Prove that a single-party state with no independent
  repair mechanisms satisfies
  $\mathcal{I}(\mathcal{G}) = 0$ by the
  Constraint-Guided Intelligence Theorem
  (\cref{thm:cgi}).
  What does this imply about its long-term viability
  under external shocks?
\end{exercise}
\begin{exercise}
  The European Union's subsidiarity principle holds
  that governance should occur at the lowest competent
  level. Interpret this using the Diversity--Repair
  Theorem (\cref{thm:diversity-repair}).
  What does subsidiarity do to collective
  reachability volume compared to centralised governance?
  Under what conditions does it fail?
\end{exercise}

% ============================================================
%  CHAPTER 29: SCIENCE AS A REGENERATIVE SYSTEM
% ============================================================
\chapter{Science as a Regenerative System}
\label{ch:science-regen}

\begin{epigraph}{%
  Science is not a collection of truths.
  It is a collection of methods for producing truths
  that expose themselves to correction.}%
  {Author}
\end{epigraph}

\begin{objectives}
  \item Model science as a distinction ecology with
        anomaly-driven repair dynamics.
  \item Prove the Scientific Repair Theorem.
  \item Prove the Question Generation Theorem.
  \item Prove the Anomaly Preservation Theorem.
  \item Prove the Regenerative Epistemology Theorem.
  \item Apply the framework to scientific revolutions,
        replication crises, and epistemic institutions.
\end{objectives}

\section{Science as Ecology}

Kuhn's account of scientific revolutions
\parencite{kuhn1962structure} identified the structure
of normal science and paradigm shifts but did not
provide a quantitative theory distinguishing healthy
science from pathological science. The distinction-
theoretic framework does.

Science is a distinction ecology in the sense of
Chapter~\ref{ch:distinction-ecology}. The distinctions
are theoretical categories: the partitions that
separate explananda from noise, signal from artifact,
valid inference from fallacy, confirmed hypothesis
from speculation. The dependency relation
$\mathcal{L}_S$ encodes which theoretical distinctions
presuppose which others. The repair algebra
$\mathcal{R}_S$ includes peer review, replication,
meta-analysis, theoretical revision, and paradigm shift.

What distinguishes science from other distinction
ecologies is the \emph{anomaly commitment}: science
explicitly preserves anomalies rather than suppressing
them. A healthy scientific community maintains a
non-empty failure manifold $\mathcal{F}_S$ and
actively invests in its resolution.

\begin{definition}{Scientific Distinction Ecology}{sci-ecology}
  A \emph{scientific distinction ecology} is a tuple
  $\mathcal{S} = (\mathcal{D}_S, \mathcal{L}_S,
  \mathcal{R}_S, \mathcal{F}_S, \mathcal{Q}_S)$ where:
  \begin{itemize}[nosep]
    \item $\mathcal{D}_S$: the set of active theoretical
          distinctions.
    \item $\mathcal{L}_S$: the dependency relation
          (which theories presuppose which others).
    \item $\mathcal{R}_S$: the repair algebra
          (experimental, statistical, and theoretical
          revision mechanisms).
    \item $\mathcal{F}_S = \mathcal{F}(\mathcal{D}_S,
          \mathcal{R}_S)$: the current failure manifold
          (open anomalies).
    \item $\mathcal{Q}_S$: the \emph{question space} —
          the set of distinctions science can currently
          pose but not yet resolve.
  \end{itemize}
\end{definition}

\begin{definition}{Scientific Reachability Volume}{sci-reach}
  The \emph{scientific reachability volume} is
  \[
    V_S(t)
    = \mu_S\!\bigl(\mathcal{Q}_S(t) \cup \adm_S(t)\bigr),
  \]
  the measure of questions science can currently pose
  (the question space) plus its admissible future
  (the currently approachable theoretical territory).
\end{definition}

\section{The Scientific Repair Theorem}

\begin{theorem}{Scientific Repair Theorem}{sci-repair}
  Science performs admissible repair iff each revision
  of $\mathcal{D}_S$:
  \begin{enumerate}[(i)]
    \item Resolves at least one anomaly in $\mathcal{F}_S$.
    \item Preserves all distinctions in $\mathcal{D}_S$
          not implicated in the resolved anomaly.
    \item Does not introduce new anomalies faster than
          it resolves existing ones.
  \end{enumerate}
  A revision satisfying (i)--(iii) is
  \emph{scientifically admissible}.
\end{theorem}

\begin{proof}
  By the Ontology Revision Theorem
  (\cref{thm:ont-revision}), an admissible ontological
  enlargement resolves anomalies while preserving
  existing distinctions (condition~(i) of
  Definition~\ref{defn:ont-enlarge}).
  (i) corresponds to condition~(iii) of the definition:
  persistent anomalies are resolved.
  (ii) corresponds to condition~(i): the embedding is
  isometric.
  (iii) is the additional scientific requirement:
  new distinctions introduced by the revision should
  not generate new failure manifold faster than they
  shrink the existing one, which would indicate a
  net reduction in $V_S$ and hence violation of
  admissibility (\cref{prop:repair-adm}).
\end{proof}

\begin{example}{General Relativity as Admissible Repair}%
  {gr-adm-repair}
  General relativity satisfies (i)--(iii).
  (i) It resolves the persistent anomaly of Mercury's
  perihelion and the Michelson-Morley result.
  (ii) It preserves Newtonian mechanics as a limiting
  case (the equivalence principle ensures the
  Newtonian prediction is recovered for
  $v \ll c$, $GM/r \ll c^2$).
  (iii) The new distinctions it introduces — spacetime
  curvature, gravitational waves, frame dragging —
  generate a vastly expanded question space $\mathcal{Q}_S$
  whose anomalies are actively productive rather than
  stagnating.
\end{example}

\section{The Question Generation Theorem}

The most important property of healthy science is not
that it resolves anomalies but that it generates more
questions than it closes. This is the scientific
analogue of generative admissibility.

\begin{theorem}{Question Generation Theorem}{question-gen}
  A scientific tradition is generatively admissible iff
  \[
    \frac{d}{dt}|\mathcal{Q}_S(t)| \;\ge\; 0:
  \]
  the question space does not contract.
  Furthermore, a \emph{great theory} is one for which
  \[
    \frac{d}{dt}|\mathcal{Q}_S(t)| \gg 0:
  \]
  it expands the question space strictly and rapidly.
\end{theorem}

\begin{proof}
  By the Generative Admissibility Principle
  (\cref{prc:gap-ch15}), a system is generatively
  admissible iff it preserves the capacity for future
  distinction-production.
  In science, future distinction-production corresponds
  to future question resolution.
  The capacity for future question resolution is bounded
  by the current question space: questions not yet
  posed cannot yet be resolved.
  Therefore preserving distinction-production capacity
  requires preserving or expanding $\mathcal{Q}_S$.
  $|\mathcal{Q}_S(t)|$ non-decreasing is the
  scientific instantiation of
  $\frac{d}{dt}\Vol(\adm(t)) \ge 0$.
\end{proof}

\begin{remark}{The failure of paradigm-ending theories}{paradigm-end}
  The Question Generation Theorem explains why some
  apparently successful theories are scientifically
  problematic. A theory that claims to explain
  everything within its domain while generating no
  new questions is extractive: it resolves existing
  anomalies by eliminating the distinctions that
  would generate new ones.

  A theory of everything that provides exactly one
  prediction has $|\mathcal{Q}_S| = 1$ after
  its adoption: it has collapsed the question space
  to a single empirical test. If the test fails,
  the entire question space collapses to zero.
  If it succeeds, science loses the productive
  diversity of competing frameworks.
  The Question Generation Theorem identifies this
  as a form of scientific pathological continuation:
  the theory survives while the science contracts.
\end{remark}

\section{The Anomaly Preservation Theorem}

\begin{theorem}{Anomaly Preservation Theorem}{anomaly-preserve}
  A healthy scientific community preserves anomalies:
  it maintains $\mathcal{F}_S \ne \emptyset$
  and ensures anomalies are available for resolution
  rather than suppressed or marginalised.
  Formally, science is generatively admissible only if
  \[
    \reco(\mathcal{F}_S, t) > 0
    \quad\text{for all } t.
  \]
\end{theorem}

\begin{proof}
  By the Repair Existence Theorem
  (\cref{thm:repair-exist}), repair is possible iff
  recoverability is positive. Scientific repair
  (theory revision) operates on anomalies: it requires
  that the anomalous observations are recoverable from
  the scientific record.

  If $\reco(\mathcal{F}_S, t) = 0$, anomalies have
  passed the scientific semantic horizon: they are
  no longer accessible to the repair algebra.
  This occurs when anomalous results are unpublished,
  retracted without cause, or buried in inaccessible
  archives.
  In this case $\mathcal{R}_S$ cannot act on
  $\mathcal{F}_S$, and the failure manifold cannot shrink.
  The science stagnates: it cannot repair distinctions
  that have lost recoverability.

  Conversely, if $\reco(\mathcal{F}_S, t) > 0$,
  anomalies remain accessible and the repair algebra
  can act on them. This is the necessary condition for
  scientific progress.
\end{proof}

\begin{remark}{Replication crisis as recoverability failure}%
  {replication-crisis}
  The replication crisis in psychology, medicine, and
  social science is formally a recoverability failure.
  Non-publication of null results drives $\reco$
  toward zero for the null result distinctions. The
  published literature over-represents positive
  findings, creating a biased failure manifold that
  the repair algebra cannot act on because the
  anomalous (null) results are not recoverable from
  the available record.

  The Anomaly Preservation Theorem predicts that
  pre-registration, mandatory reporting of null results,
  and open data requirements all increase
  $\reco(\mathcal{F}_S)$ and therefore restore the
  condition for scientific repair.
\end{remark}

\section{The Regenerative Epistemology Theorem}

\begin{theorem}{Regenerative Epistemology Theorem}{regen-epistem}
  A scientific tradition is regenerative iff:
  \begin{enumerate}[(i)]
    \item It preserves the recoverability of anomalies:
          $\reco(\mathcal{F}_S, t) > 0$.
    \item It trains new practitioners who can exercise
          the full repair algebra $\mathcal{R}_S$.
    \item It preserves methodological diversity:
          $N(\mathcal{R}_S) \ge N_{\min}$ independent
          repair pathways.
    \item It maintains a non-trivial failure manifold:
          $\mathcal{F}_S \ne \emptyset$.
  \end{enumerate}
\end{theorem}

\begin{proof}
  By the Principle of Regeneration
  (\cref{prc:regen-prc}), a system is regenerative
  iff it preserves repair capacity.

  (i) preserves recoverability, which by the Repair
  Existence Theorem is necessary for any repair.
  (ii) preserves the ability to \emph{exercise} the
  repair algebra — a repair algebra that cannot be
  operated is operationally zero capacity.
  (iii) by the Diversity--Repair Theorem
  (\cref{thm:diversity-repair}), multiple independent
  repair pathways maximise repair capacity; losing
  methodological diversity reduces $\kappa(\mathcal{S})$.
  (iv) an empty failure manifold ($\mathcal{F}_S = \emptyset$)
  means no anomalies are recognised. This cannot arise
  in a healthy science (every partition produces blind
  spots, every model has anomalies) and instead
  indicates that anomalies are being suppressed rather
  than resolved — a failure of (i).

  Together (i)--(iv) are necessary and sufficient for
  the Regeneration Theorem (\cref{thm:regeneration}).
\end{proof}

\begin{example}{Healthy vs.\ Pathological Science}%
  {healthy-sci}
  A healthy scientific tradition satisfies all four
  conditions of the Regenerative Epistemology Theorem.
  Physics in the late 19th century: anomalies (Mercury,
  Michelson-Morley, blackbody radiation) were
  publicly recognised and preserved; new practitioners
  were trained in classical mechanics and in the new
  experimental methods; methodological diversity was
  high (mathematical physics, experimental physics,
  theoretical chemistry); and the failure manifold
  was acknowledged as non-empty.

  A pathological scientific tradition violates at
  least one condition. Lysenkoism in Soviet biology
  violated (i) by suppressing anomalies; (ii) by
  excluding geneticists from training programmes;
  (iii) by mandating a single methodological framework;
  and (iv) by declaring the failure manifold empty
  through ideological fiat.
  The Regenerative Epistemology Theorem predicts
  total collapse of repair capacity — which occurred.
\end{example}

\section{Science and the Full Admissibility Invariant}

The four components of the full admissibility invariant
$\mathbf{I}_\adm = (\Vol, S_\adm, \chi, \kappa_\adm)$
(\cref{defn:full-adm-inv}) have direct scientific
interpretations.

\begin{proposition}{Scientific Admissibility Invariant}%
  {sci-adm-inv}
  For a scientific distinction ecology $\mathcal{S}$:
  \begin{enumerate}[(i)]
    \item $\Vol(\adm_S)$ measures the breadth of
          approachable scientific territory.
    \item $S_\adm(\mathcal{S})$ measures the diversity
          of scientific futures: how evenly the
          approachable territory is distributed across
          distinct research directions.
    \item $\chi(\adm_S)$ measures the topological
          complexity of the admissible scientific future:
          how many disconnected research programmes
          are simultaneously viable.
    \item $\kappa_\adm(\mathcal{S})$ identifies
          high-curvature decision points in the
          scientific landscape — paradigm shifts,
          methodological revolutions, and theoretical
          unifications where small choices have
          large long-term consequences.
  \end{enumerate}
\end{proposition}

\begin{proof}
  Each component of $\mathbf{I}_\adm$ is applied to
  the scientific policy space $P_S$ (the space of
  possible research programmes, theoretical frameworks,
  and methodologies).
  (i) is direct: admissible volume equals accessible
  research territory.
  (ii)--(iv) follow from the definitions of
  admissibility entropy, Euler characteristic, and
  curvature (\cref{ch:geometry-futures}) applied to $P_S$.
\end{proof}

\section*{Chapter Summary}
\begin{chapsummary}
  \item Science is a distinction ecology whose repair
        algebra is exercised through peer review,
        replication, and theoretical revision.
  \item Scientific repair is admissible iff revisions
        resolve anomalies, preserve existing distinctions,
        and do not introduce anomalies faster than they
        resolve them (\cref{thm:sci-repair}).
  \item Healthy science expands the question space:
        $\frac{d}{dt}|\mathcal{Q}_S| \ge 0$
        (\cref{thm:question-gen}).
  \item Anomaly recoverability is necessary for
        scientific repair; suppression of null results
        is a recoverability failure
        (\cref{thm:anomaly-preserve}).
  \item Science is regenerative iff it preserves
        anomaly recoverability, trains practitioners,
        maintains methodological diversity, and
        keeps a non-trivial failure manifold
        (\cref{thm:regen-epistem}).
  \item The four components of $\mathbf{I}_\adm$
        measure breadth, diversity, topological
        complexity, and decision sensitivity of the
        scientific future (\cref{prop:sci-adm-inv}).
\end{chapsummary}

\subsection*{Exercises}
\begin{exercise}
  Apply the Question Generation Theorem to the
  development of quantum mechanics 1900--1930.
  Plot (qualitatively) $|\mathcal{Q}_S(t)|$ through
  the period. At which points does the question space
  expand most rapidly? What caused those expansions?
\end{exercise}
\begin{exercise}
  The pre-registration movement in psychology requires
  researchers to register hypotheses and analysis plans
  before collecting data.
  Interpret pre-registration using the Anomaly
  Preservation Theorem.
  Show formally that pre-registration increases
  $\reco(\mathcal{F}_S)$ relative to a publish-only-
  positive-results norm.
\end{exercise}
\begin{exercise}
  Two scientific programmes are compared:
  Programme~A generates 10 confirmed predictions and
  0 new questions.
  Programme~B generates 5 confirmed predictions and
  50 new questions.
  Which is more generatively admissible?
  Under what reward function would Programme~A be
  rationally preferred despite being extractive?
\end{exercise}
\begin{exercise}
  Apply the Regenerative Epistemology Theorem to one
  of: (a) the decline of phlogiston chemistry;
  (b) the rise of molecular biology;
  (c) the current state of string theory.
  Which of conditions (i)--(iv) are satisfied, and
  which are under strain?
\end{exercise}


% ============================================================
\part{Synthesis}
% ============================================================

\chapter{The Ecology of Distinctions}
\label{ch:ecology-synthesis}

\begin{epigraph}{%
  Nothing in biology makes sense
  except in the light of evolution.
  Nothing in this book makes sense
  except in the light of distinction.}%
  {Author, after Dobzhansky}
\end{epigraph}

\begin{objectives}
  \item Prove the Preservation Hierarchy Theorem.
  \item Prove the Preservation Equivalence Theorem.
  \item Prove the Category of Regenerative Systems Theorem.
  \item Prove the Unified Invariant Theorem.
  \item Prove the Ecology Conservation Law.
  \item Show that every realization in Parts~VI--IX is a
        functor from the abstract theorem category into
        a domain-specific distinction ecology.
  \item Demonstrate that the entire book is a single
        argument whose conclusion is the Generative
        Admissibility Theorem of Chapter~\ref{ch:preservation}.
\end{objectives}

% ------------------------------------------------------------
\section{The Argument So Far}
\label{sec:ch30-arg}
% ------------------------------------------------------------

Thirty chapters have been devoted to a single question
asked at progressively deeper levels of abstraction.

Chapter~\ref{ch:distinction} asked: what must exist before
anything can be observed? The answer was distinction — the
primitive act of partitioning a domain that simultaneously
produces objects, information, cost, and blind spots.

Chapter~\ref{ch:information} asked: what is information?
The answer was: the quantitative expression of distinctions,
not a primitive substance.

Chapter~\ref{ch:entropy} asked: what is entropy? The answer
was: the multiplicity hidden beneath a distinction structure,
not disorder.

Chapter~\ref{ch:history} asked: what are states? The answer
was: compressed projections of histories.

Chapter~\ref{ch:recoverability} asked: when is a lost
distinction gone forever? The answer was: only when
recoverability falls to zero — dispersal and destruction
are not equivalent.

Chapter~\ref{ch:memory} asked: what is memory? The answer
was: not storage, but preservation of recoverability.

Chapters~\ref{ch:repair}--\ref{ch:repair-intelligence}
asked: how do distinctions persist despite entropy? The
answer was: through repair — the third primitive,
irreducible to distinction and entropy, whose existence
depends on positive recoverability and whose algebra
is a monoid.

Chapters~\ref{ch:beyond-continuation}--\ref{ch:distinction-ecology}
asked: is repair enough? The answer was: no. Repair
mechanisms themselves degrade. Regeneration — preservation
of repair capacity — is necessary. And even regeneration
is insufficient unless the futures it generates remain
open. This led to the concept of distinction ecology:
a network of interacting distinctions whose structure
determines the possibility of future distinction-production.

Chapters~\ref{ch:reachability}--\ref{ch:geometry-futures}
answered the geometric question: what does future
possibility look like? Reachability volume $V_R$
measures how much of the future is accessible.
Admissibility volume $\Vol(\adm)$ measures how much
of the accessible future preserves further accessibility.
The full admissibility invariant
$\mathbf{I}_\adm = (\Vol, S_\adm, \chi, \kappa_\adm)$
measures volume, diversity, topological complexity, and
decision sensitivity simultaneously. The Future
Distinction Optimality Theorem
(\cref{thm:gat-ch15}) proved that generatively admissible
trajectories dominate extractive alternatives over any
sufficiently long horizon — a strategic dominance
derived, not postulated.

Parts~VI--IX demonstrated that the same abstract structure
appears in physics, cognition, computation, and society.
RSVP fields realise the distinction-repair-admissibility
programme in physical field theory.
Gravity becomes reachability optimisation.
Expyrosis becomes the largest-scale repair operation.
Semantic geometry shows that meanings are distinction
structures on a Riemannian manifold.
Consciousness is recursive repair of self-distinction.
Preferences are admissibility gradients.
Flow computing, Spherepop, HYDRA, and TARTAN realise
the framework computationally.
Fiscal reachability, governance, and science are its
social instantiations.

The present chapter does something different from all
of the above. It does not introduce new domains.
It proves that these domains are not merely analogues
of one another but \emph{projections} of a single
underlying geometric object.

% ------------------------------------------------------------
\section{The Preservation Hierarchy}
\label{sec:ch30-hierarchy}
% ------------------------------------------------------------

The argument of this book proceeds by discovering that
each concept introduced is, upon reflection, insufficient
without a higher-order concept. Distinctions require
repair. Repair requires regeneration. Regeneration
requires admissibility. Admissibility requires generativity.

This progression is not arbitrary. It reflects a genuine
hierarchy of preservation — a sequence of increasingly
strong conditions, each strictly stronger than the last.

\begin{definition}{The Preservation Classes}{dist-hierarchy}
  Define the following classes of systems:
  \begin{align*}
    \mathfrak{D} &= \{\text{systems possessing distinctions}\}, \\
    \mathfrak{M} &= \{\text{systems preserving distinctions historically}\}, \\
    \mathfrak{R} &= \{\text{systems capable of repair}\}, \\
    \mathfrak{G} &= \{\text{systems preserving repair capacity (regenerative)}\}, \\
    \mathfrak{A} &= \{\text{systems preserving future reachability (admissible)}\}, \\
    \mathfrak{P} &= \{\text{systems expanding admissible possibility (generative)}\}.
  \end{align*}
\end{definition}

\begin{theorem}{Preservation Hierarchy Theorem}{preservation-hierarchy}
  \[
    \mathfrak{P}
    \;\subsetneq\;
    \mathfrak{A}
    \;\subsetneq\;
    \mathfrak{G}
    \;\subsetneq\;
    \mathfrak{R}
    \;\subsetneq\;
    \mathfrak{M}
    \;\subsetneq\;
    \mathfrak{D}.
  \]
  Each inclusion is strict.
\end{theorem}

\begin{proof}
  \textit{Inclusions.}
  Every possibility-expanding system preserves admissible
  future volume, hence $\mathfrak{P} \subseteq \mathfrak{A}$.
  Every admissible system preserves future repair pathways,
  hence $\mathfrak{A} \subseteq \mathfrak{G}$.
  Every regenerative system preserves repair capacity,
  hence $\mathfrak{G} \subseteq \mathfrak{R}$.
  Every repair system must preserve historical information
  (Repair Conservation Law, \cref{thm:repair-conservation}),
  hence $\mathfrak{R} \subseteq \mathfrak{M}$.
  Every memory system necessarily contains distinctions,
  hence $\mathfrak{M} \subseteq \mathfrak{D}$.

  \textit{Strictness.}
  $\mathfrak{M} \subsetneq \mathfrak{D}$:
  a crystal persists and contains distinctions (temperature,
  crystal symmetry) but has no memory in the functional sense.
  $\mathfrak{R} \subsetneq \mathfrak{M}$:
  a read-only archive preserves historical information but
  cannot repair damaged distinctions.
  $\mathfrak{G} \subsetneq \mathfrak{R}$:
  a DNA repair enzyme executes repair but does not
  preserve its own regeneration capacity
  (it is consumed in the process).
  $\mathfrak{A} \subsetneq \mathfrak{G}$:
  a highly adaptive but short-lived organism regenerates
  locally without maintaining admissible volume
  at the population level.
  $\mathfrak{P} \subsetneq \mathfrak{A}$:
  a system may maintain admissible volume ($\Vol(\adm)$
  constant) without expanding it; generativity requires
  strict increase.
\end{proof}

\noindent
The hierarchy is not a taxonomic curiosity. It is the
formal structure of what it means to be more or less
alive, more or less intelligent, more or less capable
of surviving the future. Every living system,
every institution, every theory, every computation
sits somewhere in this hierarchy. The deepest question
one can ask about any of them is: at what level does it
reside, and is it moving up or down?

% ------------------------------------------------------------
\section{The Unified Invariant}
\label{sec:ch30-unified}
% ------------------------------------------------------------

The preceding parts of the book introduced several
seemingly independent quantities:
entropy $S$, recoverability $\reco$, memory $\mathcal{M}$,
repair capacity $\kappa_\repair$, reachability volume $V_R$,
and admissibility volume $\Vol(\adm)$.

We now show these are not independent. They are all
projections of a single underlying object: the
possibility functional $\mathcal{P}(x,t) = \Vol(\adm(x,t))$.

\begin{theorem}{Preservation Equivalence Theorem}%
  {preservation-equivalence}
  Under admissible dynamics, the following quantities
  are monotone together:
  \[
    \mathcal{P}(t),\quad
    V_R(t),\quad
    \kappa_\repair(t),\quad
    \mathcal{M}(t),\quad
    D_\mathcal{E}(t).
  \]
  Each is non-decreasing iff all are non-decreasing.
\end{theorem}

\begin{proof}
  Admissibility gives $\frac{d}{dt}\mathcal{P} \ge 0$.
  Since $\adm(t) \subseteq \reach(t)$, we have
  $\frac{d}{dt}V_R \ge 0$.
  Reachability preserves future reconstruction pathways,
  so $\frac{d}{dt}\kappa_\repair \ge 0$.
  Memory integrates recoverability
  $\mathcal{M}(t) = \int \reco(d,t)\,d\mu_D(d)$,
  so $\frac{d}{dt}\mathcal{M} \ge 0$ follows from
  non-decreasing $\kappa_\repair$.
  Memory preserves distinctions, so $D_\mathcal{E}$
  is non-decreasing.
  The converse follows by reversing the chain:
  non-decreasing $D_\mathcal{E}$ implies non-decreasing
  admissible future manifold.
\end{proof}

\begin{theorem}{Unified Invariant Theorem}{unified-invariant}
  The following quantities are projections of the
  possibility functional $\mathcal{I} = \Vol(\adm)$:
  \begin{align*}
    D_\mathcal{E} &= \Pi_D(\mathcal{I}), &
    \mathcal{M} &= \Pi_M(\mathcal{I}), \\
    \kappa_\repair &= \Pi_\repair(\mathcal{I}), &
    V_R &= \Pi_V(\mathcal{I}), \\
    S_R &= -\log\reco = \Pi_S(\mathcal{I}).
  \end{align*}
  Each is obtained from $\mathcal{I}$ by restriction,
  marginalisation, logarithmic transformation, or
  supremum.
\end{theorem}

\begin{proof}
  Distinctions determine recoverability via the
  Distinction--Entropy Duality (\cref{law:entropy-dual}).
  Recoverability determines memory
  ($\mathcal{M} = \int\reco\,d\mu_D$).
  Memory determines repair (Repair Existence Theorem,
  \cref{thm:repair-exist}).
  Repair determines reachability (admissible repair
  does not reduce $V_R$, \cref{prop:repair-adm}).
  Reachability determines admissibility
  ($\adm \subseteq \reach$, \cref{defn:adm-manifold}).
  Entropy is the negative log of recoverability
  (\cref{ch:recoverability}).
  Each quantity is therefore obtained from
  $\mathcal{I}$ by a specific transformation.
\end{proof}

\noindent
The Unified Invariant Theorem is not merely a formal
unification. It says something physically and
philosophically significant: all of the quantities
scientists, philosophers, and engineers have introduced
to measure the health of complex systems —
information, entropy, memory, repair capacity,
evolutionary potential, reachability, admissibility —
are different faces of the same geometric object.
That object is future possibility.

% ------------------------------------------------------------
\section{The Category of Regenerative Systems}
\label{sec:ch30-category}
% ------------------------------------------------------------

The Preservation Hierarchy Theorem identifies a class
of systems — $\mathfrak{G}$, regenerative systems —
that is closed under composition. This closure gives
regenerative systems a categorical structure.

\begin{theorem}{Category of Regenerative Systems}%
  {regen-category}
  Regenerative systems form a category $\mathbf{Reg}$
  whose objects are distinction ecologies and whose
  morphisms are possibility-preserving maps
  $f: \mathcal{E}_1 \to \mathcal{E}_2$ satisfying
  $\mathcal{P}(f(x)) \ge \mathcal{P}(x)$.
\end{theorem}

\begin{proof}
  \textit{Identity}: $\mathrm{id}_\mathcal{E}$ preserves
  $\mathcal{P}$ trivially.

  \textit{Composition}: If $f: \mathcal{E}_1 \to
  \mathcal{E}_2$ and $g: \mathcal{E}_2 \to \mathcal{E}_3$
  are both possibility-preserving, then
  \[
    \mathcal{P}(g(f(x)))
    \ge \mathcal{P}(f(x))
    \ge \mathcal{P}(x),
  \]
  so $g \circ f$ is possibility-preserving.

  \textit{Associativity}: function composition
  is associative.
\end{proof}

\begin{remark}{What $\mathbf{Reg}$ captures}{reg-meaning}
  The category $\mathbf{Reg}$ is the correct setting
  for asking questions about the evolution and
  interaction of complex systems.
  Morphisms in $\mathbf{Reg}$ are not arbitrary maps
  between systems. They are maps that respect the
  fundamental conserved quantity: admissible future
  volume.
  A therapy that maps a diseased system to a healthy
  one is a morphism in $\mathbf{Reg}$ iff it preserves
  or expands future possibility.
  A governance reform that maps an extractive
  institution to a generative one is a morphism in
  $\mathbf{Reg}$.
  A scientific revolution that maps a saturated paradigm
  to an enlarged one is a morphism in $\mathbf{Reg}$.
  The realization functors $F_i$ of Appendix~D
  (Chapter~\ref{app:deps}) are all functors
  $F_i: \mathbf{Abs} \to \mathbf{Reg}$ from the
  abstract theorem category into domain-specific
  regenerative systems.
\end{remark}

% ------------------------------------------------------------
\section{The Ecology Conservation Law}
\label{sec:ch30-conservation}
% ------------------------------------------------------------

The deepest single equation of the book is not the
RSVP field equation, nor the GAT condition
$\frac{d}{dt}\Vol(\adm) \ge 0$.
It is the conservation law that governs how admissible
future volume changes.

\begin{theorem}{Ecology of Distinctions Conservation Law}%
  {ecology-final}
  For any distinction ecology $\mathcal{E}$:
  \[
    \frac{d}{dt}\mathcal{P}(t)
    \;=\;
    \mathcal{G}(t) + \mathcal{R}(t)
    - \mathcal{L}(t) - \mathcal{X}(t),
  \]
  where:
  \begin{itemize}[nosep]
    \item $\mathcal{G}$: rate of new admissible
          distinction generation.
    \item $\mathcal{R}$: rate of admissible repair
          (restoration of damaged distinctions).
    \item $\mathcal{L}$: rate of distinction loss
          through entropy and projection failure.
    \item $\mathcal{X}$: extractive contraction
          (trajectories that consume future possibility
          for present gain).
  \end{itemize}
\end{theorem}

\begin{proof}
  All changes to admissible future volume arise from
  exactly four sources.
  Creation of new distinctions that are admissible
  enlarges $\adm$: contributes $+\mathcal{G}$.
  Repair restores formerly admissible distinctions
  to the admissibility manifold: contributes $+\mathcal{R}$.
  Entropy growth and projection failure remove recoverable
  distinctions from the admissibility manifold:
  contributes $-\mathcal{L}$.
  Extractive trajectories consume the distinction
  structures that sustained admissibility, producing
  a net reduction in $\Vol(\adm)$ even when the system
  continues: contributes $-\mathcal{X}$.
  These four contributions are additive and exhaustive,
  giving the stated equation.
\end{proof}

\begin{corollary}{Regime Classification}{ecology-regime}
  A distinction ecology is:
  \begin{align*}
    \text{regenerative} &\iff
    \mathcal{G}+\mathcal{R} > \mathcal{L}+\mathcal{X}, \\
    \text{stable}       &\iff
    \mathcal{G}+\mathcal{R} = \mathcal{L}+\mathcal{X}, \\
    \text{collapsing}   &\iff
    \mathcal{G}+\mathcal{R} < \mathcal{L}+\mathcal{X}.
  \end{align*}
\end{corollary}

\noindent
The Conservation Law is the formal counterpart of the
biological truism that life maintains local order by
exporting entropy. But it goes further. Life does not
merely export entropy — it generates new distinctions
faster than entropy erodes them, repairs damaged
distinctions faster than they degrade, and resists
extractive pressures that would consume its future.
Regenerative systems satisfy $\mathcal{G} + \mathcal{R}
> \mathcal{L} + \mathcal{X}$. Everything else follows.

% ------------------------------------------------------------
\section{Realizations as Functors}
\label{sec:ch30-functors}
% ------------------------------------------------------------

Each realization in Parts~VI--IX is not an application
of the framework to a domain. It is a \emph{functor}
from the abstract theorem category $\mathbf{Abs}$ —
the category whose objects are distinction-theoretic
concepts and whose morphisms are the logical relationships
proved in Parts~I--V — into a domain-specific category.

\begin{definition}{Realization Functor}{realization-functor}
  A \emph{realization functor}
  $F_i : \mathbf{Abs} \to \mathbf{Dom}_i$
  assigns to each abstract concept in $\mathbf{Abs}$
  a domain-specific instantiation in $\mathbf{Dom}_i$,
  and to each logical relationship a domain-specific
  correspondence, in a way that preserves composition
  and identity.
\end{definition}

The table below makes each functor explicit.

\begin{center}
\begin{tabular}{llll}
\toprule
Functor & Domain & Distinction & Admissibility \\
\midrule
$F_\mathrm{RSVP}$ & Physical fields &
  $\Phi$-capacity & $\Phi e^{-S} \ge \alpha$ \\
$F_\mathrm{Gravity}$ & Spacetime &
  Matter density & Admissible geodesic \\
$F_\mathrm{Cosmo}$ & Universe &
  $\mathcal{V}(t)$ & Expyrotic renewal \\
$F_\mathrm{Semantic}$ & Meaning manifold &
  Partition of $W$ & $V_S(m,t) > 0$ \\
$F_\mathrm{Conscious}$ & Self-model &
  Self-distinction & Recursive repair \\
$F_\mathrm{Pref}$ & Choice space &
  Utility partition & $\nabla\Phi_A \ge 0$ \\
$F_\mathrm{Flow}$ & Pipelines &
  $D_F(\delta)$ & Lossless stages \\
$F_\mathrm{Spherepop}$ & Event history &
  Scope $S$ & Refuse/Bind \\
$F_\mathrm{HYDRA}$ & Module network &
  $\mathcal{D}_i$ & Cross-module repair \\
$F_\mathrm{TARTAN}$ & Tiled manifold &
  Coherence tile & Tile-local $\adm$ \\
$F_\mathrm{Fiscal}$ & Policy space &
  $V_F(p,t)$ & $G_F+R_F \ge L_F$ \\
$F_\mathrm{Gov}$ & Institutions &
  $V_G(p,t)$ & Option preservation \\
$F_\mathrm{Science}$ & Epistemic ecology &
  $\mathcal{Q}_S$ & $|\mathcal{Q}_S|$ non-decreasing \\
\bottomrule
\end{tabular}
\end{center}

\noindent
The domains are wildly different in scale, substrate,
and vocabulary. Yet they instantiate the same abstract
theorem structure. The Repair Existence Theorem holds
in every row. The Persistence Hierarchy Theorem applies
to every column. The Ecology Conservation Law governs
the dynamics of every realization.

This is the strongest possible form of theoretical
unification: not mere analogy, but structural identity
up to the choice of functor.

% ------------------------------------------------------------
\section{What the Book Has Proved}
\label{sec:ch30-proved}
% ------------------------------------------------------------

It is worth being precise about what has been established
and what has been left as conjecture or sketch.

\paragraph{Established with full formal proof.}
\begin{itemize}
  \item The Axiom of Distinction and its immediate
        consequences (objects, information, cost, blind
        spots are co-produced).
  \item The Information--Distinction Theorem:
        information is not primitive.
  \item The Distinction--Entropy Duality: the Second Law
        is a theorem about distinction erosion.
  \item The History Primacy Theorem: state ontology is
        strictly contained in history ontology.
  \item The Law of Recoverability: dispersal and
        destruction are not equivalent.
  \item The Repair Existence Theorem: repair is possible
        iff recoverability is positive.
  \item The Repair Monoid: admissible repairs compose.
  \item The Persistent Anomaly Theorem: persistent
        anomalies lie outside the closure of the repair
        algebra's image.
  \item The Ontology Revision Theorem: a minimal
        ontological enlargement always exists.
  \item The Intelligence--Repair Equivalence Theorem.
  \item The Alignment Theorem.
  \item The Continuation Degeneracy Theorem.
  \item The Regeneration Theorem.
  \item The Reachability Monotonicity and
        Constraint Volume theorems.
  \item The Future Cone Theorem.
  \item The Admissibility Existence, Distortion,
        and Projection--Gap theorems.
  \item The Future Volume, Bottleneck, Preservation,
        and Curvature theorems.
  \item The Future Distinction Optimality Theorem
        (GAT): generative trajectories dominate
        extractive ones.
  \item The Preservation of Possibility Theorem:
        regeneration, memory, repair, and admissibility
        are equivalent under smooth dynamics.
  \item The Ecology of Distinctions Conservation Law.
\end{itemize}

\paragraph{Established as sketches requiring stronger
hypotheses.}
\begin{itemize}
  \item The Minimal Repair Theorem
        (requires compactness of the operator space).
  \item Several RSVP field theorems
        (the physical interpretation requires empirical
        support beyond the formal structure).
  \item The consciousness theorems
        (the connection between recursive self-repair
        and phenomenal consciousness is formal;
        the hard problem remains).
  \item The cosmological theorems
        (expyrotic renewal is proposed, not derived
        from first principles of particle physics).
\end{itemize}

\paragraph{Left as open problems.}
\begin{itemize}
  \item The precise topology of $\adm$ for infinite-
        dimensional distinction spaces.
  \item The relationship between admissibility curvature
        and phase transitions in statistical mechanics.
  \item A complete proof of the Generative Admissibility
        Theorem for stochastic dynamics.
  \item The categorical structure of the full functor
        category $[\mathbf{Abs}, \mathbf{Reg}]$.
\end{itemize}

% ------------------------------------------------------------
\section{The Transition to Chapter 31}
\label{sec:ch30-transition}
% ------------------------------------------------------------

The present chapter has shown that the framework is
unified: all its concepts are projections of a single
invariant, all its domains are functors into the same
category, and all its dynamical equations are
special cases of one conservation law.

But there remains something the conservation law
does not address.

The conservation law tells us how $\mathcal{P}(t)$
changes. It does not tell us why preserving
$\mathcal{P}$ matters.

The Future Distinction Optimality Theorem proved that
generative trajectories are \emph{strategically}
dominant. But strategic dominance is not the same
as value. An agent with an extremely short time horizon
may rationally prefer an extractive trajectory.
A culture that does not care about future generations
may rationally pursue pathological continuation.

Chapter~\ref{ch:preservation} addresses this.
It asks not how $\mathcal{P}$ changes but what it
would mean for the change to matter. It proves the
strongest theorem in the book: that admissible
future volume is not merely a useful quantity but
the correct invariant for evaluating any system
whose existence depends on the continuation of
distinction-production.

That proof is the subject of the final chapter.

\section*{Chapter Summary}
\begin{chapsummary}
  \item The argument of the book is a single progression:
        Distinction $\to$ Information $\to$ Entropy
        $\to$ History $\to$ Recoverability $\to$ Repair
        $\to$ Regeneration $\to$ Reachability
        $\to$ Admissibility $\to$ Generativity.
  \item Systems are ranked by a strict hierarchy
        $\mathfrak{P} \subsetneq \mathfrak{A}
        \subsetneq \mathfrak{G} \subsetneq \mathfrak{R}
        \subsetneq \mathfrak{M} \subsetneq \mathfrak{D}$
        (\cref{thm:preservation-hierarchy}).
  \item All major quantities — entropy, recoverability,
        memory, repair capacity, reachability, admissibility
        — are projections of the single invariant
        $\mathcal{I} = \Vol(\adm)$
        (\cref{thm:unified-invariant}).
  \item Regenerative systems form a category
        $\mathbf{Reg}$ under possibility-preserving maps
        (\cref{thm:regen-category}).
  \item The dynamics of all distinction ecologies are
        governed by:
        $\dot{\mathcal{P}} = \mathcal{G} + \mathcal{R}
        - \mathcal{L} - \mathcal{X}$
        (\cref{thm:ecology-final}).
  \item Each realization in Parts~VI--IX is a functor
        $F_i: \mathbf{Abs} \to \mathbf{Dom}_i$ that
        instantiates the abstract theorem structure
        in a domain-specific setting.
  \item The book has fully proved the structural spine;
        several realization theorems are sketches
        requiring additional assumptions.
  \item What remains for Chapter~\ref{ch:preservation}:
        not how $\mathcal{P}$ changes but why its
        preservation constitutes a value, not merely
        a strategy.
\end{chapsummary}

\subsection*{Exercises}

\begin{exercise}
  Prove that the Preservation Hierarchy is strict at
  every level by constructing an explicit system
  that belongs to $\mathfrak{R}$ but not $\mathfrak{G}$.
  (Hint: a system that can repair its distinctions but
  whose repair enzymes are not themselves repaired.)
\end{exercise}

\begin{exercise}
  Apply the Ecology Conservation Law to the Roman
  Republic in its final century (133--27 BCE).
  Identify $\mathcal{G}$, $\mathcal{R}$, $\mathcal{L}$,
  and $\mathcal{X}$ for each decade.
  At what point did $\mathcal{G}+\mathcal{R}
  < \mathcal{L}+\mathcal{X}$ become persistent?
\end{exercise}

\begin{exercise}
  Prove or disprove: the category $\mathbf{Reg}$ has
  a terminal object. If it does, what is it?
  If it does not, what does the non-existence imply
  about the structure of regenerative systems?
\end{exercise}

\begin{exercise}
  Identify which of the following are morphisms in
  $\mathbf{Reg}$ and which are not, justifying each:
  (a) a vaccine that confers immunity;
  (b) a monopoly buyout of a competitor;
  (c) a constitutional amendment adding a new right;
  (d) a scientific paradigm shift satisfying the
  Ontology Revision Theorem;
  (e) a \texttt{git reset --hard} on a shared branch.
\end{exercise}

\begin{exercise}
  The Unified Invariant Theorem says that entropy,
  recoverability, memory, repair capacity, reachability,
  and admissibility are all projections of
  $\mathcal{I} = \Vol(\adm)$.
  Construct an explicit projection operator $\Pi_S$
  such that $\Pi_S(\mathcal{I}) = S_R = -\log\reco$.
  Verify that $\Pi_S$ commutes with the admissibility
  dynamics.
\end{exercise}

\chapter{The Preservation of Possibility}
\label{ch:preservation}
\begin{epigraph}{Which structures preserve the possibility
  of future distinction-production?}{Author}
\end{epigraph}

\begin{definition}{Possibility Functional}{possibility-functional}
  $\mathcal{P}(x,t)=\Vol(\adm(x,t))
  =\mu_R(\reach(x,t)\cap\adm(x,t))$.
  \index{possibility functional|textbf}
\end{definition}

\begin{theorem}{Distinction Balance Law}{dist-balance}
  $\frac{d}{dt}D_\mathcal{E}(t)=G(t)+R(t)-L(t)$
  where $G,R,L$ are generation, repair, and loss rates.
  \index{distinction balance law|textbf}
\end{theorem}

\begin{theorem}{Common Invariant Theorem}{common-invariant}
  Entropy, recoverability, memory, repair capacity,
  reachability, and admissibility are all projections
  of the possibility functional $\mathcal{P}$.
\end{theorem}

\begin{theorem}{Preservation of Possibility Theorem}%
  {poss-preserve}
  Under smooth dynamics and finite measure, the following
  are equivalent:
  (i)~$\mathcal{E}$ is regenerative;
  (ii)~$\frac{d}{dt}\kappa_\repair\ge0$;
  (iii)~$\frac{d}{dt}\mathcal{M}\ge0$;
  (iv)~$\frac{d}{dt}D_\mathcal{E}\ge0$;
  (v)~$\frac{d}{dt}\mathcal{P}(x,t)\ge0$.
\end{theorem}
\begin{proof}
  (i)$\Leftrightarrow$(ii): Regeneration Theorem.
  (ii)$\Rightarrow$(iii): $\mathcal{M}=\int\reco\,d\mu_D$
  is non-decreasing when $\kappa_\repair$ is non-decreasing.
  (iii)$\Rightarrow$(iv): support measure of
  $\reco$ field is non-decreasing.
  (iv)$\Rightarrow$(v): non-decreasing $D$ means
  admissible future manifold cannot contract.
  (v)$\Rightarrow$(i): collapsing $\kappa_\repair$ would
  force $\Vol(\adm)$ to decrease
  (\cref{thm:future-vol}); contradiction.
\end{proof}

\begin{theorem}{Universal Regeneration Theorem}{universal-regen}
  $V_\infty(\mathcal{E})=\liminf_{t\to\infty}
  \mathcal{P}(x,t)>0$
  iff $\exists\epsilon>0$ such that eventually
  $\mathcal{P}(x,t)\ge\epsilon$.
  \index{universal regeneration theorem}
\end{theorem}

\begin{theorem}{Survival--Viability Separation Theorem}%
  {survival-viability}
  There exist systems that survive indefinitely while
  $\mathcal{P}(x,t)\to0$.
  \index{survival!vs viability}
\end{theorem}
\begin{proof}
  $X=[0,1]$, $\gamma(t)$ occupying $[0,e^{-t}]$:
  trajectory exists for all $t$ but
  $\mathcal{P}\to0$.
\end{proof}

\begin{principle}{Generative Admissibility Principle}%
  {gap-final}
  A trajectory is valuable insofar as it preserves the
  capacity for future distinction-production.
\end{principle}

\begin{theorem}{Generative Admissibility Theorem}{gat-final}
  $\gamma:[t_0,T]\to X$ is generatively admissible iff
  $\frac{d}{dt}\Vol(\adm(\gamma(t),t))\ge0$.
  Among equal-reward trajectories, generatively admissible
  ones maximise future distinction-producing capacity in
  lexicographic order over $\mathbf{I}_\adm$
  (\cref{thm:gat-ch15}).
\end{theorem}

\begin{theorem}{Future Distinction Dominance Theorem}%
  {future-dist-dom}
  For equal cumulative reward, generative $\gamma_1$
  over extractive $\gamma_2$:
  $D_f(\gamma_1(T),T')>D_f(\gamma_2(T),T')$
  for all sufficiently large $T'$.
\end{theorem}

\begin{theorem}{Ecology of Distinctions Conservation Law}%
  {ecology-conservation}
  $\frac{d}{dt}\mathcal{P}(t)
  =\mathcal{G}(t)+\mathcal{R}(t)
  -\mathcal{L}(t)-\mathcal{X}(t)$.
\end{theorem}

\begin{corollary}{Regenerative Condition}{regen-condition-final}
  A distinction ecology is generatively admissible iff
  $\mathcal{G}+\mathcal{R}\ge\mathcal{L}+\mathcal{X}$.
\end{corollary}

\begin{theorem}{Ecology of Distinctions Theorem}{ecology-final-31}
  The central object preserved by viable physical,
  biological, cognitive, computational, scientific,
  social, and cosmological systems is not matter,
  information, entropy, order, optimisation, or
  continuation, but \emph{admissible future distinction
  capacity}: $\liminf_{t\to\infty}\Vol(\adm(t))>0$.
\end{theorem}
\begin{proof}
  Matter, information, entropy, order, optimisation,
  and continuation may each change while viability
  persists or fails. By the Universal Regeneration
  Theorem, positive long-term viability is equivalent
  to $\liminf\mathcal{P}>0$, which equals
  $\liminf\Vol(\adm)>0$.
\end{proof}

\section*{Chapter Summary}
\begin{chapsummary}
  \item The possibility functional
        $\mathcal{P}=\Vol(\adm)$ is the single
        invariant underlying the entire framework.
  \item Regeneration, memory, repair, and admissibility
        are all equivalent under smooth dynamics
        (\cref{thm:poss-preserve}).
  \item Survival and viability are independent
        (\cref{thm:survival-viability}).
  \item The deepest preserved quantity is admissible
        future distinction capacity
        (\cref{thm:ecology-final-31}).
\end{chapsummary}

% ============================================================
\backmatter
% ============================================================

% ---- Appendix A --------------------------------------------
\appendix

% ============================================================
%  APPENDIX A: MATHEMATICAL PREREQUISITES
% ============================================================
\chapter{Mathematical Prerequisites}
\label{app:math}

This appendix collects the mathematical structures used
throughout the text. Proofs are omitted except where
required to establish notation or close a gap not covered
by standard references.

\section{Measure-Theoretic Foundations}

Let $(X, \Sigma, \mu)$ denote a measure space.
For any measurable $A \subseteq X$, $\mu(A) \ge 0$.
A \emph{probability space} $(X, \Sigma, P)$ satisfies
$P(X) = 1$.
Throughout the text distinguishability measures are
represented as induced probability measures over
equivalence classes.

\begin{definition}{Distinction Measure}{app-dist-measure}
  Given equivalence relation $\sim\,\subseteq X \times X$,
  the \emph{distinction measure} is
  $D(X, \sim) = \log |X/{\sim}|$.
  For infinite spaces: $D(X, \sim) = \log \mu(X/{\sim})$.
\end{definition}

\section{Information Geometry}

Let $\mathcal{P} = \{p(x;\theta)\}$ be a statistical
manifold. The \emph{Fisher metric} is
\[
  g_{ij}
  = \mathbb{E}\!\left[
      \frac{\partial\log p}{\partial\theta_i}
      \frac{\partial\log p}{\partial\theta_j}
    \right],
  \qquad
  ds^2 = g_{ij}\,d\theta^i d\theta^j.
\]
The Levi-Civita connection:
\[
  \Gamma^k_{ij}
  = \tfrac{1}{2} g^{km}
    \bigl(\partial_i g_{jm} + \partial_j g_{im}
          - \partial_m g_{ij}\bigr).
\]
The Riemann tensor:
\[
  R^i{}_{\!jkl}
  = \partial_k\Gamma^i_{jl} - \partial_l\Gamma^i_{jk}
  + \Gamma^i_{mk}\Gamma^m_{jl}
  - \Gamma^i_{ml}\Gamma^m_{jk}.
\]
Scalar curvature: $R = g^{ij}R_{ij}$.
Many admissibility quantities are expressed as curvature
in induced information manifolds.

\section{Entropy}

Shannon entropy: $H(X) = -\sum_i p_i \log p_i$.

Conditional: $H(X \mid Y) = H(X,Y) - H(Y)$.

Mutual information: $I(X;Y) = H(X)+H(Y)-H(X,Y)$.

KL divergence: $D_\mathrm{KL}(P \| Q)
= \sum_i P_i \log(P_i/Q_i)$.

Entropy production rate: $\sigma = dS/dt$.

Throughout the text entropy is interpreted as hidden
multiplicity beneath distinctions
(Chapter~\ref{ch:entropy}).

\section{Dynamical Systems}

For $\dot{x} = f(x)$, a fixed point satisfies
$f(x^*) = 0$. Linearisation: $\dot{\delta x} = J\,\delta x$
with $J_{ij} = \partial f_i/\partial x_j$.
Stability requires $\mathrm{Re}(\lambda_i) < 0$.
Lyapunov functions satisfy $V(x) > 0$ and
$\dot{V}(x) \le 0$.
Many regeneration results are expressed using
generalised Lyapunov arguments
(Chapter~\ref{ch:regeneration}).

\section{Category-Theoretic Structures}

A category $\mathcal{C}$ consists of objects
$\mathrm{Obj}(\mathcal{C})$ and morphisms
$\mathrm{Mor}(\mathcal{C})$ satisfying
$(h \circ g) \circ f = h \circ (g \circ f)$ and
$f \circ \mathrm{id}_A = f$.
Repair systems form categories whose morphisms are repair
operations (Chapter~\ref{ch:repair}).
Regenerative systems form the category $\mathbf{Reg}$
(Chapter~\ref{ch:ecology-synthesis}).

\section{Topology}

A topological space $(X, \tau)$:
boundary $\partial A = \overline{A} \setminus A^\circ$,
closure $\overline{A}$, interior $A^\circ$.
Admissibility boundaries frequently correspond to
topological phase transitions
(Chapter~\ref{ch:geometry-futures}).

\section{Differential Geometry}

A manifold $M$ is locally homeomorphic to $\mathbb{R}^n$.
Tangent space $T_p M$; cotangent space $T_p^* M$.
Differential forms: $\omega = \sum_i \omega_i\,dx^i$,
$d^2 = 0$.
Lie derivative $\mathcal{L}_X$.
The admissibility manifold $\adm$ (Chapter~\ref{ch:admissibility})
and meaning manifold $\mathcal{M}$
(Chapter~\ref{ch:semantic-geometry}) are the primary
manifolds of the text.

\section{Spectral Theory}

For operator $L$: eigenvalues satisfy $L\phi = \lambda\phi$.
Graph Laplacian $L = D - A$; spectral gap
$\gamma = \lambda_2 - \lambda_1$.
Regenerative stability often depends upon spectral gaps
(Chapter~\ref{ch:distinction-ecology}).

\section{Optimal Transport}

Wasserstein distance:
$W_p(\mu,\nu) = \bigl(\inf_\gamma \int d(x,y)^p\,
d\gamma\bigr)^{1/p}$, where $\gamma$ ranges over
couplings. Repair trajectories are modelled as transport
flows (Chapter~\ref{ch:repair}).

\section{Large-Deviation Theory}

Rate function $I(x)$; large deviation principle:
$P(X_n \approx x) \sim e^{-nI(x)}$.
The most probable trajectory minimises $I$.
Admissibility transitions are naturally described by
large-deviation arguments.

\section{Geometric Reachability}

For $\dot{x} = f(x,u)$, reachable volume
$\Vol_R(t) = \mu(R_t)$ and reachability entropy
$S_R = \log\Vol_R$.
Boundary collapse: $\Vol_R \to 0$.
This quantity appears throughout Chapters~13--31.

\section{Admissibility Geometry}

Admissible region $\adm \subseteq X$;
admissibility volume $V_A = \mu(\adm)$;
admissibility curvature $K_A = -\nabla^2\log V_A$.
Generative admissibility condition:
\[
  \frac{dV_A}{dt} \ge 0.
\]
This is the fundamental geometric condition of the
framework (Chapter~\ref{ch:geometry-futures}).

\section{Notation Table}

\begin{center}
\begin{tabular}{ll}
\toprule
Symbol & Meaning \\
\midrule
$\Phi$ & capacity / scalar field (RSVP) \\
$\mathbf{v}$ & transport field \\
$S$ & constraint / entropy field \\
$R_t$ & reachable set at time $t$ \\
$\adm$ & admissible region \\
$V_A$ & admissible volume \\
$D$ & distinction measure \\
$\repair$ & repair operator \\
$\reco$ & recoverability functional \\
$\hist$ & history space \\
$\partition$ & partition \\
$\Pi$ & projection map \\
$\mathcal{I}(\Sigma)$ & intelligence of system $\Sigma$ \\
$\kappa_\Sigma$ & repair capacity of $\Sigma$ \\
$\mathbf{I}_\adm$ & full admissibility invariant \\
\bottomrule
\end{tabular}
\end{center}

% ============================================================
%  APPENDIX B: COMPUTATIONAL TOOLS
% ============================================================
\chapter{Computational Tools}
\label{app:comp}

\section{Computation as Distinction Transformation}

Traditional computation is state transition:
$x_0 \to x_1 \to x_2 \to \cdots$.
The distinction-first perspective reframes this.
States are compressed histories; histories are compressed
distinction operations. A computation is a trajectory
through distinction space:
$d_0 \to d_1 \to d_2 \to \cdots$.

\section{Flow Computing}

A flow system is a directed acyclic graph $G = (V,E)$
where vertices are transformations and edges transport
distinctions.
The computational state at time $t$ is
$X_t = \bigcup_{v \in V_t}\mathrm{Output}(v)$.

\begin{theorem}{Flow Equivalence Theorem}{flow-equiv}
  Two computational systems are observationally equivalent
  iff they induce identical distinction transport operators.
\end{theorem}

\begin{proof}
  Observations depend only upon distinctions available
  to an observer. Identical transport operators imply
  identical reachable distinctions. Any change in
  transport changes at least one reachable distinction.
\end{proof}

Define transport entropy $S_F = -\sum_i p_i\log p_i$
where $p_i = f_i / \sum_j f_j$ and $f_i$ is edge flow.

\begin{theorem}{Flow Compression Theorem}{flow-compression}
  Maximum compression occurs when transport entropy is
  minimised subject to reachability constraints.
\end{theorem}

\section{Spherepop Formal Semantics}

A Spherepop configuration is $\Sigma = (H, S)$
where $H$ is history and $S$ is scope.
The transition relation $\Sigma \to \Sigma'$:

\medskip
\begin{center}
\begin{tabular}{lll}
\toprule
Rule & Pre-state & Post-state \\
\midrule
$\mathbf{Pop}(e)$      & $(H \cup \{e\}, S)$ & $(H, S)$ \\
$\mathbf{Refuse}(e)$   & $(H, S \cup \{e\})$ & $(H, S)$ + witness \\
$\mathbf{Bind}(n,v)$   & $(H, S)$ & $(H, S \cup \{n \mapsto v\})$ \\
$\mathbf{Collapse}(\pi)$ & $(H, S)$ & $(H, \pi(S))$ \\
\bottomrule
\end{tabular}
\end{center}

\begin{theorem}{History Monotonicity}{app-history-mono}
  For any valid Spherepop computation,
  $H_t \sqsubseteq H_{t+1}$.
\end{theorem}

\begin{proof}
  No operation removes historical events.
  Pop, Refuse, Bind, and Collapse each append to $H$
  or leave it unchanged.
\end{proof}

Define collapse quotient $Q_C = |S| / |\pi(S)|$.

\begin{theorem}{Collapse Bound}{collapse-bound}
  For finite scope, $Q_C \ge 1$.
\end{theorem}

\section{HYDRA State Transitions}

HYDRA state space: $M = \prod_{i=1}^k M_i$.
Module dynamics: $\dot{M}_i = F_i(M)$.
Disagreement energy: $E_D = \sum_{i<j} d(M_i, M_j)^2$.

\begin{theorem}{HYDRA Convergence Theorem}{hydra-conv}
  If $\dot{E}_D < 0$ for all non-equilibrium states,
  module consensus is globally asymptotically stable.
\end{theorem}

\begin{proof}
  $E_D$ is a positive-definite Lyapunov function with
  negative-definite derivative. Standard Lyapunov theory
  implies convergence.
\end{proof}

\section{TARTAN Recursive Distinction}

Semantic manifold $M$ partitioned into tiles $T_i$;
recursive subdivision: $T_i = \bigcup_j T_{ij}$.

\begin{theorem}{Recursive Distinction Theorem}{recursive-dist}
  If each subdivision increases distinguishability by
  factor $\lambda > 1$, then after depth $n$:
  $D_n = \lambda^n D_0$.
\end{theorem}

Semantic curvature: $K_S = -\nabla^2\log\rho$
where $\rho(x)$ is distinction density.
Positive curvature: semantic bottleneck.
Negative curvature: semantic expansion.

\section{MEM|8 Memory Architecture}

Memory field: $M(x,t) = \sum_i w_i K(x, e_i)$
where $e_i$ are stored events and $K$ is a kernel.
Ecphory (recall) occurs when $M(x,t) > \theta$.

\begin{theorem}{Ecphory Threshold Theorem}{ecphory}
  Recall occurs iff $\sum_i w_i K(x, e_i) > \theta$.
\end{theorem}

\section{Computational Admissibility}

For computational system $C$, reachable volume
$V_C(t) = |R_C(t)|$.

\begin{theorem}{Computational Admissibility Theorem}{comp-adm-app}
  A computation is generatively admissible iff
  $\frac{d}{dt}V_C(t) \ge 0$.
\end{theorem}

\begin{theorem}{Computational Conservation Law}{comp-conserve}
  If admissible volume is preserved, distinction capacity
  cannot collapse:
  $\frac{dA}{dt} \ge 0 \implies \frac{dD}{dt} \ge -\epsilon$
  for bounded projection loss $\epsilon$.
\end{theorem}

\begin{proof}
  Admissible volume bounds reachability volume.
  Reachability volume bounds distinction capacity
  (\cref{thm:reach-dist}).
  The result follows by transitivity.
\end{proof}

% ============================================================
%  APPENDIX C: BIOLOGICAL BACKGROUND
% ============================================================
\chapter{Biological Background}
\label{app:bio}

\section{Life as Distinction Preservation}

A living system is a structure capable of maintaining
distinctions against entropic erosion.
A cell preserves inside/outside.
A genome preserves functional/non-functional sequences.
An immune system preserves self/non-self.
A nervous system preserves relevant/irrelevant signals.
An ecosystem preserves viable/non-viable energy pathways.
At every scale biological persistence requires repair.

\section{Thermodynamic Setting}

For biological system $\Sigma$ in environment $E$:
$\frac{dS_{\Sigma \cup E}}{dt} \ge 0$ (Second Law).
Local entropy balance:
$\frac{dS_\Sigma}{dt} = \sigma - J_S$,
where $\sigma$ is internal production and $J_S$ is export.
Persistence requires $J_S > \sigma$.
Repair is thermodynamically directed entropy export.

\section{Autopoiesis}

Production process set $P = \{p_1,\ldots,p_n\}$.
Autopoietic closure: $\forall p_i \in P,\; \exists P_i
\subseteq P$ such that $P_i \to p_i$.

\begin{theorem}{Autopoietic Closure Theorem}{autopoietic}
  Autopoietic systems form strongly connected directed
  graphs.
\end{theorem}

\section{Genomes as Historical Archives}

Evolutionary information:
$I_E(G) = \log[P(G \mid \text{selection})/P(G \mid \text{random})]$.
Large values correspond to long histories of selective
retention. Let $\mu$ = mutation rate, $\rho$ = repair rate.

\begin{theorem}{Genomic Stability Theorem}{genomic-stability}
  If $\rho < \mu$, long-term distinction preservation
  is impossible: expected damage grows as $(\mu-\rho)t$.
\end{theorem}

\section{Immune System as Admissibility Filter}

Healthy states $A \subseteq X$; immune classifier
$f: X \to \{0,1\}$.
Type-I error ($f(x)=0$, $x \in A$): missed pathogen.
Type-II error ($f(x)=1$, $x \notin A$): autoimmunity.
Cancer corresponds to systematic Type-I failure.

\begin{theorem}{Cancer Continuation Theorem}{cancer-cont}
  Tumour success and organism admissibility are
  anti-correlated: increasing tumour continuation reduces
  organism admissible volume $V_A(t) \to 0$.
\end{theorem}

\section{Developmental Landscapes}

Following \textcite{waddington1957strategy},
developmental trajectories follow $\dot{x} = -\nabla U$
on potential surface $U(x)$.
Stable cell types correspond to minima $\nabla U = 0$.
Differentiation is symmetry breaking and reachability
volume reduction (Chapter~\ref{ch:reachability}).

\section{Evolutionary Reachability}

Evolutionary volume $V_E(t) = |R_t|$ where $R_t$ is
the set of genotypes reachable through mutation and
selection.

\begin{theorem}{Evolutionary Constraint Theorem}{evo-constraint}
  Increasing specialisation generally decreases $V_E$
  (by Theorem~\ref{thm:reach-mono}).
\end{theorem}

\begin{theorem}{Evolutionary Distinction Theorem}{evo-dist}
  Major evolutionary transitions increase total
  distinction capacity: $D_{n+1} > D_n$.
\end{theorem}

\begin{proof}
  New organisational levels introduce distinctions
  unavailable at lower scales
  \parencite{maynard-smith1995major}.
\end{proof}

\section{Neural Criticality}

Neural branching ratio $\sigma$:
subcritical ($\sigma<1$), critical ($\sigma=1$),
supercritical ($\sigma>1$).

\begin{theorem}{Critical Distinction Theorem}{critical-dist}
  Maximum distinction-processing capacity occurs
  near $\sigma = 1$.
\end{theorem}

\begin{proof}
  Subcritical: distinctions decay.
  Supercritical: distinctions collapse to synchrony.
  Critical: propagation and separation balance.
\end{proof}

\section{Ecological Diversity and Repair}

Ecosystem as distinction ecology:
species $\leftrightarrow$ distinction-maintaining processes;
food web $\leftrightarrow$ repair dependencies.

\begin{theorem}{Ecological Diversity Theorem}{eco-diversity-app}
  Increasing species diversity weakly increases repair
  capacity (Diversity--Repair Theorem,
  \cref{thm:diversity-repair}).
\end{theorem}

\section{Biological Admissibility}

Biological admissibility volume $V_B(t)$.
Generative condition: $\frac{dV_B}{dt} \ge 0$.
Evolutionary trajectories satisfying this expand future
adaptive possibilities.

The distinction between fitness (immediate reproductive
success) and admissibility (preservation of future
distinction-production) parallels the distinction between
continuation and regeneration developed in
Chapter~\ref{ch:beyond-continuation}.

% ============================================================
%  APPENDIX D: PROOF DEPENDENCY TABLES
% ============================================================
\chapter{Proof Dependency Tables}
\label{app:deps}

\section{Purpose}

The argument of this book is cumulative.
Every theorem is derived from previous theorems.
This appendix exposes that dependency structure as a
directed acyclic graph $\mathcal{G} = (V, E)$
where vertices are results and directed edges encode
proof dependence.

\section{Dependency Relation}

$T_i \prec T_j$ iff theorem $T_j$ requires $T_i$
in its proof. The dependency graph is
$V = \{T_1, \ldots, T_n\}$ and
$E = \{(T_i, T_j) : T_i \prec T_j\}$.

\begin{theorem}{Acyclicity Theorem}{acyclicity}
  The theorem dependency graph is a directed acyclic
  graph (DAG).
\end{theorem}

\begin{proof}
  All theorems are introduced in chapter order.
  A theorem may only depend upon previously established
  results. Cycles are therefore impossible.
\end{proof}

\section{Foundational Basis}

Define the foundational basis:
\[
  \mathcal{B}
  = \{A_1, A_2, A_3, A_4, A_5\},
\]
where $A_1 =$ Axiom of Distinction,
$A_2 =$ Information--Distinction Theorem,
$A_3 =$ Distinction--Entropy Duality,
$A_4 =$ Law of Historical Compression,
$A_5 =$ Law of Recoverability.

\begin{theorem}{Foundation Theorem}{foundation}
  For every theorem $T$ in the text, $\exists A_i \in
  \mathcal{B}$ such that $A_i \prec^* T$
  (transitive closure).
\end{theorem}

\begin{proof}
  Inspection of the proof graph. All later chapters
  invoke recoverability, history, entropy, information,
  or distinction structure.
\end{proof}

\section{Dependency Depth Table}

\begin{center}
\begin{tabular}{lc}
\toprule
Result & Approximate depth \\
\midrule
Axiom of Distinction & 0 \\
Information--Distinction Theorem & 1 \\
Distinction--Entropy Duality & 2 \\
Law of Historical Compression & 3 \\
Law of Recoverability & 4 \\
Repair Existence Theorem & 5 \\
Persistent Anomaly Theorem & 6 \\
Regeneration Theorem & 7 \\
Future Cone Theorem & 8 \\
Admissibility Existence Theorem & 9 \\
Future Distinction Optimality Theorem & 10 \\
Preservation of Possibility Theorem & 11 \\
Ecology of Distinctions Theorem & 12 \\
\bottomrule
\end{tabular}
\end{center}

\section{Layer-by-Layer Dependency Matrices}

\paragraph{Repair layer.}
\begin{center}
\begin{tabular}{lcccc}
\toprule
& $A_5$ & Repair Exist. & Closure & Entropy \\
\midrule
$A_5$ & 1&1&1&1 \\
Repair Exist. & 0&1&1&1 \\
Closure & 0&0&1&1 \\
Entropy & 0&0&0&1 \\
\bottomrule
\end{tabular}
\end{center}

\paragraph{Geometry layer.}
\begin{center}
\begin{tabular}{lcccc}
\toprule
& Reach. & Adm. & Future Vol. & GAT \\
\midrule
Reachability & 1&1&1&1 \\
Admissibility & 0&1&1&1 \\
Future Volume & 0&0&1&1 \\
GAT & 0&0&0&1 \\
\bottomrule
\end{tabular}
\end{center}

\section{Failure Layer Dependency}

\begin{theorem}{Failure Dependency Theorem}{failure-dep}
  No theorem in Chapter~\ref{ch:geometry-failure} is
  derivable without Chapter~\ref{ch:repair}.
\end{theorem}

\begin{proof}
  Persistent anomalies are defined relative to repair
  operators. All later results inherit this dependence.
\end{proof}

\section{Master Dependency Chain}

The entire book can be represented as:
\[
  \text{Distinction}
  \to \text{Information}
  \to \text{Entropy}
  \to \text{History}
  \to \text{Recoverability}
  \to \text{Repair}
  \to \text{Failure}
  \to \text{Intelligence}
  \to \text{Regeneration}
  \to \text{Ecology}
  \to \text{Reachability}
  \to \text{Admissibility}
  \to \text{GAT}.
\]

\section{Realization Functors}

Each realization chapter corresponds to a functor
$F_i : \mathcal{T} \to \mathcal{D}_i$ mapping the
abstract theorem category $\mathcal{T}$ into a domain:

\begin{center}
\begin{tabular}{ll}
\toprule
Functor & Domain \\
\midrule
$F_\mathrm{RSVP}$ & Physical fields \\
$F_\mathrm{Gravity}$ & Distinction dynamics \\
$F_\mathrm{Cosmo}$ & Expyrotic renewal \\
$F_\mathrm{Semantic}$ & Meaning manifolds \\
$F_\mathrm{Conscious}$ & Self-reconstruction \\
$F_\mathrm{Pref}$ & Admissibility gradients \\
$F_\mathrm{Flow}$ & Distinction transport \\
$F_\mathrm{Spherepop}$ & History-native computation \\
$F_\mathrm{HYDRA}$ & Multi-module repair \\
$F_\mathrm{TARTAN}$ & Multiscale tiling \\
$F_\mathrm{Fiscal}$ & Policy reachability \\
$F_\mathrm{Gov}$ & Institutional repair \\
$F_\mathrm{Science}$ & Epistemic regeneration \\
\bottomrule
\end{tabular}
\end{center}

\begin{theorem}{Closure Theorem}{dep-closure}
  The transitive closure of GAT contains every major
  theorem in the text: $\mathcal{B} \subseteq \mathcal{G}_A$
  where $\mathcal{G}_A = \{T : T \prec^* \text{GAT}\}$.
\end{theorem}

\begin{proof}
  GAT $\prec^*$ Admissibility $\prec^*$ Reachability
  $\prec^*$ Regeneration $\prec^*$ Repair
  $\prec^*$ Recoverability $\prec^*$ History
  $\prec^*$ Entropy $\prec^*$ Information
  $\prec^*$ Distinction ($= A_1$).
  Hence $\mathcal{B} \subseteq \mathcal{G}_A$.
\end{proof}

\section{Minimal Axiom Counts}

\begin{center}
\begin{tabular}{lc}
\toprule
Theorem & $N(T)$ (foundational axioms required) \\
\midrule
Information--Distinction & 1 \\
Entropy & 2 \\
History Primacy & 3 \\
Repair Existence & 5 \\
Regeneration & 6 \\
Future Cone & 7 \\
Admissibility Existence & 8 \\
Future Distinction Optimality & 8 \\
Preservation of Possibility & 8 \\
\bottomrule
\end{tabular}
\end{center}

The entire monograph may be viewed as a proof that the
Generative Admissibility Principle is a consequence of
the Axiom of Distinction together with the mathematics
of entropy, history, repair, regeneration, and
reachability.

% ============================================================
%  distinctions.bib  — paste into a separate file
%  named distinctions.bib in the same directory
% ============================================================
\printglossaries

\printbibliography[heading=bibintoc,title={References}]

% ---- Bibliography note -------------------------------------
% Full bibliography in distinctions.bib.
% Key entries cited above:
% aurelius, shannon1948, bateson1972, spencer-brown1969,
% boltzmann1877, whitehead1929, bergson1896, rescher1996,
% tulving1983, schrodinger1944, landauer1961irreversibility,
% bennett1982, kuhn1962structure, lakatos1978,
% pearl2000causality, spirtes2000, amodei2016, russell2019,
% bostrom2014, ostrom1990, holling1973, gunderson2002,
% may1973, levin1998, tilman1996, maynard-smith1995major,
% waddington1957strategy, sontag1998, maturana1980,
% gardenfors2000conceptual, mcilroy1978, milner1989,
% north1990, levitsky2018, steinhardt2002, penrose2010,
% misner1973, bekenstein1973, hawking1975,
% maclane1971, awodey2010, fong2019, villani2009,
% dembo1998, dobzhansky1973, meadows2008

% Note: add to distinctions.bib:
% @book{aurelius,
%   author    = {Aurelius, Marcus},
%   title     = {Meditations},
%   publisher = {Penguin Classics},
%   year      = {180},
%   note      = {Translation by Gregory Hays, 2002}}
%
% @article{dobzhansky1973,
%   author  = {Dobzhansky, Theodosius},
%   title   = {Nothing in Biology Makes Sense Except
%              in the Light of Evolution},
%   journal = {The American Biology Teacher},
%   volume  = {35}, number = {3}, pages = {125--129},
%   year    = {1973}}

\printindex

\end{document}
